\documentclass[a4paper,12pt]{article}
\usepackage[utf8]{inputenc}
\usepackage[T1]{fontenc}
\usepackage[italian]{babel}
\usepackage[sfdefault]{atkinson}
\renewcommand*\familydefault{\sfdefault}
\usepackage{float}
\usepackage{microtype}
\usepackage{geometry}
\usepackage{setspace}
\usepackage{enumitem}
\usepackage{titlesec}
\usepackage{chngpage}
\usepackage{tocloft}
\usepackage{longtable}
\usepackage{ltablex}
\usepackage{array}
\usepackage{graphicx}
\usepackage{fancyhdr}
\usepackage[table]{xcolor}
\usepackage{color,soul}
\usepackage[most]{tcolorbox}
\usepackage{nameref}
\usepackage[colorlinks=true]{hyperref}

% --- INSERISCI QUESTO NEL PREAMBOLO DEL TUO DOCUMENTO ---
\usepackage{amssymb} % per il simbolo \checkmark
\usepackage{pifont}

% Macro per semplificare e alleggerire il codice della tabella
\newcommand{\ok}{\textcolor{green!60!black}{$\checkmark$ Superato}}
\newcommand{\ko}{\textcolor{red!80!black}{$\times$ Mancante}}
\newcommand{\warn}[1]{\textcolor{orange!90!black}{\textbf{!} Multiplo (#1)}}
% Use URL-style line breaking to keep long test file names inside table borders.
\newcommand{\f}[1]{\nolinkurl{#1}\newline}
\newcommand{\q}[1]{\textit{#1}}
\newcommand{\sep}{\newline \vspace{2mm} \hrule \vspace{2mm}}
% --------------------------------------------------------

% Colori ZPUS - Verde, Nero, Bianco
\definecolor{zpusgreen}{RGB}{4, 138, 55}
\definecolor{zpusdarkgreen}{RGB}{0, 100, 0}
\definecolor{zpusblack}{RGB}{0, 0, 0}
\definecolor{zpuswhite}{RGB}{255, 255, 255}
\definecolor{zpuslightgray}{RGB}{245, 245, 245}

\hypersetup{
    linkcolor=black,
    urlcolor=blue
}

\definecolor{lightblack}{gray}{0.35}
\newcommand{\glossario}[1]{\textit{#1}\textsubscript{\textbf{\textit{\textcolor{lightblack}{G}}}}}
\newcommand{\ped}[1]{\textsubscript{#1}}

\pagestyle{fancy}
\setlength{\headwidth}{\textwidth}
\fancyhfoffset[L,R]{0pt}
\lhead{\rightmark}
\rhead{7-ZPUs}
\lfoot{Piano di Qualifica}
\rfoot{\thepage}
\cfoot{}
\renewcommand{\headrulewidth}{0.8pt}
\renewcommand{\footrulewidth}{0.8pt}

\renewcommand{\contentsname}{Indice}
\renewcommand{\listoftables}{Elenco delle tabelle}
\renewcommand{\listoffigures}{Elenco delle immagini}

\geometry{margin=2.5cm}
\setstretch{1.1}

\titleclass{\subsubsubsection}{straight}[\subsection]

\newcounter{subsubsubsection}[subsubsection]
\renewcommand\thesubsubsubsection{\thesubsubsection.\arabic{subsubsubsection}}
\renewcommand\theparagraph{\thesubsubsubsection.\arabic{paragraph}} % optional; useful if paragraphs are to be numbered

\titleformat{\subsubsubsection}
  {\normalfont\normalsize\bfseries}{\thesubsubsubsection}{1em}{}
\titlespacing*{\subsubsubsection}
{0pt}{3.25ex plus 1ex minus .2ex}{1.5ex plus .2ex}

\makeatletter
\renewcommand\paragraph{\@startsection{paragraph}{5}{\z@}\%
  {3.25ex \@plus1ex \@minus.2ex}\%
  {-1em}\%
  {\normalfont\normalsize\bfseries}}
\renewcommand\subparagraph{\@startsection{subparagraph}{6}{\parindent}\%
  {3.25ex \@plus1ex \@minus .2ex}\%
  {-1em}\%
  {\normalfont\normalsize\bfseries}}
\def\toclevel@subsubsubsection{4}
\def\toclevel@paragraph{5}
%\def\toclevel@paragraph{6}
\def\toclevel@subparagraph{6}
\def\l@subsubsubsection{\@dottedtocline{4}{7em}{4em}}
\def\l@paragraph{\@dottedtocline{5}{10em}{5em}}
\def\l@subparagraph{\@dottedtocline{6}{14em}{6em}}
\makeatother

\setcounter{secnumdepth}{4}
\setcounter{tocdepth}{4}

\titleformat{\section}{\LARGE\bfseries}{\thesection}{1em}{}
\titleformat{\subsection}{\Large\bfseries}{\thesubsection}{1em}{}
\titleformat{\subsubsection}{\large\bfseries}{\thesubsubsection}{1em}{}
\titleformat{\paragraph}{\large\bfseries}{\theparagraph}{1em}{}
\titleformat{\subparagraph}{\normalsize\bfseries}{\thesubparagraph}{1em}{}

\setlength\parindent{0pt}

\begin{document}

\begin{center}
    \includegraphics[width=9.5cm]{../../assets/logo7ZPUs.jpg}\\
    \small\hspace{10cm} 7zpus.swe@gmail.com\\
    \vspace{0.5cm}
    \LARGE \textbf{Appendice - Piano di Qualifica}\\
    \vspace{0.3cm}
    \normalsize \textbf{Versione:} 2.0.0\\
    \textbf{Responsabile:} Georgescu Diana \\
    \textbf{Destinatari:} Sanmarco Informatica S.p.A., Prof. Tullio Vardanega, Prof. Riccardo Cardin, Team 7-ZPUs \\
\end{center}

\vspace{0.3cm}
\hrule
\vspace{0.5cm}


\section{Test di unità}

\setlength{\LTleft}{-0.4cm}
\renewcommand{\arraystretch}{1.2}
\rowcolors{2}{gray!15}{white}
\begin{longtable}{|p{2.0cm}|p{4.3cm}|p{7.7cm}|p{1.5cm}|}
\hline
\rowcolor{zpusgreen!30}
\textbf{ID Test} & \textbf{Test} & \textbf{Posizioni (File e Titolo)} & \textbf{Stato} \\
\hline
\endfirsthead
\rowcolor{zpusgreen!30}
\textbf{ID Test} & \textbf{Test} & \textbf{Posizioni (File e Titolo)} & \textbf{Stato} \\
\hline
\endhead
TU-1 & DipDAO > saves dip with UNKNOWN status & saved.getId() non deve essere null; saved.getUuid() deve essere "dip-1"; saved.getIntegrityStatus() deve essere IntegrityStatusEnum.UNKNOWN & \ok \\
TU-2 & DipDAO > updates existing dip on uuid conflict & count deve essere 1; second.getId() deve essere first.getId(); second.getIntegrityStatus() deve essere IntegrityStatusEnum.UNKNOWN & \ok \\
TU-3 & DipDAO > gets by id and uuid & byId?.getUuid() deve essere "dip-3"; byUuid?.getId() deve essere saved.getId() & \ok \\
TU-4 & DipDAO > filters by status and updates integrity & valid.map((d) => d.getUuid()) deve essere uguale a ["dip-a"]; invalid.map((d) => d.getUuid()) deve essere uguale a ["dip-b"] & \ok \\
TU-5 & DocumentClassDAO > saves and gets a document class & saved.getId() non deve essere null; found?.getUuid() deve essere "dc-uuid"; found?.getDipUuid() deve essere "dip-uuid" & \ok \\
TU-6 & DocumentClassDAO > updates by uuid conflict without duplicating rows & rows deve avere lunghezza 1; rows[0].name deve essere "Nome 2" & \ok \\
TU-7 & DocumentClassDAO > gets by dipId, status and search & dao.getByDipId(dipId) deve avere lunghezza 2; dao.getByStatus(IntegrityStatusEnum.VALID) deve avere lunghezza 1; dao.search("Contr") deve avere lunghezza 1; dao.search("NON-EXISTING") deve avere lunghezza 0 & \ok \\
TU-8 & DocumentClassDAO > computes aggregated integrity status & integritya deve essere IntegrityStatusEnum.VALID; integrityb deve essere IntegrityStatusEnum.VALID; integritya2 deve essere IntegrityStatusEnum.INVALID; integrityb2 deve essere IntegrityStatusEnum.INVALID & \ok \\
TU-9 & DocumentDAO > saves and retrieves document with metadata tree & found non deve essere null; found?.getUuid() deve essere "doc-a"; found?.getProcessId() deve essere processId; found?.getMetadata().findNodeByName("Titolo")?.getStringValue() deve essere "A" & \ok \\
TU-10 & DocumentDAO > gets by process and status and updates integrity & dao.getByProcessId(processId) deve avere lunghezza 2; dao.getByStatus(IntegrityStatusEnum.VALID) deve avere lunghezza 1; dao.getByStatus(IntegrityStatusEnum.INVALID) deve avere lunghezza 1 & \ok \\
TU-11 & DocumentDAO > searchDocumentSemantic uses document\_vector and returns top matches & results deve avere lunghezza 2; results[0].document.getUuid() deve essere "doc-sem-a"; results[0].score deve essere maggiore di results[1].score & \ok \\
TU-12 & DocumentDAO > getIndexedDocumentsCount reads from document\_vector & dao.getIndexedDocumentsCount() deve essere 1 & \ok \\
TU-13 & FileDAO > saves and retrieves file & found non deve essere null; found?.getFilename() deve essere "main.xml"; found?.getDocumentId() deve essere documentId; found?.getIsMain() deve essere true & \ok \\
TU-14 & FileDAO > updates on (document\_id, path) conflict without duplicating rows & rows deve avere lunghezza 1; rows[0].id deve essere updated.getId(); rows[0].filename deve essere "main-v2.xml"; rows[0].hash deve essere "hash-new" & \ok \\
TU-15 & FileDAO > gets by document and status and updates integrity & dao.getByDocumentId(documentId) deve avere lunghezza 2; dao.getByStatus(IntegrityStatusEnum.VALID) deve avere lunghezza 1; dao.getByStatus(IntegrityStatusEnum.INVALID) deve avere lunghezza 1 & \ok \\
TU-16 & DipMapper > maps persistence row to domain & dip.getId() deve essere 10; dip.getUuid() deve essere "dip-uuid"; dip.getIntegrityStatus() deve essere IntegrityStatusEnum.VALID & \ok \\
TU-17 & DipMapper > falls back to UNKNOWN for invalid status & dip.getIntegrityStatus() deve essere IntegrityStatusEnum.UNKNOWN & \ok \\
TU-18 & DipMapper > maps domain to dto & dto.id deve essere 42; dto.uuid deve essere "dip-uuid"; dto.integrityStatus deve essere IntegrityStatusEnum.INVALID & \ok \\
TU-19 & DipMapper > throws when converting to dto with null id & () => DipMapper.toDTO(dip) deve lanciare "Cannot convert Dip to DTO: id is null" & \ok \\
TU-20 & DipMapper > maps domain to persistence model & model deve essere uguale a { uuid: "dip-uuid", integrityStatus: IntegrityStatusEnum.VALID, } & \ok \\
TU-21 & DocumentClassMapper > maps persistence row to domain & entity.getId() deve essere 7; entity.getDipId() deve essere 3; entity.getDipUuid() deve essere "dip-uuid"; entity.getUuid() deve essere "dc-uuid"; entity.getIntegrityStatus() deve essere IntegrityStatusEnum.VALID & \ok \\
TU-22 & DocumentClassMapper > uses UNKNOWN integrity when status is missing & entity.getIntegrityStatus() deve essere IntegrityStatusEnum.UNKNOWN & \ok \\
TU-23 & DocumentClassMapper > maps domain to dto & dto deve essere uguale a { id: 12, dipId: 5, uuid: "dc-uuid", name: "Classe", timestamp: "2026-01-01T00:00:00Z", integrityStatus: IntegrityStatusEnum.INVALID, } & \ok \\
TU-24 & DocumentClassMapper > uses -1 dipId in dto when dipId is null & dto.dipId deve essere -1 & \ok \\
TU-25 & DocumentClassMapper > throws when converting to dto with null id & () => DocumentClassMapper.toDTO(entity) deve lanciare "Cannot convert to DTO: DocumentClass entity is not yet persisted and has no ID." & \ok \\
TU-26 & DocumentClassMapper > maps domain to search result including name & result deve essere uguale a { id: 21, uuid: "dc-uuid", integrityStatus: IntegrityStatusEnum.VALID, name: "Classe Risorse Umane", type: "CLASSE", } & \ok \\
TU-27 & DocumentClassMapper > throws when converting to search result with null id & () => DocumentClassMapper.toSearchResult(entity) deve lanciare "Cannot convert to SearchResult: DocumentClass entity is not yet persisted and has no ID." & \ok \\
TU-28 & DocumentClassMapper > maps domain to persistence model & model deve essere uguale a { dipUuid: "dip-uuid", uuid: "dc-uuid", integrityStatus: IntegrityStatusEnum.VALID, name: "Classe", timestamp: "2026-01-01T00:00:00Z", } & \ok \\
TU-29 & DocumentMapper > maps persistence row to domain & entity.getId() deve essere 4; entity.getUuid() deve essere "doc-uuid"; entity.getProcessId() deve essere 8; entity.getProcessUuid() deve essere "proc-uuid"; entity.getIntegrityStatus() deve essere IntegrityStatusEnum.VALID; entity.getMetadata().getName() deve essere "Documento"; entity.getMetadata().getChildren() deve avere lunghezza 1 & \ok \\
TU-30 & DocumentMapper > uses UNKNOWN when status is missing or invalid & DocumentMapper.fromPersistence( missingStatus, metadataRows, ).getIntegrityStatus() deve essere IntegrityStatusEnum.UNKNOWN; DocumentMapper.fromPersistence( invalidStatus, metadataRows, ).getIntegrityStatus() deve essere IntegrityStatusEnum.UNKNOWN & \ok \\
TU-31 & DocumentMapper > maps domain to dto & dto.id deve essere 10; dto.processId deve essere 20; dto.uuid deve essere "doc-uuid"; dto.integrityStatus deve essere IntegrityStatusEnum.INVALID; dto.metadata.name deve essere "Documento" & \ok \\
TU-32 & DocumentMapper > throws when converting to dto with null id & () => DocumentMapper.toDTO(entity) deve lanciare "Cannot convert to DTO: Document entity is not yet persisted and has no ID." & \ok \\
TU-33 & DocumentMapper > uses -1 processId in dto when processId is null & dto.processId deve essere -1 & \ok \\
TU-34 & DocumentMapper > maps domain to persistence model & model.uuid deve essere "doc-uuid"; model.integrityStatus deve essere IntegrityStatusEnum.VALID; model.processUuid deve essere "proc-uuid"; model.metadata.map((m) => m.getName()) deve essere uguale a ["Titolo", "Nome"] & \ok \\
TU-35 & DocumentMapper > metadataToJson preserves repeated sibling nodes as arrays & Array.isArray(ruolo) deve essere true; ruolo deve avere lunghezza 2; ruolo[0].Altro.PF.Cognome deve essere "Rossi"; ruolo[0].Altro.PF.Nome deve essere "Paolo"; ruolo[1].ResponsabileGestioneDocumentale.PF.Cognome deve essere "Colombo" & \ok \\
TU-36 & DocumentMapper > metadataJsonToRoot restores repeated nodes from arrays & soggetti non deve essere null; ruoloNodes deve avere lunghezza 2; root.findNodeByName("Altro")?.findNodeByName("Cognome")?.getStringValue() deve essere "Rossi"; root .findNodeByName("ResponsabileGestioneDocumentale") ?.findNodeByName("Cognome") ?.getStringValue() deve essere "Colombo" & \ok \\
TU-37 & FileMapper > maps persistence row to domain & entity.getId() deve essere 9; entity.getUuid() deve essere "file-uuid"; entity.getFilename() deve essere "main.xml"; entity.getIsMain() deve essere true; entity.getDocumentId() deve essere 3; entity.getDocumentUuid() deve essere "doc-uuid"; entity.getIntegrityStatus() deve essere IntegrityStatusEnum.VALID & \ok \\
TU-38 & FileMapper > falls back to UNKNOWN for invalid status & entity.getIsMain() deve essere false; entity.getIntegrityStatus() deve essere IntegrityStatusEnum.UNKNOWN & \ok \\
TU-39 & FileMapper > maps domain to dto & dto deve essere uguale a { id: 20, documentId: 33, filename: "main.xml", path: "/pkg/main.xml", hash: "hash", integrityStatus: IntegrityStatusEnum.INVALID, isMain: true, } & \ok \\
TU-40 & FileMapper > uses -1 documentId in dto when documentId is null & dto.documentId deve essere -1 & \ok \\
TU-41 & FileMapper > throws when converting to dto with null id & () => FileMapper.toDTO(entity) deve lanciare "Cannot convert to DTO: File entity is not yet persisted and has no ID." & \ok \\
TU-42 & FileMapper > maps domain to persistence model & model deve essere uguale a { filename: "main.xml", path: "/pkg/main.xml", hash: "hash", integrityStatus: IntegrityStatusEnum.VALID, isMain: 0, uuid: "file-uuid", documentUuid: "doc-uuid", } & \ok \\
TU-43 & MetadataKeyMapper > toPascalCase preserves canonical XML keys with underscores & MetadataKeyMapper.toPascalCase( "DocumentoInformatico.Soggetti.Ruolo.Destinatario.PG.CodiceFiscale\_PartitaIva", ) deve essere "DocumentoInformatico.Soggetti.Ruolo.Destinatario.PG.CodiceFiscale\_PartitaIva" & \ok \\
TU-44 & MetadataKeyMapper > toPascalCase still normalizes snake\_case and kebab-case segments & MetadataKeyMapper.toPascalCase("soggetti.codice\_fiscale") deve essere "Soggetti.CodiceFiscale"; MetadataKeyMapper.toPascalCase("soggetti.codice-fiscale") deve essere "Soggetti.CodiceFiscale" & \ok \\
TU-45 & MetadataKeyMapper > fromLegacyFilters drops null, empty and whitespace-only string values & group.logicOperator deve essere "AND"; group.items deve essere uguale a [ { path: "Soggetto.Nome", operator: "EQ", value: "Mario", }, ] & \ok \\
TU-46 & MetadataKeyMapper > fromLegacyFilters keeps meaningful non-string values & group.items deve essere uguale a [ { path: "Flags.Attivo", operator: "EQ", value: false, }, { path: "Metrics.Count", operator: "EQ", value: 0, }, ] & \ok \\
TU-47 & MetadataKeyMapper > mapGroup preserves CodiceFiscale\_PartitaIva inside ELEM\_MATCH payload & nested.items[0].path deve essere "Destinatario.PG.CodiceFiscale\_PartitaIva" & \ok \\
TU-48 & MetadataMapper.fromFlatRows > throws when rows are empty & () => MetadataMapper.fromFlatRows([]) deve lanciare "Cannot build metadata tree from empty rows" & \ok \\
TU-49 & MetadataMapper.fromFlatRows > builds a synthetic root when multiple top-level rows exist & root non deve essere null; root?.getName() deve essere "root"; root?.getType() deve essere MetadataType.COMPOSITE; root?.getChildren() deve avere lunghezza 2; subject?.getName() deve essere "Soggetto"; subject?.getType() deve essere MetadataType.COMPOSITE; subject?.getChildren() deve avere lunghezza 1; subject?.getChildren()[0].getName() deve essere "Nome"; subject?.getChildren()[0].getStringValue() deve essere "Mario" & \ok \\
TU-50 & MetadataMapper.fromFlatRows > normalizes lowercase metadata types & root?.getType() deve essere MetadataType.COMPOSITE; root?.getChildren() deve avere lunghezza 1; root?.getChildren()[0].getType() deve essere MetadataType.NUMBER & \ok \\
TU-51 & MetadataMapper.fromFlatRows > throws when a parent row is missing & () => MetadataMapper.fromFlatRows(rows) deve lanciare "Invalid metadata tree: missing parent row for id 100 (parent\_id=999)" & \ok \\
TU-52 & MetadataMapper.fromFlatRows > throws when duplicate row ids are present & () => MetadataMapper.fromFlatRows(rows) deve lanciare "Invalid metadata tree: duplicate row id 1" & \ok \\
TU-53 & MetadataMapper.fromFlatRows > throws when no root nodes are found & () => MetadataMapper.fromFlatRows(rows) deve lanciare "Invalid metadata tree: no root nodes found" & \ok \\
TU-54 & MetadataMapper.fromFlatRows > throws when metadata type is invalid & () => MetadataMapper.fromFlatRows(rows) deve lanciare "Invalid metadata type: invalid" & \ok \\
TU-55 & MetadataMapper.flatten > returns the same node for string metadata & flat deve avere lunghezza 1; flat[0].getName() deve essere "Titolo"; flat[0].getStringValue() deve essere "Doc A" & \ok \\
TU-56 & MetadataMapper.flatten > returns leaf descendants for composite metadata & flat deve avere lunghezza 2; flat.map((node) => node.getName()) deve essere uguale a ["A", "C"] & \ok \\
TU-57 & ProcessMapper > maps persistence row to domain & entity.getId() deve essere 3; entity.getDocumentClassId() deve essere 9; entity.getDocumentClassUuid() deve essere "dc-uuid"; entity.getUuid() deve essere "proc-uuid"; entity.getIntegrityStatus() deve essere IntegrityStatusEnum.VALID; entity.getMetadata().getChildren() deve avere lunghezza 1 & \ok \\
TU-58 & ProcessMapper > uses UNKNOWN when status is missing or invalid & ProcessMapper.fromPersistence(missingStatus, metadataRows).getIntegrityStatus() deve essere IntegrityStatusEnum.UNKNOWN; ProcessMapper.fromPersistence(invalidStatus, metadataRows).getIntegrityStatus() deve essere IntegrityStatusEnum.UNKNOWN & \ok \\
TU-59 & ProcessMapper > maps domain to dto & dto.id deve essere 8; dto.documentClassId deve essere 11; dto.uuid deve essere "proc-uuid"; dto.integrityStatus deve essere IntegrityStatusEnum.INVALID; dto.metadata.name deve essere "Process" & \ok \\
TU-60 & ProcessMapper > uses -1 documentClassId in dto when id exists but class id is null & dto.documentClassId deve essere -1 & \ok \\
TU-61 & ProcessMapper > throws when converting to dto with null id & () => ProcessMapper.toDTO(entity) deve lanciare "Cannot convert to DTO: Process entity is not yet persisted and has no ID." & \ok \\
TU-62 & ProcessMapper > maps domain to persistence model & model.documentClassUuid deve essere "dc-uuid"; model.uuid deve essere "proc-uuid"; model.integrityStatus deve essere IntegrityStatusEnum.VALID; model.metadata.map((m) => m.getName()) deve essere uguale a ["Fase", "Step"] & \ok \\
TU-63 & MetadataHelper > persists recursive metadata tree & rows deve avere lunghezza 4; rows[0] deve essere uguale a { owner\_id: 1, parent\_id: null, name: "Root", value: "", type: "COMPOSITE", }; rows[1].name deve essere "Title"; rows[2].name deve essere "Soggetto"; rows[3].name deve essere "Nome" & \ok \\
TU-64 & MetadataHelper > loads metadata rows ordered by id & loaded deve avere lunghezza 2; loaded.map((r) => r.name) deve essere uguale a ["A", "B"] & \ok \\
TU-65 & MetadataHelper > accepts metadata arrays in saveMetadata & count deve essere 2 & \ok \\
TU-66 & ProcessDAO > saves and retrieves process with metadata tree & found non deve essere null; found?.getUuid() deve essere "proc-a"; found?.getDocumentClassId() deve essere documentClassId; found?.getMetadata().getName() deve essere "Process"; found?.getMetadata().findNodeByName("Versione")?.getStringValue() deve essere "v1" & \ok \\
TU-67 & ProcessDAO > gets by documentClassId and status and updates integrity & dao.getByDocumentClassId(documentClassId) deve avere lunghezza 2; dao.getByStatus(IntegrityStatusEnum.VALID) deve avere lunghezza 1; dao.getByStatus(IntegrityStatusEnum.INVALID) deve avere lunghezza 1 & \ok \\
TU-68 & DocumentMetadataQueryBuilder > returns empty query for empty filter group & result.sql deve essere ""; result.params deve essere uguale a [] & \ok \\
TU-69 & DocumentMetadataQueryBuilder > handles a single EQ condition & result.sql deve essere "(json\_extract(document.metadata, '\$.Doc.Type') = ? COLLATE NOCASE)"; result.params deve essere uguale a ["Invoice"] & \ok \\
TU-70 & DocumentMetadataQueryBuilder > converts boolean EQ values to SQLite-safe numeric params & result.sql deve essere "(json\_extract(document.metadata, '\$.DocumentoInformatico.Riservato') = ? COLLATE NOCASE)"; result.params deve essere uguale a [1] & \ok \\
TU-71 & DocumentMetadataQueryBuilder > handles a single GT condition & result.sql deve essere "(json\_extract(document.metadata, '\$.Amount') > ?)"; result.params deve essere uguale a [100] & \ok \\
TU-72 & DocumentMetadataQueryBuilder > handles a single LT condition & result.sql deve essere "(json\_extract(document.metadata, '\$.Amount') < ?)"; result.params deve essere uguale a [50] & \ok \\
TU-73 & DocumentMetadataQueryBuilder > handles a single LIKE condition & result.sql deve essere "(json\_extract(document.metadata, '\$.Name') LIKE ? COLLATE NOCASE)"; result.params deve essere uguale a ["\%test\%"] & \ok \\
TU-74 & DocumentMetadataQueryBuilder > handles a single IN condition & result.sql deve essere "(json\_extract(document.metadata, '\$.Status') IN (?, ?, ?))"; result.params deve essere uguale a ["A", "B", "C"] & \ok \\
TU-75 & DocumentMetadataQueryBuilder > returns empty for IN condition with empty array & result.sql deve essere ""; result.params deve essere uguale a [] & \ok \\
TU-76 & DocumentMetadataQueryBuilder > handles AND group with multiple scalar conditions & result.sql deve essere "(json\_extract(document.metadata, '\$.Doc.Type') = ? COLLATE NOCASE AND json\_extract(document.metadata, '\$.Status') IN (?, ?))"; result.params deve essere uguale a ["Invoice", "Paid", "Sent"] & \ok \\
TU-77 & DocumentMetadataQueryBuilder > handles OR group with multiple scalar conditions & result.sql deve essere "(json\_extract(document.metadata, '\$.Type') = ? COLLATE NOCASE OR json\_extract(document.metadata, '\$.Type') = ? COLLATE NOCASE)"; result.params deve essere uguale a ["A", "B"] & \ok \\
TU-78 & DocumentMetadataQueryBuilder > handles deeply nested groupings & result.sql deve essere "(json\_extract(document.metadata, '\$.Year') > ? AND (json\_extract(document.metadata, '\$.Dept') = ? COLLATE NOCASE OR json\_extract(document.metadata, '\$.Dept') = ? COLLATE NOCASE))"; result.params deve essere uguale a [2020, "IT", "HR"] & \ok \\
TU-79 & DocumentMetadataQueryBuilder > handles ELEM\_MATCH group for JSON arrays & result.sql deve essere "(EXISTS (SELECT 1 FROM json\_each(document.metadata, '\$.Items') AS el WHERE (json\_extract(el.value, '\$.Code') = ? COLLATE NOCASE AND json\_extract(el.value, '\$.Price') > ?)))"; result.params deve essere uguale a ["XYZ", 50] & \ok \\
TU-80 & DocumentMetadataQueryBuilder > throws when ELEM\_MATCH value is not a valid SearchGroup & () => builder.build(query) deve lanciare /ELEM\_MATCH requires a nested SearchGroup value/ & \ok \\
TU-81 & DocumentMetadataQueryBuilder > skips conditions with null or undefined value & result.sql deve essere "(json\_extract(document.metadata, '\$.Name') = ? COLLATE NOCASE)"; result.params deve essere uguale a ["Valid"] & \ok \\
TU-82 & DocumentMetadataQueryBuilder > skips conditions with whitespace-only string value & result.sql deve essere "(json\_extract(document.metadata, '\$.Dept') = ? COLLATE NOCASE)"; result.params deve essere uguale a ["IT"] & \ok \\
TU-83 & DocumentMetadataQueryBuilder > sanitizes IN arrays by dropping empty and whitespace-only values & result.sql deve essere "(json\_extract(document.metadata, '\$.Status') IN (?, ?))"; result.params deve essere uguale a ["Paid", "Sent"] & \ok \\
TU-84 & DocumentMetadataQueryBuilder > skips ELEM\_MATCH when nested conditions only contain whitespace-only values & result.sql deve essere ""; result.params deve essere uguale a [] & \ok \\
TU-85 & DocumentMetadataQueryBuilder > normalizes path by stripping start dollars and dots, and validates it & result.sql deve essere "(json\_extract(document.metadata, '\$.My.Path') = ? COLLATE NOCASE)"; result.params deve essere uguale a ["X"]; () => builder.build(badQuery1) deve lanciare /SearchCondition path cannot be empty/; () => builder.build(badQuery2) deve lanciare /Invalid SearchCondition path segment/ & \ok \\
TU-86 & VectorDAO > should save and retrieve a vector correctly & retrievedVector non deve essere null; retrievedVector?.getDocumentId() deve essere 1; retrievedVector?.getEmbedding() deve essere uguale a new Float32Array([0.1, 0.2, 0.3]) & \ok \\
TU-87 & VectorDAO > should return null for non-existent document ID & result deve essere null & \ok \\
TU-88 & VectorDAO > should return empty array when searching similar vectors with topK <= 0 & result deve essere uguale a [] & \ok \\
TU-89 & Dip entity > constructor > TU-F-browsing-01: getUuid() should assegna l'uuid passato & dip.getUuid() deve essere "abc-123" & \ok \\
TU-90 & Dip entity > constructor > TU-F-browsing-02: imposta() should imposta integrityStatus a UNKNOWN di default & dip.getIntegrityStatus() deve essere IntegrityStatusEnum.UNKNOWN & \ok \\
TU-91 & Dip entity > constructor > TU-F-browsing-03: id() should id è null finché non viene salvato nel DB & dip.getId() deve essere null & \ok \\
TU-92 & Dip entity > setIntegrityStatus > TU-F-browsing-05: aggiorna() should aggiorna lo stato di integrità & dip.getIntegrityStatus() deve essere IntegrityStatusEnum.VALID & \ok \\
TU-93 & Document entity > constructor > TU-F-browsing-08: assegna() should assegna uuid, metadata e processUuid & doc.getUuid() deve essere "doc-uuid"; doc.getMetadata() deve essere meta; doc.getProcessUuid() deve essere "process-uuid" & \ok \\
TU-94 & Document entity > constructor > TU-F-browsing-09: imposta() should imposta integrityStatus a UNKNOWN di default & doc.getIntegrityStatus() deve essere IntegrityStatusEnum.UNKNOWN & \ok \\
TU-95 & Document entity > constructor > TU-F-browsing-10: id() should id è null finché non viene persistito & doc.getId() deve essere null & \ok \\
TU-96 & Document entity > setIntegrityStatus > TU-F-browsing-13: aggiorna() should aggiorna lo stato di integrità & doc.getIntegrityStatus() deve essere IntegrityStatusEnum.INVALID & \ok \\
TU-97 & DocumentClass entity > constructor > TU-F-browsing-16: assegna() should assegna dipUuid, uuid, name e timestamp & dc.getDipUuid() deve essere "dip-uuid"; dc.getUuid() deve essere "dc-uuid"; dc.getName() deve essere "Contratti"; dc.getTimestamp() deve essere "2024-01-01T00:00:00Z" & \ok \\
TU-98 & DocumentClass entity > constructor > TU-F-browsing-17: imposta() should imposta integrityStatus a UNKNOWN di default & dc.getIntegrityStatus() deve essere IntegrityStatusEnum.UNKNOWN & \ok \\
TU-99 & DocumentClass entity > constructor > TU-F-browsing-18: id() should id è null finché non viene persistito & dc.getId() deve essere null & \ok \\
TU-100 & DocumentClass entity > setIntegrityStatus > TU-F-browsing-21: aggiorna() should aggiorna lo stato di integrità & dc.getIntegrityStatus() deve essere IntegrityStatusEnum.INVALID & \ok \\
TU-101 & constructor > TU-F-browsing-24: assegna() should assegna filename, path, isMain, uuid e documentUuid & file.getFilename() deve essere "documento.xml"; file.getPath() deve essere "/dip/subdir/documento.xml"; file.getIsMain() deve essere true; file.getDocumentUuid() deve essere "doc-uuid"; file.getUuid() deve essere "uuid" & \ok \\
TU-102 & constructor > TU-F-browsing-25: imposta() should imposta integrityStatus a UNKNOWN di default & file.getIntegrityStatus() deve essere IntegrityStatusEnum.UNKNOWN & \ok \\
TU-103 & constructor > TU-F-browsing-26: id() should id è null finché non viene persistito & file.getId() deve essere null & \ok \\
TU-104 & constructor > TU-F-browsing-27: isMain() should isMain = false viene mantenuto & file.getIsMain() deve essere false & \ok \\
TU-105 & setIntegrityStatus > TU-F-browsing-30: aggiorna() should aggiorna lo stato di integrità & file.getIntegrityStatus() deve essere IntegrityStatusEnum.INVALID & \ok \\
TU-106 & Process entity > constructor > TU-F-browsing-33: assegna() should assegna documentClassUuid, uuid e metadata & proc.getDocumentClassUuid() deve essere "dc-uuid"; proc.getUuid() deve essere "proc-uuid"; proc.getMetadata() deve essere meta & \ok \\
TU-107 & Process entity > constructor > TU-F-browsing-34: imposta() should imposta integrityStatus a UNKNOWN di default & proc.getIntegrityStatus() deve essere IntegrityStatusEnum.UNKNOWN & \ok \\
TU-108 & Process entity > constructor > TU-F-browsing-35: id() should id è null finché non viene persistito & proc.getId() deve essere null & \ok \\
TU-109 & Process entity > setIntegrityStatus > TU-F-browsing-38: aggiorna() should aggiorna lo stato di integrità & proc.getIntegrityStatus() deve essere IntegrityStatusEnum.INVALID & \ok \\
TU-110 & BrowsingIpcAdapter > TU-S-browsing-111: register() should resolve all browsing use-cases and register all browse channels handled by adapter & resolveMock deve essere stato chiamato con DocumentoUC.GET\_BY\_ID; resolveMock deve essere stato chiamato con DocumentoUC.GET\_BY\_PROCESS; resolveMock deve essere stato chiamato con DocumentoUC.GET\_BY\_STATUS; resolveMock deve essere stato chiamato con FileUC.GET\_BY\_ID; resolveMock deve essere stato chiamato con FileUC.GET\_BY\_DOCUMENT; resolveMock deve essere stato chiamato con FileUC.GET\_BY\_STATUS; resolveMock deve essere stato chiamato con ProcessUC.GET\_BY\_ID; resolveMock deve essere stato chiamato con ProcessUC.GET\_BY\_STATUS; resolveMock deve essere stato chiamato con ProcessUC.GET\_BY\_DOCUMENT\_CLASS; resolveMock deve essere stato chiamato con DocumentClassUC.GET\_BY\_DIP\_ID; resolveMock deve essere stato chiamato con DocumentClassUC.GET\_BY\_STATUS; resolveMock deve essere stato chiamato con DocumentClassUC.GET\_BY\_ID; resolveMock deve essere stato chiamato con DipUC.GET\_BY\_ID; resolveMock deve essere stato chiamato con DipUC.GET\_BY\_STATUS; handleMock deve essere stato chiamato 15 volte; handlers.has(IpcChannels.BROWSE\_GET\_DOCUMENT\_BY\_ID) deve essere true; handlers.has(IpcChannels.BROWSE\_GET\_DOCUMENTS\_BY\_PROCESS) deve essere true; handlers.has(IpcChannels.BROWSE\_GET\_DOCUMENTS\_BY\_STATUS) deve essere true; handlers.has(IpcChannels.BROWSE\_GET\_FILE\_BY\_ID) deve essere true; handlers.has(IpcChannels.BROWSE\_GET\_FILE\_BUFFER\_BY\_ID) deve essere true; handlers.has(IpcChannels.BROWSE\_GET\_FILE\_BY\_DOCUMENT) deve essere true; handlers.has(IpcChannels.BROWSE\_GET\_FILE\_BY\_STATUS) deve essere true; handlers.has(IpcChannels.BROWSE\_GET\_PROCESS\_BY\_ID) deve essere true; handlers.has(IpcChannels.BROWSE\_GET\_PROCESS\_BY\_STATUS) deve essere true; handlers.has(IpcChannels.BROWSE\_GET\_PROCESS\_BY\_DOCUMENT\_CLASS) deve essere true; handlers.has(IpcChannels.BROWSE\_GET\_DOCUMENT\_CLASS\_BY\_DIP\_ID) deve essere true; handlers.has(IpcChannels.BROWSE\_GET\_DOCUMENT\_CLASS\_BY\_STATUS) deve essere true; handlers.has(IpcChannels.BROWSE\_GET\_DOCUMENT\_CLASS\_BY\_ID) deve essere true; handlers.has(IpcChannels.BROWSE\_GET\_DIP\_BY\_ID) deve essere true; handlers.has(IpcChannels.BROWSE\_GET\_DIP\_BY\_STATUS) deve essere true; handlers.has(IpcChannels.BROWSE\_GET\_DIP\_BY\_DOCUMENT\_CLASS) deve essere false & \ok \\
TU-111 & BrowsingIpcAdapter > TU-S-browsing-112: register() should map all by-id channels to DTO when entity exists & getDocByIdUC.execute deve essere stato chiamato con 10; getFileByIdUC.execute deve essere stato chiamato con 11; getProcessByIdUC.execute deve essere stato chiamato con 12; getDocClassByIdUC.execute deve essere stato chiamato con 13; getDipByIdUC.execute deve essere stato chiamato con 14; documentToDtoMock deve essere stato chiamato con docEntity; fileToDtoMock deve essere stato chiamato con fileEntity; processToDtoMock deve essere stato chiamato con processEntity; documentClassToDtoMock deve essere stato chiamato con docClassEntity; dipToDtoMock deve essere stato chiamato con dipEntity; docResult deve essere uguale a { id: 1, dto: "doc" }; fileResult deve essere uguale a { id: 2, dto: "file" }; processResult deve essere uguale a { id: 3, dto: "process" }; docClassResult deve essere uguale a { id: 4, dto: "doc-class" }; dipResult deve essere uguale a { id: 5, dto: "dip" } & \ok \\
TU-112 & BrowsingIpcAdapter > TU-S-browsing-113: register() should return null for all by-id channels when entity is not found & docResult deve essere null; fileResult deve essere null; processResult deve essere null; docClassResult deve essere null; dipResult deve essere null; documentToDtoMock non deve essere stato chiamato; fileToDtoMock non deve essere stato chiamato; processToDtoMock non deve essere stato chiamato; documentClassToDtoMock non deve essere stato chiamato; dipToDtoMock non deve essere stato chiamato & \ok \\
TU-113 & BrowsingIpcAdapter > TU-S-browsing-114: register() should map document list channels to DTO and forward arguments & getDocByProcessUC.execute deve essere stato chiamato con 200; getDocByStatusUC.execute deve essere stato chiamato con IntegrityStatusEnum.VALID; documentToDtoMock deve essere stato chiamato con doc1; documentToDtoMock deve essere stato chiamato con doc2; byProcessResult deve essere uguale a [{ mapped: doc1 }, { mapped: doc2 }]; byStatusResult deve essere uguale a [{ mapped: doc1 }] & \ok \\
TU-114 & BrowsingIpcAdapter > TU-S-browsing-115: register() should map file list channels to DTO and forward arguments & getFileByDocUC.execute deve essere stato chiamato con 300; getFileByStatusUC.execute deve essere stato chiamato con IntegrityStatusEnum.UNKNOWN; fileToDtoMock deve essere stato chiamato con file1; fileToDtoMock deve essere stato chiamato con file2; byDocumentResult deve essere uguale a [{ mapped: file1 }, { mapped: file2 }]; byStatusResult deve essere uguale a [{ mapped: file2 }] & \ok \\
TU-115 & BrowsingIpcAdapter > TU-S-browsing-116: register() should map process list channels to DTO and forward arguments & getProcessByStatusUC.execute deve essere stato chiamato con IntegrityStatusEnum.INVALID; getProcessByDocumentClassUC.execute deve essere stato chiamato con 400; processToDtoMock deve essere stato chiamato con process1; processToDtoMock deve essere stato chiamato con process2; byStatusResult deve essere uguale a [ { mapped: process1 }, { mapped: process2 }, ]; byDocClassResult deve essere uguale a [{ mapped: process2 }] & \ok \\
TU-116 & BrowsingIpcAdapter > TU-S-browsing-117: register() should map document-class list channels to DTO and forward arguments & getDocClassByDipIdUC.execute deve essere stato chiamato con 500; getDocClassByStatusUC.execute deve essere stato chiamato con IntegrityStatusEnum.VALID; documentClassToDtoMock deve essere stato chiamato con docClass1; documentClassToDtoMock deve essere stato chiamato con docClass2; byDipResult deve essere uguale a [{ mapped: docClass1 }, { mapped: docClass2 }]; byStatusResult deve essere uguale a [{ mapped: docClass2 }] & \ok \\
TU-117 & BrowsingIpcAdapter > TU-S-browsing-118: register() should map dip-by-status channel to DTO and forward argument & getDipByStatusUC.execute deve essere stato chiamato con IntegrityStatusEnum.UNKNOWN; dipToDtoMock deve essere stato chiamato con dip1; dipToDtoMock deve essere stato chiamato con dip2; result deve essere uguale a [{ mapped: dip1 }, { mapped: dip2 }] & \ok \\
TU-118 & CheckIntegrityIpcAdapter > TU-S-browsing-107: register() should resolve all integrity use-cases and register all integrity channels & resolveMock deve essere stato chiamato con DocumentoUC.CHECK\_INTEGRITY\_STATUS; resolveMock deve essere stato chiamato con FileUC.CHECK\_INTEGRITY\_STATUS; resolveMock deve essere stato chiamato con ProcessUC.CHECK\_INTEGRITY\_STATUS; resolveMock deve essere stato chiamato con DocumentClassUC.CHECK\_INTEGRITY\_STATUS; resolveMock deve essere stato chiamato con DipUC.CHECK\_INTEGRITY\_STATUS; handleMock deve essere stato chiamato 5 volte; handleMock.mock.calls.some( (call) => call[0] === IpcChannels.CHECK\_DOCUMENT\_INTEGRITY\_STATUS, ) deve essere true; handleMock.mock.calls.some( (call) => call[0] === IpcChannels.CHECK\_FILE\_INTEGRITY\_STATUS, ) deve essere true; handleMock.mock.calls.some( (call) => call[0] === IpcChannels.CHECK\_PROCESS\_INTEGRITY\_STATUS, ) deve essere true; handleMock.mock.calls.some( (call) => call[0] === IpcChannels.CHECK\_DOCUMENT\_CLASS\_INTEGRITY\_STATUS, ) deve essere true; handleMock.mock.calls.some( (call) => call[0] === IpcChannels.CHECK\_DIP\_INTEGRITY\_STATUS, ) deve essere true & \ok \\
TU-119 & CheckIntegrityIpcAdapter > TU-S-browsing-108: register() should delegate non-file integrity checks to the resolved use-cases & checkDocumentIntegrityUC.execute deve essere stato chiamato con 10; checkProcessIntegrityUC.execute deve essere stato chiamato con 11; checkDocumentClassIntegrityUC.execute deve essere stato chiamato con 12; checkDipIntegrityUC.execute deve essere stato chiamato con 13; documentResult deve essere IntegrityStatusEnum.VALID; processResult deve essere IntegrityStatusEnum.INVALID; documentClassResult deve essere IntegrityStatusEnum.UNKNOWN; dipResult deve essere IntegrityStatusEnum.VALID & \ok \\
TU-120 & CheckIntegrityIpcAdapter > TU-S-browsing-109: register() should await and return file integrity result & checkFileIntegrityUC.execute deve essere stato chiamato con 22; result deve essere IntegrityStatusEnum.UNKNOWN & \ok \\
TU-121 & CheckIntegrityIpcAdapter > TU-S-browsing-110: register() should propagate rejection from file integrity use-case & getRegisteredHandler(handleMock, IpcChannels.CHECK\_FILE\_INTEGRITY\_STATUS)( {}, 99, ) deve lanciare "File with id 99 not found" & \ok \\
TU-122 & FileViewerIpcAdapter > registra tutti i channel IPC attesi & registeredChannels deve contenere IpcChannels.FILE\_DOWNLOAD & \ok \\
TU-123 & FileViewerIpcAdapter > FILE\_DOWNLOAD > chiama exportFileUC.execute con fileId & exportFileUC.execute deve essere stato chiamato con 1 & \ok \\
TU-124 & FileViewerIpcAdapter > FILE\_DOWNLOAD > ritorna ExportResult.ok() se l'esportazione riesce & result.success deve essere true & \ok \\
TU-125 & FileViewerIpcAdapter > FILE\_DOWNLOAD > ritorna ExportResult.fail() se il file non esiste & result.success deve essere false; result.errorCode deve essere "NOT\_FOUND" & \ok \\
TU-126 & FileViewerIpcAdapter > FILE\_DOWNLOAD > ritorna ExportResult.fail() se la scrittura fallisce & result.success deve essere false; result.errorCode deve essere "WRITE\_ERROR" & \ok \\
TU-127 & FileViewerIpcAdapter > FILE\_DOWNLOAD > propaga eccezioni di exportFileUC.execute & ipcMain.invoke(IpcChannels.FILE\_DOWNLOAD, 1) deve lanciare "export exploded" & \ok \\
TU-128 & SearchIpcAdapter > registra tutti i channel IPC attesi & registeredChannels deve contenere IpcChannels.SEARCH\_CLASSES; registeredChannels deve contenere IpcChannels.SEARCH\_PROCESSES; registeredChannels deve contenere IpcChannels.SEARCH\_DOCUMENTS; registeredChannels deve contenere IpcChannels.SEARCH\_SEMANTIC; registeredChannels deve contenere IpcChannels.SEARCH\_FULLTEXT; registeredChannels deve contenere IpcChannels.SEARCH\_GET\_AI\_STATE; registeredChannels deve contenere IpcChannels.SEARCH\_CUSTOM\_METADATA\_KEYS & \ok \\
TU-129 & SearchIpcAdapter > SEARCH\_CLASSES chiama execute con il nome e ritorna i DTO & searchClassesUC.execute deve essere stato chiamato con "Contratti"; result deve essere uguale a [DocumentClassMapper.toSearchResult(dc)] & \ok \\
TU-130 & SearchIpcAdapter > SEARCH\_CLASSES usa stringa vuota se il nome non è fornito & searchClassesUC.execute deve essere stato chiamato con "" & \ok \\
TU-131 & SearchIpcAdapter > SEARCH\_CLASSES ritorna array vuoto se nessuna classe trovata & result deve essere uguale a [] & \ok \\
TU-132 & SearchIpcAdapter > SEARCH\_CLASSES propaga eccezioni della use-case & ipcMain.invoke(IpcChannels.SEARCH\_CLASSES, "Contratti") deve lanciare "classes failed" & \ok \\
TU-133 & SearchIpcAdapter > SEARCH\_PROCESSES chiama execute con uuid e ritorna i DTO & searchProcessiUC.execute deve essere stato chiamato con "proc-uuid-abc"; result deve essere uguale a [ProcessMapper.toSearchResult(proc)] & \ok \\
TU-134 & SearchIpcAdapter > SEARCH\_PROCESSES usa stringa vuota se uuid non è fornito & searchProcessiUC.execute deve essere stato chiamato con "" & \ok \\
TU-135 & SearchIpcAdapter > SEARCH\_PROCESSES ritorna array vuoto se nessun processo trovato & result deve essere uguale a [] & \ok \\
TU-136 & SearchIpcAdapter > SEARCH\_PROCESSES propaga eccezioni della use-case & ipcMain.invoke(IpcChannels.SEARCH\_PROCESSES, "uuid-1") deve lanciare "processes failed" & \ok \\
TU-137 & SearchIpcAdapter > SEARCH\_DOCUMENTS chiama execute e propaga i risultati & searchDocumentsUC.execute deve essere stato chiamato; result deve essere uguale a [DocumentMapper.toSearchResult(doc, null)] & \ok \\
TU-138 & SearchIpcAdapter > SEARCH\_DOCUMENTS ritorna array vuoto se nessun documento trovato & result deve essere uguale a [] & \ok \\
TU-139 & SearchIpcAdapter > SEARCH\_DOCUMENTS propaga rejection della use-case & ipcMain.invoke(IpcChannels.SEARCH\_DOCUMENTS, filters) deve lanciare "documents failed" & \ok \\
TU-140 & SearchIpcAdapter > SEARCH\_SEMANTIC chiama execute con la query e ritorna SearchResult[] & searchSemanticUC.execute deve essere stato chiamato con "ricerca semantica"; result deve essere uguale a [DocumentMapper.toSearchResult(doc, 0.91)] & \ok \\
TU-141 & SearchIpcAdapter > SEARCH\_SEMANTIC ritorna array vuoto se nessun documento simile trovato & result deve essere uguale a [] & \ok \\
TU-142 & SearchIpcAdapter > SEARCH\_SEMANTIC passa stringa vuota quando la query è vuota & searchSemanticUC.execute deve essere stato chiamato con "" & \ok \\
TU-143 & SearchIpcAdapter > SEARCH\_SEMANTIC propaga rejection della use-case & ipcMain.invoke(IpcChannels.SEARCH\_SEMANTIC, query) deve lanciare "semantic failed" & \ok \\
TU-144 & SearchIpcAdapter > SEARCH\_FULLTEXT chiama searchDocumentsUC con filtro sul nome & calledQuery.filter.items[0] deve essere uguale a { path: "NomeDelDocumento", operator: "LIKE", value: "\%fattura.pdf\%", } & \ok \\
TU-145 & SearchIpcAdapter > SEARCH\_FULLTEXT ritorna Document[] dal searchDocumentsUC & result deve essere uguale a expectedResults.map((doc) => DocumentMapper.toSearchResult(doc, null)) & \ok \\
TU-146 & SearchIpcAdapter > SEARCH\_FULLTEXT ritorna array vuoto se nessun documento trovato & result deve essere uguale a [] & \ok \\
TU-147 & SearchIpcAdapter > SEARCH\_FULLTEXT propaga rejection della use-case & ipcMain.invoke(IpcChannels.SEARCH\_FULLTEXT, query) deve lanciare "fulltext failed" & \ok \\
TU-148 & SearchIpcAdapter > SEARCH\_GET\_AI\_STATE ritorna IDLE quando il modello non è caricato & result.status deve essere "IDLE"; result.progressPercentage deve essere 0; result.lastIndexedAt deve essere null; result.indexedDocuments deve essere 0 & \ok \\
TU-149 & SearchIpcAdapter > SEARCH\_GET\_AI\_STATE ritorna READY quando il modello è caricato & result.status deve essere "READY"; result.progressPercentage deve essere 100; result.lastIndexedAt deve essere null; result.indexedDocuments deve essere 42 & \ok \\
TU-150 & SearchIpcAdapter > SEARCH\_GET\_AI\_STATE riflette il conteggio aggiornato dei documenti indicizzati & result.indexedDocuments deve essere 150 & \ok \\
TU-151 & SearchIpcAdapter > SEARCH\_GET\_AI\_STATE propaga eccezioni durante la lettura dello stato AI & ipcMain.invoke(IpcChannels.SEARCH\_GET\_AI\_STATE) deve lanciare "ai state unavailable" & \ok \\
TU-152 & SearchIpcAdapter > SEARCH\_CUSTOM\_METADATA\_KEYS ritorna le chiavi custom distinte per dip & getCustomMetadataKeysUC.execute deve essere stato chiamato con 7; result deve essere uguale a ["DataProtocollo", "NomeCliente", "NumeroPratica"] & \ok \\
TU-153 & SearchIpcAdapter > SEARCH\_CUSTOM\_METADATA\_KEYS usa null quando dipId non e fornito & getCustomMetadataKeysUC.execute deve essere stato chiamato con null; result deve essere uguale a ["NomeCliente"] & \ok \\
TU-154 & DipPersistenceAdapter > TU-F-browsing-44: save() should save crea un dip con status UNKNOWN & dao.save deve essere stato chiamato con input; result deve essere savedDip & \ok \\
TU-155 & DipPersistenceAdapter > TU-F-browsing-45: getById() should getById e getByUuid restituiscono l'entita & dao.getById deve essere stato chiamato con 12; dao.getByUuid deve essere stato chiamato con "dip-2"; byId?.getUuid() deve essere "dip-2"; byUuid?.getId() deve essere 12 & \ok \\
TU-156 & DipPersistenceAdapter > TU-F-browsing-46: updateIntegrityStatus() should updateIntegrityStatus e getByStatus funzionano & dao.updateIntegrityStatus deve essere stato chiamato con 13, IntegrityStatusEnum.VALID; dao.getByStatus deve essere stato chiamato con IntegrityStatusEnum.VALID; found deve avere lunghezza 1; found[0].getUuid() deve essere "dip-3" & \ok \\
TU-157 & DipPersistenceAdapter > TU-F-browsing-47: getByStatus() should ritorna array vuoto se nessuna riga trovata per getByStatus & dao.getByStatus deve essere stato chiamato con IntegrityStatusEnum.VALID; found deve avere lunghezza 0 & \ok \\
TU-158 & DipPersistenceAdapter > TU-F-browsing-48: save() should save fallback if changes is 0 & dao.save deve essere stato chiamato con dip; saved.getId() deve essere 99 & \ok \\
TU-159 & DipPersistenceAdapter > TU-F-browsing-49: save() should save fallback if lastInsertRowid is 0 but changes > 0 & dao.save deve essere stato chiamato con dip; saved.getId() deve essere 88; saved.getUuid() deve essere "dip-update" & \ok \\
TU-160 & DocumentClassPersistenceAdapter > TU-F-browsing-50: save() should save e getById funzionano & dao.save deve essere stato chiamato con input; dao.getById deve essere stato chiamato con 21; found?.getDipId() deve essere 10; found?.getUuid() deve essere "dc-1" & \ok \\
TU-161 & DocumentClassPersistenceAdapter > TU-F-browsing-51: getByDipId() should getByDipId, getByStatus e updateIntegrityStatus funzionano & repo.getByDipId(20) deve avere lunghezza 1; dao.updateIntegrityStatus deve essere stato chiamato con 31, IntegrityStatusEnum.VALID; byStatus deve avere lunghezza 1 & \ok \\
TU-162 & DocumentClassPersistenceAdapter > TU-F-browsing-52: search() should search restituisce risultati o null & dao.search deve essere stato chiamato con "Verbali"; dao.search deve essere stato chiamato con "inesistente"; found non deve essere null; found?.[0].getName() deve contenere "Verbali"; notFound deve essere null & \ok \\
TU-163 & DocumentClassPersistenceAdapter > TU-F-browsing-53: save() should save fallback per lastInsertRowid falsy e exception per fallimento & () => repo.save(input).toThrowError("Failed to save DocumentClass with uuid=dc-err") & \ok \\
TU-164 & DocumentClassPersistenceAdapter > TU-F-browsing-54: getById() should return null quando la classe non esiste & dao.getById deve essere stato chiamato con 404; result deve essere null & \ok \\
TU-165 & DocumentPersistenceAdapter > TU-F-browsing-55: save() should successfully persist a document with complex metadata & dao.save deve essere stato chiamato con input; result deve essere saved & \ok \\
TU-166 & DocumentPersistenceAdapter > TU-S-browsing-56: save() should handle updating an existing document & result.getUuid() deve essere "doc-unique"; dao.save deve essere stato chiamato 1 volte & \ok \\
TU-167 & DocumentPersistenceAdapter > TU-F-browsing-57: getByProcessId() should return a list of documents & dao.getByProcessId deve essere stato chiamato con 1; results deve avere lunghezza 1; results[0].getUuid() deve essere "doc-p1" & \ok \\
TU-168 & DocumentPersistenceAdapter > TU-F-browsing-58: getByProcessId() should return an empty array & results deve avere lunghezza 0 & \ok \\
TU-169 & DocumentPersistenceAdapter > TU-F-browsing-59: getByStatus() should return documents & dao.getByStatus deve essere stato chiamato con IntegrityStatusEnum.VALID; results deve avere lunghezza 1 & \ok \\
TU-170 & DocumentPersistenceAdapter > TU-F-browsing-60: getByStatus() should return an empty array & repo.getByStatus(IntegrityStatusEnum.VALID) deve avere lunghezza 0 & \ok \\
TU-171 & DocumentPersistenceAdapter > TU-S-browsing-61: updateIntegrityStatus() should successfully update the status & dao.updateIntegrityStatus deve essere stato chiamato con 7, IntegrityStatusEnum.VALID & \ok \\
TU-172 & DocumentPersistenceAdapter > TU-S-browsing-62: save() should fallback to select & result.getId() deve essere 33 & \ok \\
TU-173 & DocumentPersistenceAdapter > TU-F-browsing-63: getById() should return null quando il documento non esiste & dao.getById deve essere stato chiamato con 999; result deve essere null & \ok \\
TU-174 & DocumentPersistenceAdapter > searchDocument delega al DAO con i filtri & dao.searchDocument deve essere stato chiamato con emptyFilters; results deve essere uguale a [] & \ok \\
TU-175 & DocumentPersistenceAdapter > searchDocumentSemantic delega al DAO & dao.searchDocumentSemantic deve essere stato chiamato con queryVector; result deve essere uguale a semantic & \ok \\
TU-176 & DocumentPersistenceAdapter > getIndexedDocumentsCount delega al DAO & dao.getIndexedDocumentsCount deve essere stato chiamato 1 volte; count deve essere 42 & \ok \\
TU-177 & EntityToDtoConverter > dipToDto > TU-F-converter-01: dipToDto() should convertire un Dip entity a DipDTO correttamente & dto.id deve essere 1; dto.uuid deve essere "test-uuid"; dto.integrityStatus deve essere IntegrityStatusEnum.VALID & \ok \\
TU-178 & EntityToDtoConverter > dipToDto > TU-F-converter-02: dipToDto() should lanciare errore se id è null & () => EntityToDtoConverter.dipToDto(dip) deve lanciare "Cannot convert Dip to DTO: id is null" & \ok \\
TU-179 & EntityToDtoConverter > dipToDto > keeps UNKNOWN integrity status & dto.integrityStatus deve essere IntegrityStatusEnum.UNKNOWN & \ok \\
TU-180 & EntityToDtoConverter > documentClassToDto > TU-F-converter-04: documentClassToDto() should convertire un DocumentClass entity a DocumentClassDTO correttamente & dto.id deve essere 2; dto.dipId deve essere 1; dto.uuid deve essere "dc-uuid"; dto.name deve essere "Ricevute"; dto.timestamp deve essere "2025-03-01T08:00:00Z"; dto.integrityStatus deve essere IntegrityStatusEnum.VALID & \ok \\
TU-181 & EntityToDtoConverter > documentClassToDto > TU-F-converter-05: documentClassToDto() should lanciare errore se id è null & () => EntityToDtoConverter.documentClassToDto(dc) deve lanciare "Cannot convert DocumentClass to DTO: id or dipId is null" & \ok \\
TU-182 & EntityToDtoConverter > documentClassToDto > TU-F-converter-06: documentClassToDto() should gestire INVALID integrity status & dto.integrityStatus deve essere IntegrityStatusEnum.INVALID & \ok \\
TU-183 & EntityToDtoConverter > documentToDto > TU-F-converter-07: documentToDto() should convertire un Document entity a DocumentDTO con metadati & dto.id deve essere 5; dto.processId deve essere 10; dto.uuid deve essere "doc-uuid"; dto.integrityStatus deve essere IntegrityStatusEnum.VALID; dto.metadata deve essere definito; dto.metadata.name deve essere "DocumentoInformatico"; dto.metadata.value deve avere lunghezza 2 & \ok \\
TU-184 & EntityToDtoConverter > documentToDto > TU-F-converter-08: documentToDto() should lanciare errore se id è null & () => EntityToDtoConverter.documentToDto(doc) deve lanciare "Cannot convert Document to DTO: id or processId is null" & \ok \\
TU-185 & EntityToDtoConverter > documentToDto > converts document with empty composite metadata & dto.metadata.name deve essere "root"; Array.isArray(dto.metadata.value) deve essere true; (dto.metadata.value as MetadataDTO[]).length deve essere 0 & \ok \\
TU-186 & EntityToDtoConverter > fileToDto > TU-F-converter-10: fileToDto() should convertire un File entity a FileDTO correttamente & dto.id deve essere 7; dto.documentId deve essere 5; dto.filename deve essere "documento.pdf"; dto.path deve essere "/docs/documento.pdf"; dto.hash deve essere "abc123hash"; dto.integrityStatus deve essere IntegrityStatusEnum.VALID; dto.isMain deve essere true & \ok \\
TU-187 & EntityToDtoConverter > fileToDto > TU-F-converter-11: fileToDto() should lanciare errore se id è null & () => EntityToDtoConverter.fileToDto(file) deve lanciare "Cannot convert File to DTO: id or documentId is null" & \ok \\
TU-188 & EntityToDtoConverter > fileToDto > converts non-main attachment file & dto.isMain deve essere false; dto.integrityStatus deve essere IntegrityStatusEnum.INVALID & \ok \\
TU-189 & EntityToDtoConverter > processToDto > TU-F-converter-13: processToDto() should convertire un Process entity a ProcessDTO con metadati & dto.id deve essere 10; dto.documentClassId deve essere 2; dto.uuid deve essere "proc-uuid"; dto.integrityStatus deve essere IntegrityStatusEnum.VALID; (dto.metadata.value as MetadataDTO[]).length deve essere 2; (dto.metadata.value[0] as MetadataDTO).name deve essere "fase"; (dto.metadata.value[0] as MetadataDTO).value deve essere "Acquisizione"; (dto.metadata.value[1] as MetadataDTO).name deve essere "ordine"; (dto.metadata.value[1] as MetadataDTO).value deve essere "1" & \ok \\
TU-190 & EntityToDtoConverter > processToDto > TU-F-converter-14: processToDto() should lanciare errore se id è null & () => EntityToDtoConverter.processToDto(proc) deve lanciare "Cannot convert Process to DTO: id or documentClassId is null" & \ok \\
TU-191 & EntityToDtoConverter > processToDto > converts process with empty metadata children & dto.metadata.name deve essere "ProcessoVuoto"; Array.isArray(dto.metadata.value) deve essere true; (dto.metadata.value as MetadataDTO[]).length deve essere 0 & \ok \\
TU-192 & EntityToDtoConverter > metadataToDto > TU-F-converter-16: metadataToDto() should convertire un Metadata semplice (STRING) a MetadataDTO & dto.name deve essere "autore"; dto.value deve essere "Mario Rossi"; dto.type deve essere MetadataType.STRING & \ok \\
TU-193 & EntityToDtoConverter > metadataToDto > TU-F-converter-17: metadataToDto() should convertire un Metadata numerico (NUMBER) a MetadataDTO & dto.name deve essere "anno"; dto.value deve essere "2025"; dto.type deve essere MetadataType.NUMBER & \ok \\
TU-194 & EntityToDtoConverter > metadataToDto > TU-F-converter-18: metadataToDto() should convertire un Metadata booleano (BOOLEAN) a MetadataDTO & dto.name deve essere "riservato"; dto.value deve essere "true"; dto.type deve essere MetadataType.BOOLEAN & \ok \\
TU-195 & EntityToDtoConverter > metadataToDto > converts COMPOSITE metadata recursively & dto.name deve essere "root"; dto.type deve essere MetadataType.COMPOSITE; Array.isArray(dto.value) deve essere true; (dto.value as any[]).length deve essere 2 & \ok \\
TU-196 & EntityToDtoConverter > metadataToDto > converts empty COMPOSITE metadata & dto.name deve essere "emptyRoot"; dto.type deve essere MetadataType.COMPOSITE; Array.isArray(dto.value) deve essere true; (dto.value as MetadataDTO[]).length deve essere 0 & \ok \\
TU-197 & EntityToDtoConverter > metadataToDto > converts deeply nested metadata & dto.name deve essere "level1"; level2.name deve essere "level2"; level3.name deve essere "level3"; deep.name deve essere "deep"; deep.value deve essere "value" & \ok \\
TU-198 & FilePersistenceAdapter > TU-F-browsing-64: save() should save e getById funzionano & dao.save deve essere stato chiamato con input; dao.getById deve essere stato chiamato con 71; found?.getFilename() deve essere "main.xml"; found?.getDocumentId() deve essere 3 & \ok \\
TU-199 & FilePersistenceAdapter > TU-F-browsing-65: getByDocumentId() should getByDocumentId, getByStatus e updateIntegrityStatus funzionano & repo.getByDocumentId(4) deve avere lunghezza 1; dao.updateIntegrityStatus deve essere stato chiamato con 72, IntegrityStatusEnum.INVALID; dao.getByStatus deve essere stato chiamato con IntegrityStatusEnum.INVALID; byStatus deve avere lunghezza 1 & \ok \\
TU-200 & FilePersistenceAdapter > TU-F-browsing-66: save() should save con fallback quando lastInsertRowid e falsy & dao.save deve essere stato chiamato con input; result.getId() deve essere 88 & \ok \\
TU-201 & FilePersistenceAdapter > TU-F-browsing-67: getById() should return null quando il file non esiste & dao.getById deve essere stato chiamato con 404; result deve essere null & \ok \\
TU-202 & LocalExportPort > should pipe data to file system & result.success deve essere true; fileContent deve essere data & \ok \\
TU-203 & LocalExportPort > should return error result on write failure & result.success deve essere false; result.errorCode deve essere "WRITE\_ERROR" & \ok \\
TU-204 & LocalPackageReaderAdapter > TU-F-Indexing-07: readDip() should throw if DiP index file is not found (Illegal path) & adapter.setDipPath(dipPath) deve lanciare "DiP index file not found in '/mock/dip/path'. Expected format: DiPIndex.<uuid>.xml" & \ok \\
TU-205 & LocalPackageReaderAdapter > TU-F-Indexing-08: readDip() should read the first alphabetically sorted DiP index file (Edge case) and return Dip & result deve essere un'istanza di Dip; result.getUuid() deve essere "dip-from-a" & \ok \\
TU-206 & LocalPackageReaderAdapter > TU-F-Indexing-09: readDocumentClasses() should yield mapped document classes & results deve avere lunghezza 2; results[0] deve essere un'istanza di DocumentClass; results[0].getDipUuid() deve essere "dip-uuid"; results[0].getUuid() deve essere "dc-1"; results[0].getName() deve essere "doc1"; results[1] deve essere un'istanza di DocumentClass; results[1].getDipUuid() deve essere "dip-uuid"; results[1].getUuid() deve essere "dc-2"; results[1].getName() deve essere "doc2" & \ok \\
TU-207 & LocalPackageReaderAdapter > TU-F-Indexing-10: readProcesses() should yield mapped processes (Legal metadata execution) & results deve avere lunghezza 1; results[0] deve essere un'istanza di Process; results[0].getDocumentClassUuid() deve essere "dc-uuid"; results[0].getUuid() deve essere "proc-meta"; results[0].getMetadata() deve essere uguale a new Metadata("meta", "value") & \ok \\
TU-208 & LocalPackageReaderAdapter > TU-F-Indexing-11: readProcesses() should handle missing metadata files gracefully by passing null (Edge case) & results deve avere lunghezza 1; results[0] deve essere un'istanza di Process; results[0].getDocumentClassUuid() deve essere "dc-uuid"; results[0].getUuid() deve essere "proc-null-metadata"; results[0].getMetadata() deve essere uguale a new Metadata("meta", "value") & \ok \\
TU-209 & LocalPackageReaderAdapter > TU-F-Indexing-12: readDocuments() should safely skip metadata fetch if relative path is null & results deve avere lunghezza 1; results[0] deve essere un'istanza di Document; results[0].getUuid() deve essere "doc-without-metadata"; results[0].getProcessUuid() deve essere "proc-uuid"; results[0].getMetadata() deve essere uguale a new Metadata("meta", "value") & \ok \\
TU-210 & LocalPackageReaderAdapter > TU-F-Indexing-14: readFiles() should fetch paths correctly from file mappers & results deve avere lunghezza 1; results[0] deve essere un'istanza di File; results[0].getFilename() deve essere "some-file"; results[0].getPath() deve essere "path/without-meta"; results[0].getHash() deve essere "hash"; results[0].getIsMain() deve essere true; results[0].getDocumentUuid() deve essere "doc-uuid" & \ok \\
TU-211 & PrintPort > crea una BrowserWindow nascosta con plugins abilitati & BrowserWindow deve essere stato chiamato con { show: false, webPreferences: { plugins: true }, } & \ok \\
TU-212 & PrintPort > carica il file con il protocollo file:// & mockLoadURL deve essere stato chiamato con "file:///tmp/doc.pdf" & \ok \\
TU-213 & PrintPort > chiama print con le opzioni fornite dopo did-finish-load & mockPrint deve essere stato chiamato con opts, any Function & \ok \\
TU-214 & PrintPort > risolve { success: true } quando la stampa riesce & result deve essere uguale a { success: true } & \ok \\
TU-215 & PrintPort > distrugge la finestra dopo una stampa riuscita & mockDestroy.toHaveBeenCalledOnce() & \ok \\
TU-216 & PrintPort > risolve { success: false, error } quando la stampa fallisce & result.success deve essere false; result.errorMessage deve essere "printer not found"; result.errorCode deve essere "PRINT\_ERROR" & \ok \\
TU-217 & PrintPort > distrugge la finestra anche quando la stampa fallisce & mockDestroy.toHaveBeenCalledOnce() & \ok \\
TU-218 & PrintPort > risolve { success: false, error } con il messaggio corretto su did-fail-load & result.success deve essere false; result.errorMessage deve essere "did-fail-load: ERR\_FILE\_NOT\_FOUND (codice: -6)"; result.errorCode deve essere "LOAD\_ERROR" & \ok \\
TU-219 & PrintPort > distrugge la finestra su did-fail-load & mockDestroy.toHaveBeenCalledOnce() & \ok \\
TU-220 & PrintPort > non chiama print se il caricamento fallisce & mockPrint non deve essere stato chiamato & \ok \\
TU-221 & ProcessPersistenceAdapter > TU-F-browsing-68: save() should save persiste processo e metadata & dao.save deve essere stato chiamato con input; result.getId() deve essere 81; result.getDocumentClassId() deve essere 11 & \ok \\
TU-222 & ProcessPersistenceAdapter > TU-F-browsing-69: getByDocumentClassId() should getByDocumentClassId, getByStatus e updateIntegrityStatus funzionano & repo.getByDocumentClassId(12) deve avere lunghezza 1; dao.updateIntegrityStatus deve essere stato chiamato con 82, IntegrityStatusEnum.VALID; dao.getByStatus deve essere stato chiamato con IntegrityStatusEnum.VALID; byStatus deve avere lunghezza 1; byStatus[0].getUuid() deve essere "proc-2" & \ok \\
TU-223 & searchDocumentSemantic — delega al DAO > propaga la query invariata al DAO & dao.searchDocumentSemantic deve essere stato chiamato con queryVector; dao.searchDocumentSemantic deve essere stato chiamato 1 volte & \ok \\
TU-224 & searchDocumentSemantic — delega al DAO > ritorna i risultati semantici senza alterazioni & result deve essere semanticResults; result deve avere lunghezza 2; result[0].score deve essere 0.91; result[1].document.getUuid() deve essere "doc-2" & \ok \\
TU-225 & searchDocumentSemantic — delega al DAO > propaga errori sollevati dal DAO & repo.searchDocumentSemantic(new Float32Array([0.7, 0.8, 0.9])) deve lanciare "semantic failed" & \ok \\
TU-226 & searchDocumentSemantic — delega al DAO > supporta chiamate multiple indipendenti & first deve avere lunghezza 1; second deve avere lunghezza 0; dao.searchDocumentSemantic deve essere stato chiamato 2 volte & \ok \\
TU-227 & SearchSemanticUC — calcolo score > score è 1.0 quando la distanza è 0 (corrispondenza perfetta) & results[0].score deve essere 1 & \ok \\
TU-228 & SearchSemanticUC — calcolo score > score è 0.0 quando la distanza è massima (nessuna similarità) & results[0].score deve essere 0 & \ok \\
TU-229 & SearchSemanticUC — calcolo score > preserva l'ordine decrescente per score restituito dal repository & scores[i] deve essere maggiore o uguale a scores[i + 1] & \ok \\
TU-230 & SearchSemanticUC — calcolo score > preserva la precisione decimale dello score & results[0].score deve essere vicino a 0.123456789, 9 & \ok \\
TU-231 & SearchSemanticUC — calcolo score > score nella fascia alta (> 0.8) indica alta rilevanza semantica & results[0].score deve essere maggiore di 0.8 & \ok \\
TU-232 & SearchSemanticUC — calcolo score > score nella fascia bassa (< 0.5) indica bassa rilevanza semantica & results[0].score deve essere minore di 0.5 & \ok \\
TU-233 & SearchSemanticUC — calcolo score > non restituisce più di 10 risultati (limite del repository) & results.length deve essere minore o uguale a 10 & \ok \\
TU-234 & SearchSemanticUC — calcolo score > tutti gli score sono nel range valido [0.0, 1.0] & r.score deve essere maggiore o uguale a 0; r.score deve essere minore o uguale a 1 & \ok \\
TU-235 & SearchSemanticUC — calcolo score > lancia un errore se il documentId restituito dal vectorRepo non esiste & uc.execute("missing doc") deve lanciare "999" & \ok \\
TU-236 & SearchSemanticUC — calcolo score > passa il vettore generato dall'adapter al vectorRepo & aiAdapter.generateEmbedding deve essere stato chiamato con "test vettore"; vectorRepo.searchSimilarVectors deve essere stato chiamato con vector, 10 & \ok \\
TU-237 & SqliteTransactionManager > starts BEGIN IMMEDIATE TRANSACTION and COMMIT on success & result deve essere "ok"; db.exec.mock.calls deve essere uguale a [ ["BEGIN IMMEDIATE TRANSACTION"], ["COMMIT"], ] & \ok \\
TU-238 & SqliteTransactionManager > starts BEGIN IMMEDIATE TRANSACTION and ROLLBACK on failure & manager.runInTransaction(async () => { throw error; }) deve essere error; db.exec.mock.calls deve essere uguale a [ ["BEGIN IMMEDIATE TRANSACTION"], ["ROLLBACK"], ] & \ok \\
TU-239 & SqliteTransactionManager > propagates caller values without altering them & result deve essere payload & \ok \\
TU-240 & DataMapper > TU-F-Indexing-136: mapDip() returns correctly mapped Dip entity from dipindex\_test.xml structure & result.getUuid() deve essere "test-dip-uuid-1234" & \ok \\
TU-241 & DataMapper > TU-F-Indexing-137: mapDocumentClass() returns correctly mapped DocumentClass entity from dipindex\_test.xml structure & result[0].getUuid() deve essere "class-1"; result[0].getName() deve essere "Class1" & \ok \\
TU-242 & DataMapper > TU-F-Indexing-138: getProcessMappers() should return correctly mapped paths for processes in the dipindex\_test.xml structure & result[0].metadataRelativePath deve essere "./class-1/aip-1/AiPInfo.aip-1.xml"; result[1].metadataRelativePath deve essere "./class-2/aip-2/AiPInfo.aip-2.xml" & \ok \\
TU-243 & DataMapper > TU-F-Indexing-138: getProcessMappers() should map processes from dipindex\_test.xml structure correctly & proc.getUuid() deve essere "aip-1"; proc.getDocumentClassUuid() deve essere "class-1"; processMeta deve essere definito; startMeta deve essere definito; startMeta.getChildren() deve essere uguale a [ expect.objectContaining({ name: "Date", value: "2025-01-01T00:00:00Z", }), expect.objectContaining({ name: "UserUUID", value: "dummy-user-uuid" }), expect.objectContaining({ name: "Source", value: "DummySource" }), ]; endMeta deve essere definito; endMeta.getChildren() deve essere uguale a [ expect.objectContaining({ name: "Date", value: "2025-01-01T01:00:00Z", }), expect.objectContaining({ name: "UserUUID", value: "dummy-user-uuid" }), expect.objectContaining({ name: "Source", value: "DummySource" }), expect.objectContaining({ name: "Status", value: "Completed" }), ]; submissionMeta deve essere definito; submissionStart?.getChildren() deve essere uguale a [ expect.objectContaining({ name: "Date", value: "2025-01-01T00:00:00Z", }), expect.objectContaining({ name: "UserUUID", value: "dummy-user-uuid" }), expect.objectContaining({ name: "Source", value: "DummySource" }), ]; submissionEnd?.getChildren() deve essere uguale a [ expect.objectContaining({ name: "Date", value: "2025-01-01T00:30:00Z", }), expect.objectContaining({ name: "UserUUID", value: "dummy-user-uuid" }), expect.objectContaining({ name: "Source", value: "DummySource" }), expect.objectContaining({ name: "Status", value: "Completed" }), ]; preservationMeta deve essere definito; preservationStart?.getChildren() deve essere uguale a [ expect.objectContaining({ name: "Date", value: "2025-01-01T00:45:00Z", }), expect.objectContaining({ name: "UserUUID", value: "dummy-user-uuid" }), expect.objectContaining({ name: "Source", value: "DummySource" }), ]; preservationEnd.getChildren() deve essere definito & \ok \\
TU-244 & DataMapper > TU-F-Indexing-139: getProcessMappers() should handle missing optional metadata gracefully & proc.getUuid() deve essere "aip-1"; proc.getDocumentClassUuid() deve essere "class-1"; proc.getMetadata() deve essere uguale a new Metadata("AiPInfo", [], MetadataType.COMPOSITE) & \ok \\
TU-245 & DataMapper > TU-F-Indexing-140: getProcessMappers() should handle unexpected metadata structures gracefully & proc.getUuid() deve essere "aip-1"; proc.getDocumentClassUuid() deve essere "class-1"; proc.getMetadata() deve essere uguale a new Metadata("AiPInfo", [], MetadataType.COMPOSITE) & \ok \\
TU-246 & DataMapper > TU-F-Indexing-141: getDocumentMappers() should map documents correctly & doc.getUuid() deve essere "doc-1"; doc.getProcessUuid() deve essere "aip-1"; doc.getIntegrityStatus() deve essere "UNKNOWN"; metadata.getName() deve essere "DocumentoInformatico"; metadata.getChildren() deve avere lunghezza 6; metadata.getChildren()[0].getName() deve essere "IdDoc"; metadata.getChildren()[0].getChildren() deve avere lunghezza 2; metadata.getChildren()[0].getChildren()[0].getName() deve essere "ImprontaCrittograficaDelDocumento"; metadata.getChildren()[0].getChildren()[0].getChildren() deve avere lunghezza 2; metadata.getChildren()[0].getChildren()[0].getChildren()[0].getName() deve essere "Impronta"; metadata.getChildren()[0].getChildren()[0].getChildren()[1].getName() deve essere "Algoritmo"; metadata.getChildren()[0].getChildren()[1].getName() deve essere "Identificativo"; metadata.getChildren()[1].getName() deve essere "Riservato"; metadata.getChildren()[2].getName() deve essere "TempoDiConservazione"; metadata.getChildren()[3].getName() deve essere "Soggetti"; metadata.getChildren()[3].getChildren() deve avere lunghezza 2; metadata.getChildren()[3].getChildren()[0].getName() deve essere "Ruolo"; metadata.getChildren()[3].getChildren()[0].getChildren()[0] deve essere uguale a new Metadata("Tipo", "Autore", MetadataType.STRING); metadata.getChildren()[4].getName() deve essere "Titolo"; metadata.getChildren()[4] deve essere uguale a new Metadata("Titolo", "Documento di test", MetadataType.STRING) & \ok \\
TU-247 & DataMapper > TU-F-Indexing-142: getDocumentMappers() should return correctly mapped paths for documents in the dipindex\_test.xml structure & result[0].metadataRelativePath deve essere "./class-1/aip-1/documents/doc-1/meta.xml"; result[1].metadataRelativePath deve essere "./class-2/aip-2/documents/doc-2/meta2.xml" & \ok \\
TU-248 & DataMapper > TU-F-Indexing-143: getDocumentMappers() should handle missing optional metadata gracefully & doc.getUuid() deve essere "doc-1"; doc.getProcessUuid() deve essere "aip-1"; doc.getIntegrityStatus() deve essere "UNKNOWN"; doc.getMetadata().getName() deve essere "Unknown"; doc.getMetadata().getChildren() deve avere lunghezza 0 & \ok \\
TU-249 & DataMapper > TU-F-Indexing-144: getFileMappers() should return correctly mapped paths for files in the dipindex\_test.xml structure & result deve avere lunghezza 5; metadataPaths.filter( (p) => p === "./class-1/aip-1/documents/doc-1/meta.xml", ) deve avere lunghezza 1; metadataPaths.filter( (p) => p === "./class-2/aip-2/documents/doc-2/meta2.xml", ) deve avere lunghezza 3; metadataPaths.filter( (p) => p === "./class-2/aip-3/documents/doc-3/meta3.xml", ) deve avere lunghezza 1 & \ok \\
TU-250 & DataMapper > TU-F-Indexing-145: getFileMappers() should map all files correctly & mappedFiles deve avere lunghezza 5; mappedFiles[0].getUuid() deve essere "prim-1"; mappedFiles[0].getFilename() deve essere "./primary.pdf"; mappedFiles[0].getPath() deve essere "class-1/aip-1/documents/doc-1/primary.pdf"; mappedFiles[0].getHash() deve essere "hash-main-prim-1"; mappedFiles[0].getIsMain() deve essere true; mappedFiles[0].getDocumentId() deve essere null; mappedFiles[1].getUuid() deve essere "prim-2"; mappedFiles[1].getHash() deve essere ""; mappedFiles[1].getIsMain() deve essere true; mappedFiles[2].getUuid() deve essere "prim-3"; mappedFiles[2].getHash() deve essere ""; mappedFiles[2].getIsMain() deve essere true; mappedFiles[3].getUuid() deve essere "att-1"; mappedFiles[3].getFilename() deve essere "./att1.pdf"; mappedFiles[3].getPath() deve essere "class-2/aip-2/documents/doc-2/att1.pdf"; mappedFiles[3].getHash() deve essere "hash-att-1"; mappedFiles[3].getIsMain() deve essere false; mappedFiles[4].getUuid() deve essere "att-2"; mappedFiles[4].getFilename() deve essere "./att2.pdf"; mappedFiles[4].getPath() deve essere "class-2/aip-2/documents/doc-2/att2.pdf"; mappedFiles[4].getHash() deve essere "hash-att-2"; mappedFiles[4].getIsMain() deve essere false & \ok \\
TU-251 & ensureArray > TU-F-browsing-76: ensureArray() should return an empty array for undefined & ensureArray(undefined) deve essere uguale a [] & \ok \\
TU-252 & ensureArray > TU-F-browsing-77: ensureArray() should return an empty array for null & ensureArray(null) deve essere uguale a [] & \ok \\
TU-253 & ensureArray > TU-F-browsing-78: ensureArray() should return the same array if an array is passed & ensureArray(input) deve essere input & \ok \\
TU-254 & ensureArray > TU-F-browsing-79: ensureArray() should wrap a single non-array value in an array & ensureArray(1) deve essere uguale a [1]; ensureArray("test") deve essere uguale a ["test"]; ensureArray({ a: 1 }) deve essere uguale a [{ a: 1 }] & \ok \\
TU-255 & FileSystemProvider > TU-F-browsing-70: readFile() should return Uint8Array when reading an existing file from disk & result deve essere un'istanza di Uint8Array; Buffer.from(result).toString("utf-8") deve essere "test content" & \ok \\
TU-256 & FileSystemProvider > TU-F-browsing-71: openReadStream() should return a readable stream that reads actual file content & content deve essere "test content" & \ok \\
TU-257 & FileSystemProvider > TU-F-browsing-72: readTextFile() should read utf-8 text from an existing file on disk & result deve essere "test content" & \ok \\
TU-258 & FileSystemProvider > TU-F-browsing-73: fileExists() should return true for an existing path on disk & result deve essere true & \ok \\
TU-259 & FileSystemProvider > TU-F-browsing-74: fileExists() should return false for a non-existing path on disk & result deve essere false & \ok \\
TU-260 & FileSystemProvider > TU-F-browsing-75: listFiles() should return directory contents from the real filesystem & result.sort((a: string, b: string) => a.localeCompare(b)) deve essere uguale a [ "file1.txt", "file2.txt", ] & \ok \\
TU-261 & XmlDipParser > TU-F-indexing-80: parse() should correctly parse metadata from XML structure & parsedObj.simpleNode deve essere "exampleText"; parsedObj.nestedNode.childNode deve essere uguale a ["childText", "child2Text"] & \ok \\
TU-262 & XmlDipParser > TU-F-indexing-81: parse() should return an empty object for empty XML & parsedObj deve essere uguale a {} & \ok \\
TU-263 & XmlDipParser > TU-F-indexing-82: parse() should handle XML with attributes & parsedObj.node["\#text"] deve essere "text" & \ok \\
TU-264 & VectorPersistenceAdapter > should call DAO save() method without altering the Vector object & mockDAO.save deve essere stato chiamato con vector & \ok \\
TU-265 & VectorPersistenceAdapter > should throw error on null Vector object & repository.saveVector(null) deve lanciare un errore & \ok \\
TU-266 & VectorPersistenceAdapter > should call dao.getByDocumentId() and return embedding & mockDAO.getByDocumentId deve essere stato chiamato con 1; result deve essere uguale a new Float32Array([0.1, 0.2, 0.3]) & \ok \\
TU-267 & VectorPersistenceAdapter > should call dao.getByDocumentId() and return null if vector not found & mockDAO.getByDocumentId deve essere stato chiamato con 1; result deve essere null & \ok \\
TU-268 & VectorPersistenceAdapter > should call dao.getVector with correct documentId and return results & mockDAO.getByDocumentId.toBeCalledWith(1); result deve essere uguale a expectedVector & \ok \\
TU-269 & VectorPersistenceAdapter > should call dao.searchSimilar with correct parameters and return results & mockDAO.searchSimilar.toBeCalledWith(queryVector, topK); result deve essere uguale a [ { documentId: 1, score: 0.9 }, { documentId: 2, score: 0.8 }, ] & \ok \\
TU-270 & download-models — ensureDirectories > crea la cartella onnx se non esiste & fs.mkdirSync deve essere stato chiamato con string containing 'paraphrase-multilingual-MiniLM-L12-v2', { recursive: true } & \ok \\
TU-271 & download-models — ensureDirectories > non ricrea la cartella se esiste già & fs.mkdirSync non deve essere stato chiamato & \ok \\
TU-272 & download-models — ensureDirectories > il path contiene Xenova, il nome del modello e onnx & mkdirCall deve contenere 'Xenova'; mkdirCall deve contenere 'paraphrase-multilingual-MiniLM-L12-v2'; mkdirCall deve contenere 'onnx' & \ok \\
TU-273 & download-models — ensureDirectories > ritorna il path base che contiene Xenova & baseDir deve contenere 'Xenova' & \ok \\
TU-274 & download-models — salto file già presenti > legge la dimensione del file esistente con statSync & fs.existsSync deve essere stato chiamato & \ok \\
TU-275 & download-models — selezione protocollo > usa https per URL https:// & https.get deve essere stato chiamato; http.get non deve essere stato chiamato & \ok \\
TU-276 & download-models — selezione protocollo > crea il WriteStream nel path corretto & fs.createWriteStream deve essere stato chiamato con string containing 'model.onnx' & \ok \\
TU-277 & download-models — download riuscito > chiama pipe sulla response per scrivere il file & capturedRes.pipe deve essere stato chiamato con ws & \ok \\
TU-278 & download-models — download riuscito > chiude il WriteStream dopo il finish & ws.close deve essere stato chiamato & \ok \\
TU-279 & download-models — download riuscito > risolve la promise con il path completo del file & result deve contenere 'model.onnx'; result deve contenere '/tmp' & \ok \\
TU-280 & download-models — gestione redirect > segue un redirect 302 verso il nuovo URL & calls deve avere lunghezza 2; calls[0] deve essere 'https://example.com/model.onnx'; calls[1] deve essere 'https://cdn.example.com/model.onnx' & \ok \\
TU-281 & download-models — gestione redirect > distrugge il WriteStream prima di seguire il redirect & ws.destroy deve essere stato chiamato & \ok \\
TU-282 & download-models — gestione errori HTTP > rigetta la promise con errore HTTP 404 & downloadFile('https://example.com/model.onnx', '/tmp', 'model.onnx') deve lanciare 'Errore HTTP 404' & \ok \\
TU-283 & download-models — gestione errori HTTP > rigetta la promise con errore HTTP 500 & downloadFile('https://example.com/model.onnx', '/tmp', 'model.onnx') deve lanciare 'Errore HTTP 500' & \ok \\
TU-284 & download-models — gestione errori HTTP > cancella il file parziale su errore HTTP & fs.unlink deve essere stato chiamato; ws.destroy deve essere stato chiamato & \ok \\
TU-285 & download-models — gestione errori HTTP > rigetta la promise su errore di rete & downloadFile('https://example.com/model.onnx', '/tmp', 'model.onnx') deve lanciare 'ECONNREFUSED' & \ok \\
TU-286 & download-models — errori WriteStream > cancella il file parziale se WriteStream emette errore & fs.unlink deve essere stato chiamato & \ok \\
TU-287 & download-models — errori WriteStream > rigetta la promise se WriteStream emette errore & downloadFile('https://example.com/model.onnx', '/tmp', 'model.onnx') deve lanciare 'ENOSPC' & \ok \\
TU-288 & EmbeddingService > limits streamed text size to avoid huge string allocations & typeof text deve essere "string"; text.length deve essere maggiore di 0; text.length deve essere minore di 10\_000\_000; fileSystemProvider.openReadStream deve essere stato chiamato 1 volte & \ok \\
TU-289 & EmbeddingService > generates embedding for small streamed content & embedding deve essere un'istanza di Float32Array; embeddingAdapter.generateEmbedding deve essere stato chiamato & \ok \\
TU-290 & EmbeddingService > caps embedding chunk count for large content & embedding deve essere un'istanza di Float32Array; embeddingAdapter.generateEmbedding deve essere stato chiamato 8 volte & \ok \\
TU-291 & HashingService > TU-S-browsing-57: checkFileIntegrity() should return VALID for a file with correct hash & result deve essere "VALID" & \ok \\
TU-292 & HashingService > TU-S-browsing-58: checkFileIntegrity() should return INVALID for a file with incorrect hash & result deve essere "INVALID" & \ok \\
TU-293 & HashingService > TU-S-browsing-59: checkFileIntegrity() should handle empty files & result deve essere "VALID" & \ok \\
TU-294 & IntegrityVerificationService > throws when file root is not found and uses transaction wrapper & h.service.checkFileIntegrityStatus(777) deve lanciare "File with id 777 not found"; h.transactionManager.runInTransaction deve essere stato chiamato 1 volte & \ok \\
TU-295 & IntegrityVerificationService > throws when document root is not found and uses transaction wrapper & h.service.checkDocumentIntegrityStatus(778) deve lanciare "Document with id 778 not found"; h.transactionManager.runInTransaction deve essere stato chiamato 1 volte & \ok \\
TU-296 & IntegrityVerificationService > throws when process root is not found and uses transaction wrapper & h.service.checkProcessIntegrityStatus(779) deve lanciare "Process with id 779 not found"; h.transactionManager.runInTransaction deve essere stato chiamato 1 volte & \ok \\
TU-297 & IntegrityVerificationService > throws when document class root is not found and uses transaction wrapper & h.service.checkDocumentClassIntegrityStatus(780) deve lanciare "DocumentClass with id 780 not found"; h.transactionManager.runInTransaction deve essere stato chiamato 1 volte & \ok \\
TU-298 & IntegrityVerificationService > throws when dip root is not found and uses transaction wrapper & h.service.checkDipIntegrityStatus(781) deve lanciare "Dip with id 781 not found"; h.transactionManager.runInTransaction deve essere stato chiamato 1 volte & \ok \\
TU-299 & IntegrityVerificationService > checkFileIntegrityStatus returns status and updates only file level & status deve essere IntegrityStatusEnum.VALID; h.fileRepo.updateIntegrityStatus deve essere stato chiamato con 1, IntegrityStatusEnum.VALID; h.documentRepo.updateIntegrityStatus non deve essere stato chiamato; h.processRepo.updateIntegrityStatus non deve essere stato chiamato; h.documentClassRepo.updateIntegrityStatus non deve essere stato chiamato; h.dipRepo.updateIntegrityStatus non deve essere stato chiamato & \ok \\
TU-300 & IntegrityVerificationService > checkDocumentIntegrityStatus returns status and updates document level & status deve essere IntegrityStatusEnum.VALID; h.documentRepo.updateIntegrityStatus deve essere stato chiamato con 10, IntegrityStatusEnum.VALID; h.fileRepo.updateIntegrityStatus deve essere stato chiamato 2 volte & \ok \\
TU-301 & IntegrityVerificationService > checkProcessIntegrityStatus returns status and updates process level & status deve essere IntegrityStatusEnum.VALID; h.processRepo.updateIntegrityStatus deve essere stato chiamato con 20, IntegrityStatusEnum.VALID; h.documentRepo.updateIntegrityStatus deve essere stato chiamato con 21, IntegrityStatusEnum.VALID; h.fileRepo.updateIntegrityStatus deve essere stato chiamato con 2101, IntegrityStatusEnum.VALID & \ok \\
TU-302 & IntegrityVerificationService > checkDocumentClassIntegrityStatus returns status and updates class level & status deve essere IntegrityStatusEnum.VALID; h.documentClassRepo.updateIntegrityStatus deve essere stato chiamato con 30, IntegrityStatusEnum.VALID; h.processRepo.updateIntegrityStatus deve essere stato chiamato con 31, IntegrityStatusEnum.VALID; h.documentRepo.updateIntegrityStatus deve essere stato chiamato con 32, IntegrityStatusEnum.VALID; h.fileRepo.updateIntegrityStatus deve essere stato chiamato con 3201, IntegrityStatusEnum.VALID & \ok \\
TU-303 & IntegrityVerificationService > checkDipIntegrityStatus with full 2x2x2x2 valid tree marks every level as VALID & status deve essere IntegrityStatusEnum.VALID; h.transactionManager.runInTransaction deve essere stato chiamato 1 volte; fileValidUpdates deve avere lunghezza 16; documentValidUpdates deve avere lunghezza 8; processValidUpdates deve avere lunghezza 4; documentClassValidUpdates deve avere lunghezza 2; h.dipRepo.updateIntegrityStatus deve essere stato chiamato con 1, IntegrityStatusEnum.VALID & \ok \\
TU-304 & IntegrityVerificationService > checkDipIntegrityStatus propagates INVALID from one file to all its parents & status deve essere IntegrityStatusEnum.INVALID; h.documentRepo.updateIntegrityStatus deve essere stato chiamato con hierarchy.invalidChain.documentId, IntegrityStatusEnum.INVALID; h.processRepo.updateIntegrityStatus deve essere stato chiamato con hierarchy.invalidChain.processId, IntegrityStatusEnum.INVALID; h.documentClassRepo.updateIntegrityStatus deve essere stato chiamato con hierarchy.invalidChain.documentClassId, IntegrityStatusEnum.INVALID; h.dipRepo.updateIntegrityStatus deve essere stato chiamato con 1, IntegrityStatusEnum.INVALID & \ok \\
TU-305 & IntegrityVerificationService > checkDipIntegrityStatus returns VALID when there are no verifiable children & status deve essere IntegrityStatusEnum.VALID; h.dipRepo.updateIntegrityStatus deve essere stato chiamato con 99, IntegrityStatusEnum.VALID; h.documentClassRepo.updateIntegrityStatus non deve essere stato chiamato; h.processRepo.updateIntegrityStatus non deve essere stato chiamato; h.documentRepo.updateIntegrityStatus non deve essere stato chiamato; h.fileRepo.updateIntegrityStatus non deve essere stato chiamato & \ok \\
TU-306 & genera un Float32Array di 384 dimensioni & result deve essere un'istanza di Float32Array; result.length deve essere 384 & \ok \\
TU-307 & il vettore è normalizzato — norma ≈ 1.0 & norm deve essere vicino a 1, 2 & \ok \\
TU-308 & tutti i valori del vettore sono nel range [-1, 1] & v deve essere maggiore o uguale a -1; v deve essere minore o uguale a 1 & \ok \\
TU-309 & testi diversi producono vettori diversi & areDifferent deve essere true & \ok \\
TU-310 & lo stesso testo produce sempre lo stesso vettore & v1[i] deve essere vicino a v2[i], 5 & \ok \\
TU-311 & frasi semanticamente simili hanno similarità > 0.7 & cosineSimilarity(v1, v2) deve essere maggiore di 0.7 & \ok \\
TU-312 & frasi semanticamente diverse hanno similarità < 0.5 & cosineSimilarity(v1, v2) deve essere minore di 0.5 & \ok \\
TU-313 & la stessa frase ha similarità = 1.0 con se stessa & cosineSimilarity(v, v) deve essere vicino a 1, 4 & \ok \\
TU-314 & frasi in italiano simili hanno similarità alta & cosineSimilarity(v1, v2) deve essere maggiore di 0.75 & \ok \\
TU-315 & documenti amministrativi simili hanno similarità > 0.7 & cosineSimilarity(v1, v2) deve essere maggiore di 0.7 & \ok \\
TU-316 & la similarità è simmetrica — sim(a,b) = sim(b,a) & cosineSimilarity(v1, v2) deve essere vicino a cosineSimilarity(v2, v1), 5 & \ok \\
TU-317 & testo multilingua — italiano e inglese simili hanno similarità > 0.6 & cosineSimilarity(v1, v2) deve essere maggiore di 0.6 & \ok \\
TU-318 & score come 1 - distanza è nel range [0, 1] & score deve essere maggiore o uguale a 0; score deve essere minore o uguale a 1 & \ok \\
TU-319 & gestisce testo molto lungo senza errori & result deve essere un'istanza di Float32Array; result.length deve essere 384 & \ok \\
TU-320 & gestisce stringa vuota senza errori & result deve essere un'istanza di Float32Array; result.length deve essere 384 & \ok \\
TU-321 & gestisce caratteri speciali e unicode & result deve essere un'istanza di Float32Array; result.length deve essere 384 & \ok \\
TU-322 & gestisce numeri e codici & result deve essere un'istanza di Float32Array; result.length deve essere 384 & \ok \\
TU-323 & genera un embedding reale dalla query e lo passa al vectorRepo come Float32Array & vectorArg deve essere un'istanza di Float32Array; vectorArg.length deve essere 384 & \ok \\
TU-324 & embedding di query simili producono vettori diversi ma vicini & v1 non deve essere uguale a v2; cosineSimilarity(v1, v2) deve essere maggiore di 0.7 & \ok \\
TU-325 & embedding di query diverse producono vettori semanticamente distanti & cosineSimilarity(v1, v2) deve essere minore di 0.5 & \ok \\
TU-326 & flusso completo SearchSemanticUC — embedding reale e risultati vectorRepo passano invariati & results deve avere lunghezza 2; results[0].score deve essere vicino a 0.9, 5; results[1].score deve essere vicino a 0.6, 5; results[0].document deve essere fakeDoc1; results[1].document deve essere fakeDoc2 & \ok \\
TU-327 & query identica genera embedding identico & Array.from(v1) deve essere uguale a Array.from(v2) & \ok \\
TU-328 & flusso completo con query in inglese — modello multilingua & results deve avere lunghezza 1; results[0].score deve essere vicino a 0.85, 5 & \ok \\
TU-329 & vettore generato da embedding reale non è nullo o tutto zero & hasNonZero deve essere true & \ok \\
TU-330 & lancia un errore se vectorRepo restituisce un documentId inesistente & uc.execute("query senza documenti") deve lanciare "999" & \ok \\
TU-331 & WordEmbedding > isInitialized() ritorna false prima di qualsiasi chiamata & adapter.isInitialized() deve essere false & \ok \\
TU-332 & WordEmbedding > isInitialized() ritorna true dopo initialize() & adapter.isInitialized() deve essere true & \ok \\
TU-333 & WordEmbedding > chiama pipeline con i parametri corretti del modello & pipeline deve essere stato chiamato con 'feature-extraction', 'paraphrase-multilingual-MiniLM-L12-v2', { quantized: true, local\_files\_only: true } & \ok \\
TU-334 & WordEmbedding > initialize() è idempotente — pipeline non viene richiamata & pipeline deve essere stato chiamato 1 volte & \ok \\
TU-335 & WordEmbedding > configura transformers env per usare modelli locali su CPU & env.allowLocalModels deve essere true; env.allowRemoteModels deve essere false; env.useBrowserCache deve essere false; env.backends.onnx.executionProviders deve essere uguale a ['cpu'] & \ok \\
TU-336 & WordEmbedding > lancia errore con path se la cartella modelli non esiste & adapter.initialize() deve lanciare '[AI] Cartella modelli non trovata' & \ok \\
TU-337 & WordEmbedding > non modifica isInitialized se initialize() fallisce & adapter.isInitialized() deve essere false & \ok \\
TU-338 & WordEmbedding > propaga errore di pipeline e mantiene initialized=false & adapter.initialize() deve lanciare 'model load failed'; adapter.isInitialized() deve essere false & \ok \\
TU-339 & WordEmbedding > consente retry di initialize dopo errore precedente & adapter.initialize() deve lanciare 'temporary load failure'; adapter.isInitialized() deve essere false; adapter.isInitialized() deve essere true; pipeline deve essere stato chiamato 2 volte & \ok \\
TU-340 & WordEmbedding > ritorna un Float32Array con i dati del modello & result deve essere un'istanza di Float32Array; result deve essere fakeVector & \ok \\
TU-341 & WordEmbedding > chiama l'embedder con pooling mean e normalize true & mockEmbedder deve essere stato chiamato con 'testo', { pooling: 'mean', normalize: true, } & \ok \\
TU-342 & WordEmbedding > inizializza automaticamente il modello alla prima chiamata & adapter.isInitialized() deve essere false; adapter.isInitialized() deve essere true & \ok \\
TU-343 & WordEmbedding > non richiama pipeline su generateEmbedding successive & pipeline deve essere stato chiamato 1 volte; mockEmbedder deve essere stato chiamato 3 volte & \ok \\
TU-344 & WordEmbedding > produce vettori di 384 dimensioni (dimensione MiniLM) & result.length deve essere 384 & \ok \\
TU-345 & WordEmbedding > testi diversi producono vettori diversi & result1[0] deve essere vicino a 0.1, 5; result2[0] deve essere vicino a 0.9, 5; result1 non deve essere result2 & \ok \\
TU-346 & WordEmbedding > propaga errori dell'embedder durante generateEmbedding & adapter.generateEmbedding('testo') deve lanciare 'inference failed' & \ok \\
TU-347 & DocumentClass use-cases > TU-S-browsing-80: execute() should GetDocumentClassByIdUC delega a repo.getById & repo.getById deve essere stato chiamato con 5; result deve essere entity & \ok \\
TU-348 & DocumentClass use-cases > TU-S-browsing-81: execute() should GetDocumentClassByDipIdUC delega a repo.getByDipId & repo.getByDipId deve essere stato chiamato con 9; result deve essere list & \ok \\
TU-349 & DocumentClass use-cases > TU-S-browsing-82: execute() should GetDocumentClassByStatusUC delega a repo.getByStatus & repo.getByStatus deve essere stato chiamato con IntegrityStatusEnum.INVALID; result deve essere list & \ok \\
TU-350 & DocumentClass use-cases > TU-S-browsing-83: execute() should CheckDocumentClassIntegrityStatusUC aggiorna lo status aggregato dai process & integrityService.checkDocumentClassIntegrityStatus deve essere stato chiamato con 1; result deve essere IntegrityStatusEnum.UNKNOWN & \ok \\
TU-351 & DocumentClass use-cases > TU-S-browsing-84: execute() should CheckDocumentClassIntegrityStatusUC lancia errore se classe inesistente & uc.execute(14) deve lanciare "DocumentClass with id 14 not found" & \ok \\
TU-352 & Dip use-cases > TU-S-browsing-85: execute() should return the correct dip & repo.getById deve essere stato chiamato con 10; result deve essere dip & \ok \\
TU-353 & Dip use-cases > TU-S-browsing-86: execute() should return null & repo.getById deve essere stato chiamato con 999; result deve essere null & \ok \\
TU-354 & Dip use-cases > TU-S-browsing-87: execute() should return an array of dips & repo.getByStatus deve essere stato chiamato con IntegrityStatusEnum.VALID; result deve essere list & \ok \\
TU-355 & Dip use-cases > TU-S-browsing-88: execute() should return an empty array & repo.getByStatus deve essere stato chiamato con IntegrityStatusEnum.INVALID; result deve essere uguale a [] & \ok \\
TU-356 & Dip use-cases > TU-S-browsing-89: execute() should successfully update and return the newly aggregated status & integrityService.checkDipIntegrityStatus deve essere stato chiamato con 1; result deve essere IntegrityStatusEnum.VALID & \ok \\
TU-357 & Dip use-cases > TU-S-browsing-90: execute() should throw an error & uc.execute(3) deve lanciare "Dip with id 3 not found" & \ok \\
TU-358 & Document use-cases > TU-S-browsing-91: execute() should GetDocumentByIdUC delega a repo.getById & repo.getById deve essere stato chiamato con 10; result deve essere entity & \ok \\
TU-359 & Document use-cases > TU-S-browsing-92: execute() should GetDocumentByProcessUC delega a repo.getByProcessId & repo.getByProcessId deve essere stato chiamato con 3; result deve essere list & \ok \\
TU-360 & Document use-cases > TU-S-browsing-93: execute() should GetDocumentByStatusUC delega a repo.getByStatus & repo.getByStatus deve essere stato chiamato con IntegrityStatusEnum.VALID; result deve essere list & \ok \\
TU-361 & Document use-cases > TU-S-browsing-94: execute() should CheckDocumentIntegrityStatusUC aggiorna lo status aggregato dai file & integrityService.checkDocumentIntegrityStatus deve essere stato chiamato con 1; result deve essere IntegrityStatusEnum.INVALID & \ok \\
TU-362 & Document use-cases > TU-S-browsing-95: execute() should CheckDocumentIntegrityStatusUC lancia errore se document inesistente & uc.execute(77) deve lanciare "Document with id 77 not found" & \ok \\
TU-363 & File use-cases > TU-S-browsing-96: GetFileByIdUC delega a repo.getById & repo.getById deve essere stato chiamato con 11; result deve essere entity & \ok \\
TU-364 & File use-cases > TU-S-browsing-97: GetFileByDocumentUC delega a repo.getByDocumentId & repo.getByDocumentId deve essere stato chiamato con 8; result deve essere list & \ok \\
TU-365 & File use-cases > TU-S-browsing-98: GetFileByStatusUC delega a repo.getByStatus & repo.getByStatus deve essere stato chiamato con IntegrityStatusEnum.UNKNOWN; result deve essere list & \ok \\
TU-366 & File use-cases > TU-S-browsing-99: CheckFileIntegrityStatusUC imposta VALID se hash coincide & result deve essere IntegrityStatusEnum.VALID; integrityService.checkFileIntegrityStatus deve essere stato chiamato con 8 & \ok \\
TU-367 & File use-cases > TU-S-browsing-100: CheckFileIntegrityStatusUC propaga UNKNOWN & result deve essere IntegrityStatusEnum.UNKNOWN; integrityService.checkFileIntegrityStatus deve essere stato chiamato con 8 & \ok \\
TU-368 & File use-cases > TU-S-browsing-101: CheckFileIntegrityStatusUC lancia errore se file inesistente & uc.execute(99) deve lanciare "File with id 99 not found" & \ok \\
TU-369 & ExportFileUC > ritorna NOT\_FOUND quando il file non esiste & result.success deve essere false; result.errorCode deve essere "NOT\_FOUND"; dialogPort.showSaveDialog non deve essere stato chiamato; exportPort.exportFile non deve essere stato chiamato & \ok \\
TU-370 & ExportFileUC > ritorna ExportResult.canceled() se il dialog viene annullato & result.success deve essere false; result.canceled deve essere true; exportPort.exportFile non deve essere stato chiamato & \ok \\
TU-371 & ExportFileUC > chiama exportPort.exportFile con lo stream e il path scelto dal dialog & openReadStreamArg deve contenere "base/dip"; openReadStreamArg deve contenere "docs/f.pdf"; exportPort.exportFile deve essere stato chiamato con stream, "/dest/f.pdf"; result deve essere exported & \ok \\
TU-372 & ExportFileUC > propaga eccezioni di exportPort.exportFile & makeUC().execute(1) deve lanciare "disk full" & \ok \\
TU-373 & ExportFilesUC > ritorna canceled:true se il dialog cartella viene annullato & result.canceled deve essere true; result.results deve avere lunghezza 0; onProgress non deve essere stato chiamato & \ok \\
TU-374 & ExportFilesUC > registra errore per file non trovato senza interrompere il loop & result.canceled deve essere false; result.results[0] deve corrispondere a { fileId: 99, success: false }; onProgress deve essere stato chiamato con 1, 1 & \ok \\
TU-375 & ExportFilesUC > esporta ogni file nel folder scelto e chiama onProgress & streamArg deve essere stream; destPathArg deve contenere "/dest"; destPathArg deve contenere "f.pdf"; result.canceled deve essere false; result.results[0] deve corrispondere a { fileId: 1, success: true }; onProgress deve essere stato chiamato con 1, 1 & \ok \\
TU-376 & PrintFileUC > ritorna NOT\_FOUND quando il file non esiste & result.success deve essere false; result.errorCode deve essere "NOT\_FOUND"; printPort.printSingle non deve essere stato chiamato & \ok \\
TU-377 & PrintFileUC > chiama printSingle con il path assoluto e le opzioni corrette & absolutePathArg deve contenere "base/dip"; absolutePathArg deve contenere "docs/report.pdf"; optsArg deve essere uguale a { silent: false, printBackground: true } & \ok \\
TU-378 & PrintFileUC > ritorna il risultato di printSingle quando la stampa riesce & result deve essere ok & \ok \\
TU-379 & PrintFileUC > propaga eccezioni di printSingle & makeUC().execute(1) deve lanciare "unexpected crash" & \ok \\
TU-380 & PrintFileUC > chiama fileRepo.getById con il fileId corretto & fileRepo.getById deve essere stato chiamato con 42 & \ok \\
TU-381 & PrintFilesUC > ritorna canceled:true se il dialog di conferma viene annullato & result.canceled deve essere true; result.results deve avere lunghezza 0; onProgress non deve essere stato chiamato & \ok \\
TU-382 & PrintFilesUC > registra errore per file non trovato senza interrompere il loop & result.canceled deve essere false; result.results[0] deve corrispondere a { fileId: 99, success: false }; onProgress deve essere stato chiamato con 1, 1 & \ok \\
TU-383 & PrintFilesUC > stampa ogni file e chiama onProgress & absolutePathArg deve contenere "base/dip"; absolutePathArg deve contenere "docs/f.pdf"; optsArg deve essere uguale a { silent: false, printBackground: true }; result.canceled deve essere false; result.results[0] deve corrispondere a { fileId: 1, success: true }; onProgress deve essere stato chiamato con 1, 1 & \ok \\
TU-384 & PrintFilesUC > chiama showConfirmPrint con il numero corretto di file & dialogPort.showConfirmPrint deve essere stato chiamato con 3 & \ok \\
TU-385 & Process use-cases > TU-S-browsing-102: execute() should return the correct process & repo.getById deve essere stato chiamato con 12; result deve essere entity & \ok \\
TU-386 & Process use-cases > TU-S-browsing-103: execute() should return list of processes & repo.getByDocumentClassId deve essere stato chiamato con 3; result deve essere list & \ok \\
TU-387 & Process use-cases > TU-S-browsing-104: execute() should return list of processes & repo.getByStatus deve essere stato chiamato con IntegrityStatusEnum.VALID; result deve essere list & \ok \\
TU-388 & Process use-cases > TU-S-browsing-105: execute() should update and return aggregated status & integrityService.checkProcessIntegrityStatus deve essere stato chiamato con 1; result deve essere IntegrityStatusEnum.VALID & \ok \\
TU-389 & Process use-cases > TU-S-browsing-106: execute() should throw an error & uc.execute(55) deve lanciare "Process with id 55 not found" & \ok \\
TU-390 & SearchDocumentalClassUC > ritorna le classi documentali trovate dal repository & results deve avere lunghezza 1; results[0].getId() deve essere 1; results[0].getName() deve essere "Contratti"; results[0].getUuid() deve essere "uuid-1" & \ok \\
TU-391 & SearchDocumentalClassUC > passa il nome al repository senza modificarlo & repo.searchDocumentalClasses deve essere stato chiamato con "Fatture 2026" & \ok \\
TU-392 & SearchDocumentalClassUC > ritorna array vuoto se nessuna classe corrisponde al nome & results deve avere lunghezza 0 & \ok \\
TU-393 & SearchDocumentalClassUC > passa stringa vuota al repository per ricerca senza filtro & repo.searchDocumentalClasses deve essere stato chiamato con ""; results deve avere lunghezza 3 & \ok \\
TU-394 & SearchDocumentalClassUC > preserva l'ordine dei risultati restituito dal repository & results.map((r) => r.getName()) deve essere uguale a ["AAA", "BBB", "CCC"] & \ok \\
TU-395 & SearchDocumentalClassUC > mantiene gli attributi dominio della classe documentale & results[0].getId() deve essere 1; results[0].getName() deve essere "Originale"; results[0].getIntegrityStatus() deve essere "UNKNOWN" & \ok \\
TU-396 & SearchDocumentalClassUC > propaga eccezioni del repository senza alterarle & () => uc.execute("Contratti") deve lanciare "document class search failed" & \ok \\
TU-397 & SearchProcessUC > ritorna i processi trovati dal repository & results deve avere lunghezza 1; results[0].getId() deve essere 1; results[0].getUuid() deve essere "proc-uuid-abc"; results[0].getMetadata().findNodeByName("name")?.getStringValue() deve essere "Proc A"; results[0].getMetadata().findNodeByName("type")?.getStringValue() deve essere "PROCESSO" & \ok \\
TU-398 & SearchProcessUC > passa l'uuid al repository senza modificarlo & repo.searchProcesses deve essere stato chiamato con "uuid-specifico-123" & \ok \\
TU-399 & SearchProcessUC > ritorna array vuoto se nessun processo corrisponde all'uuid & results deve avere lunghezza 0 & \ok \\
TU-400 & SearchProcessUC > passa stringa vuota al repository per ricerca senza filtro & repo.searchProcesses deve essere stato chiamato con ""; results deve avere lunghezza 2 & \ok \\
TU-401 & SearchProcessUC > preserva l'ordine dei risultati restituito dal repository & results.map((r) => r.getId()) deve essere uguale a [1, 2, 3] & \ok \\
TU-402 & SearchProcessUC > usa stringa vuota per name e type se i metadati non sono presenti & results[0].getMetadata().findNodeByName("name") deve essere null; results[0].getMetadata().findNodeByName("type") deve essere null & \ok \\
TU-403 & SearchProcessUC > accetta uuid parziali per ricerche per prefisso & repo.searchProcesses deve essere stato chiamato con "proc-" & \ok \\
TU-404 & SearchProcessUC > propaga eccezioni del repository senza alterarle & () => uc.execute("proc-uuid") deve lanciare "process search failed" & \ok \\
TU-405 & SearchDocumentsUC > mappa correttamente uuid, nome e tipoDocumento nel SearchResult & results deve avere lunghezza 1; results[0].getUuid() deve essere "uuid-1"; results[0] .getMetadata() .findNodeByName("NomeDelDocumento") ?.getStringValue() deve essere "fattura.pdf"; results[0] .getMetadata() .findNodeByName("tipoDocumento") ?.getStringValue() deve essere "DOCUMENTO INFORMATICO" & \ok \\
TU-406 & SearchDocumentsUC > ritorna array vuoto se il repository non trova risultati & results deve avere lunghezza 0 & \ok \\
TU-407 & SearchDocumentsUC > usa stringa vuota per name e type se i metadati sono assenti & results[0].getMetadata().findNodeByName("nome") deve essere null; results[0].getMetadata().findNodeByName("tipoDocumento") deve essere null & \ok \\
TU-408 & SearchDocumentsUC > passa i filtri al repository senza modificarli & repo.searchDocument deve essere stato chiamato con filters & \ok \\
TU-409 & SearchDocumentsUC > mappa correttamente più documenti restituiti dal repository & results deve avere lunghezza 3; results.map((r) => r.getUuid()) deve essere uguale a [ "uuid-a", "uuid-b", "uuid-c", ] & \ok \\
TU-410 & SearchDocumentsUC > ignora metadati con chiavi diverse da nome e tipoDocumento & results[0] .getMetadata() .findNodeByName("NomeDelDocumento") ?.getStringValue() deve essere "contratto.pdf"; results[0] .getMetadata() .findNodeByName("tipoDocumento") ?.getStringValue() deve essere "DOCUMENTO AMMINISTRATIVO INFORMATICO" & \ok \\
TU-411 & SearchDocumentsUC > usa il primo metadato corrispondente quando ci sono duplicati (nome/tipoDocumento) & results[0] .getMetadata() .findNodeByName("NomeDelDocumento") ?.getStringValue() deve essere "primo-nome.pdf"; results[0] .getMetadata() .findNodeByName("tipoDocumento") ?.getStringValue() deve essere "TIPO-1" & \ok \\
TU-412 & SearchDocumentsUC > propaga errori del repository durante la ricerca normale & () => uc.execute(emptyFilters) deve lanciare "db unavailable" & \ok \\
TU-413 & SearchSemanticUC > ritorna coppie documento+score dal repository & results deve avere lunghezza 1; results[0].document.getUuid() deve essere "uuid-s1"; results[0].score deve essere 0.92 & \ok \\
TU-414 & SearchSemanticUC > lo score semantico è sempre un numero, mai null & typeof results[0].score deve essere "number"; results[0].score non deve essere null & \ok \\
TU-415 & SearchSemanticUC > ritorna array vuoto se nessun documento supera la soglia di similarità & results deve avere lunghezza 0 & \ok \\
TU-416 & SearchSemanticUC > genera embedding dalla query e passa il vettore al repository & aiAdapter.generateEmbedding deve essere stato chiamato con "ricerca semantica avanzata"; vectorRepo.searchSimilarVectors deve essere stato chiamato con embedding, 10 & \ok \\
TU-417 & SearchSemanticUC > preserva l'ordine dei risultati restituito dal repository & results[0].score deve essere 0.95; results[1].score deve essere 0.78; results[2].score deve essere 0.61 & \ok \\
TU-418 & SearchSemanticUC > gestisce correttamente score ai valori limite (0.0 e 1.0) & results[0].score deve essere 1; results[1].score deve essere 0 & \ok \\
TU-419 & SearchSemanticUC > ritorna comunque documento e score anche con metadati assenti & results[0].document.getUuid() deve essere "uuid-empty"; results[0].score deve essere 0.88 & \ok \\
TU-420 & SearchSemanticUC > preserva il documento originale anche con metadati duplicati & results[0].document.getUuid() deve essere "uuid-dup"; results[0].score deve essere 0.5 & \ok \\
TU-421 & SearchSemanticUC > propaga errori del repository durante la ricerca semantica & uc.execute("query") deve lanciare "semantic fail" & \ok \\
TU-422 & IndexDip use-case performance > TU-F-Indexing-26: execute() should measures end-to-end indexing performance on a real DiP & true deve essere true; true deve essere true; result.success deve essere true; durationsMs deve avere lunghezza runs & \ok \\
TU-423 & IndexDip > TU-F-Indexing-27: execute() should orchestrate reader and repository writes parsing real XML & result deve essere uguale a { success: true }; dips deve avere lunghezza 1; (dips[0] as any).uuid deve essere "dip-uuid-1"; docClasses deve avere lunghezza 1; (docClasses[0] as any).uuid deve essere "dc-1"; processes deve avere lunghezza 1; (processes[0] as any).uuid deve essere "proc-1"; documents deve avere lunghezza 1; (documents[0] as any).uuid deve essere "doc-1"; files deve avere lunghezza 1; (files[0] as any).filename deve essere "./primary.pdf"; embeddingService.generateDocumentEmbedding deve essere stato chiamato 1 volte; calledFile.getPath() deve contenere "primary.pdf"; vectorRepository.saveVector deve essere stato chiamato 1 volte; vectorRepository.saveVector deve essere stato chiamato con any Vector & \ok \\
TU-424 & IndexDip > TU-F-Indexing-28: execute() should throw on structurally invalid empty files and avoid persistence & useCase.execute(dipPath) deve lanciare un errore; dips deve avere lunghezza 0; docClasses deve avere lunghezza 0 & \ok \\
TU-425 & IndexDip > TU-F-Indexing-29: execute() should continue indexing when vector generation fails & result deve essere uguale a { success: true }; files deve avere lunghezza 1; embeddingService.generateDocumentEmbedding deve essere stato chiamato 1 volte; vectorRepository.saveVector non deve essere stato chiamato & \ok \\
TU-426 & appConfig cache providers > risolve CACHE\_SERVICE\_TOKEN sulla stessa istanza di IpcCacheService & cacheByToken deve essere cacheByClass & \ok \\
TU-427 & App /routes > should redirect from empty path to /browse & location.path() deve essere '/browse' & \ok \\
TU-428 & App /routes > should have a lazy loaded search route & searchRoute deve essere truthy; searchRoute?.loadChildren deve essere definito; typeof searchRoute.loadChildren deve essere 'function' & \ok \\
TU-429 & App /routes > should have a lazy loaded detail route & detailRoute deve essere truthy; detailRoute?.loadChildren deve essere definito; typeof detailRoute.loadChildren deve essere 'function' & \ok \\
TU-430 & App /routes > should load verification routes when navigating to /integrity-dashboard & integrityRoute deve essere truthy; integrityRoute?.loadChildren deve essere definito; typeof integrityRoute.loadChildren deve essere 'function' & \ok \\
TU-431 & App /routes > should redirect from /dip to /browse & location.path() deve essere '/browse' & \ok \\
TU-432 & App /routes > should redirect from an unknown path to /browse & location.path() deve essere '/browse' & \ok \\
TU-433 & App > should create the app & app deve essere truthy & \ok \\
TU-434 & App > should render router outlet & compiled.querySelector('router-outlet') non deve essere null & \ok \\
TU-435 & ElectronLoggingGateway > dovrebbe chiamare il bridge IPC corretto e completarsi con successo & mockBridge.invoke deve essere stato chiamato con 'ipc:log:write', mockLogEntry & \ok \\
TU-436 & ElectronLoggingGateway > dovrebbe propagare l'errore se la scrittura IPC fallisce & errorResult deve essere mockError; mockBridge.invoke deve essere stato chiamato 1 volte & \ok \\
TU-437 & IpcErrorHandlerService > createError() > dovrebbe creare un AppError con la corretta mappatura (es. DOC\_BLOB\_LOAD\_ERROR) & error.code deve essere ErrorCode.DOC\_BLOB\_LOAD\_ERROR; error.message deve essere 'File non trovato'; error.source deve essere 'DocFacade'; error.category deve essere ErrorCategory.IO; error.severity deve essere ErrorSeverity.ERROR; error.recoverable deve essere true & \ok \\
TU-438 & IpcErrorHandlerService > createError() > dovrebbe creare un AppError fatale e non recuperabile per IPC\_ERROR & error.severity deve essere ErrorSeverity.FATAL; error.recoverable deve essere false & \ok \\
TU-439 & IpcErrorHandlerService > handle() / toErrorCode() > dovrebbe estrarre correttamente il codice se fornito nell'oggetto raw & result.code deve essere ErrorCode.DOC\_FORMAT\_UNSUPPORTED; result.message deve essere 'Formato .xyz non supportato'; result.source deve essere 'ElectronRenderer'; result.category deve essere ErrorCategory.IO & \ok \\
TU-440 & IpcErrorHandlerService > handle() / toErrorCode() > dovrebbe fare fallback su SEARCH\_ENGINE\_ERROR se l'oggetto non ha un codice riconosciuto & result.code deve essere ErrorCode.SEARCH\_ENGINE\_ERROR; result.message deve essere 'Errore misterioso nel motore di ricerca'; result.detail deve essere rawError.stack & \ok \\
TU-441 & IpcErrorHandlerService > handle() / toErrorCode() > dovrebbe dedurre IPC\_ERROR se il messaggio contiene la stringa "ipc" & result.code deve essere ErrorCode.IPC\_ERROR; result.category deve essere ErrorCategory.IPC & \ok \\
TU-442 & IpcErrorHandlerService > handle() / toErrorCode() > dovrebbe gestire stringhe primitive in modo sicuro & result.code deve essere ErrorCode.SEARCH\_ENGINE\_ERROR; result.message deve essere 'Qualcosa è andato storto'; result.source deve essere 'IPC Gateway' & \ok \\
TU-443 & IpcErrorHandlerService > Casi limite (Edge Cases) e Coverage > dovrebbe usare la configurazione di default in createError per un codice non mappato (righe 67-77) & error.category deve essere ErrorCategory.IPC; error.severity deve essere ErrorSeverity.ERROR; error.recoverable deve essere false & \ok \\
TU-444 & IpcErrorHandlerService > Casi limite (Edge Cases) e Coverage > dovrebbe fare fallback se l'oggetto raw ha una proprietà code non presente nell'enum (riga 98) & result.code deve essere ErrorCode.SEARCH\_ENGINE\_ERROR & \ok \\
TU-445 & IpcCacheService > dovrebbe salvare e recuperare un valore se non è scaduto & result deve essere uguale a { id: 1, name: 'Documento' } & \ok \\
TU-446 & IpcCacheService > dovrebbe restituire null e invalidare la cache se il TTL è scaduto & result deve essere null & \ok \\
TU-447 & IpcCacheService > dovrebbe restituire null per una chiave mai inserita & result deve essere null & \ok \\
TU-448 & IpcCacheService > dovrebbe invalidare esplicitamente una chiave specifica & cacheService.get('manual-key') deve essere null & \ok \\
TU-449 & IpcCacheService > dovrebbe invalidare solo le chiavi che corrispondono a un prefisso (invalidatePrefix) & cacheService.get('search/123') deve essere null; cacheService.get('search/456') deve essere null; cacheService.get('doc/789') deve essere 'dati documento' & \ok \\
TU-450 & TelemetryService > trackEvent() dovrebbe inviare un log di livello INFO senza payload & mockLoggingChannel.writeLog deve essere stato chiamato con { level: LogLevel.INFO, event: TelemetryEvent.SEARCH\_EXECUTED, timestamp: MOCK\_TIME, payload: null, } & \ok \\
TU-451 & TelemetryService > trackTiming() dovrebbe inviare un log di livello INFO con la durata nel payload & mockLoggingChannel.writeLog deve essere stato chiamato con { level: LogLevel.INFO, event: TelemetryMetric.SEARCH\_LATENCY\_MS, timestamp: MOCK\_TIME, payload: { durationMs: 350 }, } & \ok \\
TU-452 & TelemetryService > trackError() dovrebbe inviare un log di livello ERROR con l'oggetto AppError & mockLoggingChannel.writeLog deve essere stato chiamato con { level: LogLevel.ERROR, event: 'APP\_ERROR', timestamp: MOCK\_TIME, payload: mockError, } & \ok \\
TU-453 & TelemetryService > dovrebbe gestire in modo sicuro gli errori del canale di logging (fire-and-forget) & () => service.trackEvent(TelemetryEvent.SEARCH\_EXECUTED) non deve lanciare un errore; console.error deve essere stato chiamato con 'Fallimento scrittura telemetria via IPC:', logError & \ok \\
TU-454 & aggregate.mapper > mapProcessDtoToAggregateDetail > should handle undefined or null dto safely and apply default fallback values & result.idAgg!.tipoAggregazione deve essere 'Fascicolo'; result.idAgg!.idAggregazione deve essere ''; result.tipologiaFascicolo deve essere 'procedimento amministrativo'; result.soggetti deve essere uguale a [{ tipoRuolo: 'Sconosciuto', denominazione: 'N/A' }]; result.dataApertura deve essere 'N/D'; result.progressivo deve essere 1; result.chiaveDescrittiva!.oggetto deve essere 'Fascicolo n. '; result.procedimentoAmministrativo!.materiaArgomentoStruttura deve essere 'Standard'; result.procedimentoAmministrativo!.procedimento deve essere 'N/A'; result.procedimentoAmministrativo!.fasi deve essere uguale a []; result.posizioneFisicaAggregazioneDocumentale deve essere 'Archivio'; result.tempoDiConservazione deve essere 10; result.indiceDocumenti deve essere uguale a []; result.customMetadata deve essere uguale a [] & \ok \\
TU-455 & aggregate.mapper > mapProcessDtoToAggregateDetail > should map standard metadata fields correctly & result.idAgg.tipoAggregazione deve essere 'Serie'; result.idAgg.idAggregazione deve essere 'AGG-001'; result.tipologiaFascicolo deve essere 'affare'; result.dataApertura deve essere '2023-01-01T00:00:00Z'; result.dataChiusura deve essere '2024-01-01T00:00:00Z'; result.progressivo deve essere 42; result.chiaveDescrittiva.oggetto deve essere 'Fascicolo di Test'; result.chiaveDescrittiva.paroleChiave deve essere 'test, mock'; result.posizioneFisicaAggregazioneDocumentale deve essere 'Scaffale A'; result.tempoDiConservazione deve essere 99; result.note deve essere 'Nota di test'; result.idAggPrimario deve essere 'PRIM-001'; result.processSummary!.uuid deve essere 'abc-uuid'; result.processSummary!.integrityStatus deve essere 'VALID'; result.processSummary!.timestamp deve essere '2023-01-01T00:00:00Z' & \ok \\
TU-456 & aggregate.mapper > mapProcessDtoToAggregateDetail > should map assignment (assegnazione) effectively & result.assegnazione.tipoAssegnazione deve essere 'Per conoscenza'; result.assegnazione.soggettoAssegnatario.denominazione deve essere 'Mario Rossi'; result.assegnazione.soggettoAssegnatario.codiceFiscale deve essere 'RSSMRA'; result.assegnazione.dataInizioAssegnazione deve essere '2023-05-01'; result.assegnazione.dataFineAssegnazione deve essere '2023-06-01' & \ok \\
TU-457 & aggregate.mapper > mapProcessDtoToAggregateDetail > Soggetti Mapping > should map Persona Fisica (PF) subject & result.soggetti deve avere lunghezza 1; result.soggetti[0] deve essere uguale a { tipoRuolo: 'Mittente', denominazione: 'Giuseppe Verdi', codiceFiscale: 'VRDGPP', indirizzoDigitale: 'giuseppe@pec.it', } & \ok \\
TU-458 & aggregate.mapper > mapProcessDtoToAggregateDetail > Soggetti Mapping > should map Persona Giuridica (PG) subject & result.soggetti deve avere lunghezza 1; result.soggetti[0] deve essere uguale a { tipoRuolo: 'Destinatario', denominazione: 'Azienda Test SRL', codiceFiscale: '12345678901', indirizzoDigitale: 'azienda@pec.it', } & \ok \\
TU-459 & aggregate.mapper > mapProcessDtoToAggregateDetail > Soggetti Mapping > should map Amministrazione subject & result.soggetti deve avere lunghezza 1; result.soggetti[0].tipoRuolo deve essere 'Produttore'; result.soggetti[0].denominazione deve essere 'IPA123' & \ok \\
TU-460 & aggregate.mapper > mapProcessDtoToAggregateDetail > Soggetti Mapping > should map RUP subject & result.soggetti deve avere lunghezza 1; result.soggetti[0].tipoRuolo deve essere 'RUP'; result.soggetti[0].denominazione deve essere 'Luca Bianchi' & \ok \\
TU-461 & aggregate.mapper > mapProcessDtoToAggregateDetail > Soggetti Mapping > should map fallback plain subject correctly & result.soggetti[0].denominazione deve essere 'Soggetto Generico'; result.soggetti[0].tipoRuolo deve essere 'Sconosciuto' & \ok \\
TU-462 & aggregate.mapper > mapProcessDtoToAggregateDetail > Document Indices Mapping > should map DocumentoAmministrativoInformatico & result.indiceDocumenti deve avere lunghezza 1; result.indiceDocumenti[0] deve essere uguale a { tipoDocumento: 'DocumentoAmministativoinformatico', identificativo: 'DOC-AMM-1', impronta: 'HASH123', } & \ok \\
TU-463 & aggregate.mapper > mapProcessDtoToAggregateDetail > Document Indices Mapping > should map DocumentoInformatico & result.indiceDocumenti deve avere lunghezza 1; result.indiceDocumenti[0] deve essere uguale a { tipoDocumento: 'Documentoinformatico', identificativo: 'DOC-INF-2', impronta: 'HASH456', } & \ok \\
TU-464 & aggregate.mapper > mapProcessDtoToAggregateDetail > Document Indices Mapping > should map unknown document types as generic Documento & result.indiceDocumenti deve avere lunghezza 1; result.indiceDocumenti[0] deve essere uguale a { tipoDocumento: 'Documento', identificativo: 'DOC-GEN-3', impronta: undefined, } & \ok \\
TU-465 & aggregate.mapper > mapProcessDtoToAggregateDetail > Fasi Mapping > should map fasi correctly & result.procedimentoAmministrativo!.fasi deve avere lunghezza 2; result.procedimentoAmministrativo!.fasi[0] deve essere uguale a { tipoFase: 'Istruttoria', dataInizio: '2023-01-01', dataFine: '2023-01-31', }; result.procedimentoAmministrativo!.fasi[1] deve essere uguale a { tipoFase: 'Decisionale', dataInizio: '2023-02-01', dataFine: '2023-02-15', } & \ok \\
TU-466 & aggregate.mapper > mapProcessDtoToAggregateDetail > Custom Metadata > should extract explicit custom data like CustomMetadata and ArchimemoData & result.customMetadata deve essere uguale a array containing [ { nome: 'CampoExtra1', valore: 'Valore 1' }, { nome: 'CampoExtra2', valore: 'Valore 2' }, ] & \ok \\
TU-467 & aggregate.mapper > mapProcessDtoToAggregateDetail > Custom Metadata > should extract generic unmapped roots missing from known blocks & result.customMetadata deve essere uguale a array containing [{ nome: 'UnknownTag123', valore: 'Value123' }] & \ok \\
TU-468 & AggregateFacade > dovrebbe restituire i dati dalla cache senza chiamare il gateway (Cache Hit) & mockCache.get deve essere stato chiamato con 'aggregate:123'; mockGateway.invoke non deve essere stato chiamato; facade.getState()().detail deve essere uguale a mockAggregateData; mockTelemetry.trackTiming deve essere stato chiamato & \ok \\
TU-469 & AggregateFacade > dovrebbe chiamare il gateway e salvare in cache se i dati non sono presenti (Cache Miss) & mockGateway.invoke deve essere stato chiamato con IpcChannels.BROWSE\_GET\_PROCESS\_BY\_ID, 123, null; mockGateway.invoke deve essere stato chiamato con IpcChannels.BROWSE\_GET\_DOCUMENTS\_BY\_PROCESS, 123, null; mockCache.set deve essere stato chiamato con 'aggregate:123', any Object, 300000; facade.getState()().detail deve essere truthy & \ok \\
TU-470 & AggregateFacade > dovrebbe gestire gli errori, mapparli tramite ErrorHandler e tracciarli con Telemetry & mockErrorHandler.handle deve essere stato chiamato con rawError; facade.getState()().error deve essere uguale a mappedError; mockTelemetry.trackError deve essere stato chiamato con mappedError & \ok \\
TU-471 & ClasseDocumentaleApiService > should be created & service deve essere truthy & \ok \\
TU-472 & ClasseDocumentaleApiService > findAll > should call electronAPI.classeDocumentale.list and return the result & mockElectronAPI.classeDocumentale.list deve essere stato chiamato; result deve essere uguale a mockClasses & \ok \\
TU-473 & ClasseDocumentaleApiService > findById > should call electronAPI.classeDocumentale.get with correct id and return the result & mockElectronAPI.classeDocumentale.get deve essere stato chiamato con 1; result deve essere uguale a mockClasse & \ok \\
TU-474 & ClasseDocumentaleApiService > create > should call electronAPI.classeDocumentale.create with correct name and return the result & mockElectronAPI.classeDocumentale.create deve essere stato chiamato con 'Nuova Classe'; result deve essere uguale a mockClasse & \ok \\
TU-475 & document.mapper > mapDocumentDtoToDetail > should handle safely empty, undefined, or null payload & result.documentId deve essere ''; result.fileName deve essere 'Documento'; result.mimeType deve essere MimeType.UNSUPPORTED; result.metadata!.identificativo deve essere 'N/A'; result.metadata!.impronta deve essere 'N/A'; result.integrityStatus deve essere undefined; result.classification!.indice deve essere 'N/A'; result.registration!.tipoRegistro deve essere 'N/A'; result.format!.tipo deve essere 'N/A'; result.verification!.firmaDigitale deve essere 'N/A'; result.attachments!.numero deve essere 0; result.attachments!.allegati deve essere uguale a []; result.subjects deve essere uguale a []; result.customMetadata deve essere uguale a [] & \ok \\
TU-476 & document.mapper > mapDocumentDtoToDetail > Identity and MimeType > should extract correct identity and fallback mimeType to extension if missing & result.documentId deve essere 'DB-ID-123'; result.aipInfo!.uuid deve essere 'UUID-456'; result.fileName deve essere 'relazione\_tecnica.pdf'; result.mimeType deve essere MimeType.PDF & \ok \\
TU-477 & document.mapper > mapDocumentDtoToDetail > Identity and MimeType > should use explicit mimetype from metadata over file extension & result.documentId deve essere 'META-ID-789'; result.fileName deve essere 'image\_no\_ext'; result.mimeType deve essere MimeType.IMAGE & \ok \\
TU-478 & document.mapper > mapDocumentDtoToDetail > Base Metadata mapping > should map core metadata properties accurately & result.metadata!.identificativo deve essere 'DOC-123'; result.metadata!.impronta deve essere 'HASH123'; result.metadata!.algoritmoImpronta deve essere 'SHA-512'; result.metadata!.oggetto deve essere 'Oggetto Test'; result.metadata!.tipoDocumentale deve essere 'Contratto'; result.metadata!.riservatezza deve essere 'true'; result.metadata!.versione deve essere '2.0'; result.metadata!.tempoDiConservazione deve essere '99'; result.metadata!.paroleChiave deve essere uguale a ['test', 'mock'] & \ok \\
TU-479 & document.mapper > mapDocumentDtoToDetail > Registration Data (DatiDiRegistrazione) > should extract flat registration data & result.registration!.tipoRegistro deve essere 'Protocollo Ufficiale'; result.registration!.data deve essere '2023-01-01'; result.registration!.numero deve essere 'PROT-001' & \ok \\
TU-480 & document.mapper > mapDocumentDtoToDetail > Registration Data (DatiDiRegistrazione) > should extract complex nested TipoRegistro & result.registration!.tipoRegistro deve essere 'Registro Interno Nested' & \ok \\
TU-481 & document.mapper > mapDocumentDtoToDetail > Registration Data (DatiDiRegistrazione) > should fallback to Repertorio\_Registro & result.registration!.tipoRegistro deve essere 'Repertorio 123' & \ok \\
TU-482 & document.mapper > mapDocumentDtoToDetail > Subjects (Soggetti) Mapping > should parse PF (Persona Fisica) subject & result.subjects deve avere lunghezza 1; result.subjects![0].tipo deve essere SubjectType.PF; result.subjects![0].ruolo deve essere 'Autore'; result.subjects![0].campiSpecifici['Nome'] deve essere 'Mario'; result.subjects![0].campiSpecifici['CodiceFiscale'] deve essere 'RSSMRA'; result.subjects![0].campiSpecifici['IndirizziDigitaliDiRiferimento'] deve essere 'mario.rossi@pec.it' & \ok \\
TU-483 & document.mapper > mapDocumentDtoToDetail > Subjects (Soggetti) Mapping > should parse PG (Persona Giuridica) subject & result.subjects![0].tipo deve essere SubjectType.PG; result.subjects![0].campiSpecifici['DenominazioneOrganizzazione'] deve essere 'Azienda SRL'; result.subjects![0].campiSpecifici['CodiceFiscale\_PartitaIva'] deve essere 'IT123456789' & \ok \\
TU-484 & document.mapper > mapDocumentDtoToDetail > Subjects (Soggetti) Mapping > should parse PAI (Pubblica Amministrazione Italiana) & result.subjects![0].tipo deve essere SubjectType.PAI; result.subjects![0].campiSpecifici['IPAAmm'] deve essere 'IPA-111' & \ok \\
TU-485 & document.mapper > mapDocumentDtoToDetail > Subjects (Soggetti) Mapping > should parse RUP mapped as Person & result.subjects![0].tipo deve essere SubjectType.PF; result.subjects![0].campiSpecifici['Nome'] deve essere 'Giulia' & \ok \\
TU-486 & document.mapper > mapDocumentDtoToDetail > Attachments (Allegati) > should extract attachment details & result.attachments!.numero deve essere 1; result.attachments!.allegati deve avere lunghezza 1; result.attachments!.allegati![0].id deve essere 'All-01'; result.attachments!.allegati![0].descrizione deve essere 'Planimetria' & \ok \\
TU-487 & document.mapper > mapDocumentDtoToDetail > Change Tracking (TracciatureModificheDocumento) > should extract change tracking details accurately & result.changeTracking!.tipo deve essere 'Aggiornamento'; result.changeTracking!.soggetto deve essere 'Anna Verdi'; result.changeTracking!.data deve essere '2023-10-10'; result.changeTracking!.idVersionePrecedente deve essere 'v1.0' & \ok \\
TU-488 & document.mapper > mapDocumentDtoToDetail > Custom Metadata > should bundle unmapped custom variables outside standard document domain & result.customMetadata deve essere uguale a array containing [ { nome: 'SistemaSorgente', valore: 'ExtApp' }, { nome: 'FieldNotPresentInXSD', valore: '1234' }, ] & \ok \\
TU-489 & DocumentFacade > dovrebbe chiamare il gateway per i metadati e usare cache a 5 min (Cache Miss) & mockGateway.invoke deve essere stato chiamato con IpcChannels.BROWSE\_GET\_DOCUMENT\_BY\_ID, 456, null; mockCache.set deve essere stato chiamato con 'document:456', any Object, 300000; facade.getState()().detail deve essere truthy & \ok \\
TU-490 & DocumentFacade > dovrebbe chiamare il gateway per il Blob e usare cache a 1 min (Cache Miss) & mockGateway.invoke deve essere stato chiamato alla posizione 1 con IpcChannels.BROWSE\_GET\_FILE\_BY\_DOCUMENT, 456, null; mockGateway.invoke deve essere stato chiamato alla posizione 2 con IpcChannels.BROWSE\_GET\_FILE\_BUFFER\_BY\_ID, 999, null; mockGateway.invoke non deve essere stato chiamato con IpcChannels.BROWSE\_GET\_FILE\_BY\_ID, 456, null; mockCache.set deve essere stato chiamato con 'blob:456', mockBlob, 60000; result deve essere uguale a mockBlob & \ok \\
TU-491 & DocumentFacade > dovrebbe gestire un errore durante il recupero del blob e lanciare un AppError & facade.getFileBlob('456') deve essere uguale a mappedError; mockTelemetry.trackError deve essere stato chiamato con mappedError & \ok \\
TU-492 & DocumentFacade > loadDocument dovrebbe usare la cache se presente invece di chiamare IPC & mockGateway.invoke non deve essere stato chiamato; facade.getState()().detail deve essere uguale a mockDocumentData & \ok \\
TU-493 & DocumentFacade > loadDocument dovrebbe gestire errore e aggiornare lo stato con AppError & facade.getState()().error deve essere uguale a mappedError; facade.getState()().loading deve essere false; mockTelemetry.trackError deve essere stato chiamato con mappedError & \ok \\
TU-494 & DocumentFacade > loadDocument dovrebbe tentare il fallback del mimeType in base al nome file & mockGateway.invoke deve essere stato chiamato 2 volte; facade.getState()().detail?.mimeType deve essere MimeType.PDF & \ok \\
TU-495 & DocumentFacade > getFileBlob dovrebbe restituire il blob dalla cache se presente & mockGateway.invoke non deve essere stato chiamato; result deve essere uguale a mockBlob & \ok \\
TU-496 & DocumentFacade > getFileBlob dovrebbe lanciare errore se id non è numerico & facade.getFileBlob('abc') deve essere uguale a mappedError; mockErrorHandler.handle deve essere stato chiamato & \ok \\
TU-497 & DocumentFacade > getFileBlob dovrebbe lanciare errore se non ci sono file associati al documento & facade.getFileBlob('456') deve essere uguale a mappedError & \ok \\
TU-498 & ExportState > dovrebbe partire con lo stato INITIAL & state.phase() deve essere ExportPhase.IDLE; state.loading() deve essere false; state.progress() deve essere 0 & \ok \\
TU-499 & ExportState > setProcessing dovrebbe attivare il caricamento e resettare il progresso & state.phase() deve essere ExportPhase.PROCESSING; state.loading() deve essere true; state.outputContext() deve essere OutputContext.SINGLE\_EXPORT; state.progress() deve essere 0 & \ok \\
TU-500 & ExportState > setProgress dovrebbe arrotondare i valori e non superare 100 & state.progress() deve essere 46; state.progress() deve essere 100 & \ok \\
TU-501 & ExportState > setSuccess dovrebbe impostare la fase SUCCESS e spegnere il loading & state.phase() deve essere ExportPhase.SUCCESS; state.result() deve essere uguale a mockResult; state.loading() deve essere false; state.progress() deve essere 100 & \ok \\
TU-502 & ExportState > setError dovrebbe impostare la fase ERROR & state.phase() deve essere ExportPhase.ERROR; state.error() deve essere uguale a mockError; state.loading() deve essere false & \ok \\
TU-503 & ExportState > reset dovrebbe riportare lo stato all'origine & state.phase() deve essere ExportPhase.IDLE; state.loading() deve essere false; state.outputContext() deve essere null & \ok \\
TU-504 & ExportIpcGateway > senza Electron bridge > exportFile ritorna ExportResult con successo=false & result.success deve essere false & \ok \\
TU-505 & ExportIpcGateway > senza Electron bridge > exportFile ritorna errorCode BRIDGE\_UNAVAILABLE & result.errorCode deve essere 'BRIDGE\_UNAVAILABLE' & \ok \\
TU-506 & ExportIpcGateway > senza Electron bridge > exportFiles ritorna canceled=false e results vuoti & result deve essere uguale a { canceled: false, results: [] } & \ok \\
TU-507 & ExportIpcGateway > senza Electron bridge > getFileDto ritorna null & result deve essere null & \ok \\
TU-508 & ExportIpcGateway > senza Electron bridge > getFilesByDocumentId ritorna array vuoto & result deve essere uguale a [] & \ok \\
TU-509 & ExportIpcGateway > senza Electron bridge > printFile ritorna success=false & result.success deve essere false & \ok \\
TU-510 & ExportIpcGateway > senza Electron bridge > printFiles ritorna canceled=false e results vuoti & result deve essere uguale a { canceled: false, results: [] } & \ok \\
TU-511 & ExportIpcGateway > senza Electron bridge > onExportProgress ritorna una funzione di unsubscribe no-op & () => unsub() non deve lanciare un errore & \ok \\
TU-512 & ExportIpcGateway > senza Electron bridge > onPrintProgress ritorna una funzione di unsubscribe no-op & () => unsub() non deve lanciare un errore & \ok \\
TU-513 & ExportIpcGateway > exportFile() > invoca il canale corretto con fileId & ipc.invoke deve essere stato chiamato con IpcChannels.FILE\_DOWNLOAD, 42 & \ok \\
TU-514 & ExportIpcGateway > exportFile() > ritorna il risultato IPC invariato & result deve essere uguale a ipcResult & \ok \\
TU-515 & ExportIpcGateway > exportFile() > propaga l'eccezione se ipc.invoke rigetta & gateway.exportFile(1) deve lanciare 'IPC crash' & \ok \\
TU-516 & ExportIpcGateway > exportFiles() > invoca il canale corretto con i fileIds & ipc.invoke deve essere stato chiamato con IpcChannels.FILE\_DOWNLOAD\_MANY, [1, 2, 3] & \ok \\
TU-517 & ExportIpcGateway > exportFiles() > ritorna la risposta IPC invariata & result deve essere uguale a ipcResult & \ok \\
TU-518 & ExportIpcGateway > exportFiles() > ritorna canceled=true se l'utente annulla & result.canceled deve essere true & \ok \\
TU-519 & ExportIpcGateway > exportFiles() > propaga l'eccezione se ipc.invoke rigetta & gateway.exportFiles([1]) deve lanciare 'IPC crash' & \ok \\
TU-520 & ExportIpcGateway > getFileDto() > invoca il canale corretto con fileId & ipc.invoke deve essere stato chiamato con 'browse:get-file-by-id', 7 & \ok \\
TU-521 & ExportIpcGateway > getFileDto() > ritorna il FileDTO dalla risposta IPC & result deve essere uguale a dto & \ok \\
TU-522 & ExportIpcGateway > getFileDto() > ritorna null se IPC risponde null & result deve essere null & \ok \\
TU-523 & ExportIpcGateway > getFilesByDocumentId() > invoca il canale corretto con documentId & ipc.invoke deve essere stato chiamato con 'browse:get-file-by-document', 10 & \ok \\
TU-524 & ExportIpcGateway > getFilesByDocumentId() > ritorna l'array di FileDTO dalla risposta IPC & result deve essere uguale a dtos; result deve avere lunghezza 2 & \ok \\
TU-525 & ExportIpcGateway > getFilesByDocumentId() > ritorna array vuoto se IPC risponde con array vuoto & result deve essere uguale a [] & \ok \\
TU-526 & ExportIpcGateway > printFile() > invoca il canale corretto con fileId & ipc.invoke deve essere stato chiamato con IpcChannels.FILE\_PRINT, 5 & \ok \\
TU-527 & ExportIpcGateway > printFile() > ritorna la risposta IPC invariata & result deve essere uguale a { success: true } & \ok \\
TU-528 & ExportIpcGateway > printFile() > ritorna success=false con error in caso di fallimento IPC & result.success deve essere false; result.error deve essere 'Stampante non trovata' & \ok \\
TU-529 & ExportIpcGateway > printFile() > propaga l'eccezione se ipc.invoke rigetta & gateway.printFile(1) deve lanciare 'IPC crash' & \ok \\
TU-530 & ExportIpcGateway > printFiles() > invoca il canale corretto con i fileIds & ipc.invoke deve essere stato chiamato con IpcChannels.FILE\_PRINT\_MANY, [10, 11] & \ok \\
TU-531 & ExportIpcGateway > printFiles() > ritorna la risposta IPC invariata & result deve essere uguale a ipcResult & \ok \\
TU-532 & ExportIpcGateway > printFiles() > ritorna canceled=true se l'utente annulla & result.canceled deve essere true & \ok \\
TU-533 & ExportIpcGateway > printFiles() > propaga l'eccezione se ipc.invoke rigetta & gateway.printFiles([1]) deve lanciare 'IPC crash' & \ok \\
TU-534 & ExportIpcGateway > onExportProgress() > registra il listener sul canale corretto & ipc.on deve essere stato chiamato con IpcChannels.FILE\_DOWNLOAD\_PROGRESS, any Function & \ok \\
TU-535 & ExportIpcGateway > onExportProgress() > ritorna la funzione di unsubscribe fornita da ipc.on & result deve essere unsub & \ok \\
TU-536 & ExportIpcGateway > onPrintProgress() > registra il listener sul canale corretto & ipc.on deve essere stato chiamato con IpcChannels.FILE\_PRINT\_PROGRESS, any Function & \ok \\
TU-537 & ExportIpcGateway > onPrintProgress() > ritorna la funzione di unsubscribe fornita da ipc.on & result deve essere unsub & \ok \\
TU-538 & ExportFacade > signal accessors > delega phase a ExportState & facade.phase deve essere state.phase & \ok \\
TU-539 & ExportFacade > signal accessors > delega outputContext a ExportState & facade.outputContext deve essere state.outputContext & \ok \\
TU-540 & ExportFacade > signal accessors > delega result a ExportState & facade.result deve essere state.result & \ok \\
TU-541 & ExportFacade > signal accessors > delega progress a ExportState & facade.progress deve essere state.progress & \ok \\
TU-542 & ExportFacade > signal accessors > delega error a ExportState & facade.error deve essere state.error & \ok \\
TU-543 & ExportFacade > signal accessors > delega loading a ExportState & facade.loading deve essere state.loading & \ok \\
TU-544 & ExportFacade > signal accessors > delega queue a ExportState & facade.queue deve essere state.queue & \ok \\
TU-545 & ExportFacade > reset() > chiama state.reset() & state.reset deve essere stato chiamato 1 volte & \ok \\
TU-546 & ExportFacade > exportFile() > chiama setProcessing con SINGLE\_EXPORT & state.setProcessing deve essere stato chiamato con OutputContext.SINGLE\_EXPORT & \ok \\
TU-547 & ExportFacade > exportFile() > chiama reset e ritorna se il risultato è canceled & state.reset deve essere stato chiamato 1 volte; state.setSuccess non deve essere stato chiamato & \ok \\
TU-548 & ExportFacade > exportFile() > chiama setSuccess con ExportResult corretto & state.setSuccess deve essere stato chiamato con object containing { outputContext: OutputContext.SINGLE\_EXPORT, successCount: 1, failedCount: 0, } & \ok \\
TU-549 & ExportFacade > exportFile() > chiama setError se ipcResult.success è false & state.setError deve essere stato chiamato 1 volte & \ok \\
TU-550 & ExportFacade > exportFile() > usa errorCode come fallback se errorMessage è assente & error.message deve essere 'ERR\_IO' & \ok \\
TU-551 & ExportFacade > exportFile() > chiama setError con recoverable=true in caso di eccezione & error.recoverable deve essere true; error.message deve essere 'IPC crash' & \ok \\
TU-552 & ExportFacade > exportFile() > gestisce err non-Error come "Errore sconosciuto" & error.message deve essere 'Errore sconosciuto' & \ok \\
TU-553 & ExportFacade > exportFiles() > chiama setProcessing con MULTI\_EXPORT & state.setProcessing deve essere stato chiamato con OutputContext.MULTI\_EXPORT & \ok \\
TU-554 & ExportFacade > exportFiles() > chiama setError se tutti i DTO sono null & state.setError deve essere stato chiamato 1 volte & \ok \\
TU-555 & ExportFacade > exportFiles() > costruisce la coda solo con i DTO validi & queue deve avere lunghezza 2; queue.map((q: any) => q.filename) deve essere uguale a ['a.pdf', 'c.pdf'] & \ok \\
TU-556 & ExportFacade > exportFiles() > chiama reset e ritorna se exportFiles è canceled & state.reset deve essere stato chiamato 1 volte; state.setSuccess non deve essere stato chiamato & \ok \\
TU-557 & ExportFacade > exportFiles() > aggiorna il progresso tramite onExportProgress callback & calls deve essere uguale a [ expect.closeTo(33.33, 1), expect.closeTo(66.67, 1), 100, ] & \ok \\
TU-558 & ExportFacade > exportFiles() > conteggia successCount e failedCount correttamente & result.successCount deve essere 2; result.failedCount deve essere 1 & \ok \\
TU-559 & ExportFacade > exportFiles() > aggiorna lo stato dell'item a "done" in caso di successo & state.updateQueueItem deve essere stato chiamato con 7, { status: 'done', error: undefined } & \ok \\
TU-560 & ExportFacade > exportFiles() > aggiorna lo stato dell'item a "error" in caso di fallimento & state.updateQueueItem deve essere stato chiamato con 5, object containing { status: 'error', error: 'IO fail' } & \ok \\
TU-561 & ExportFacade > exportFiles() > setSuccess ha outputContext MULTI\_EXPORT & result.outputContext deve essere OutputContext.MULTI\_EXPORT & \ok \\
TU-562 & ExportFacade > exportFiles() > chiama setError se exportFiles lancia eccezione & state.setError deve essere stato chiamato 1 volte & \ok \\
TU-563 & ExportFacade > printDocument() > chiama setProcessing con SINGLE\_PRINT & state.setProcessing deve essere stato chiamato con OutputContext.SINGLE\_PRINT & \ok \\
TU-564 & ExportFacade > printDocument() > chiama printFile per i formati stampabili & gateway.printFile deve essere stato chiamato con 1 & \ok \\
TU-565 & ExportFacade > printDocument() > chiama setSuccess con SINGLE\_PRINT dopo la stampa & result.outputContext deve essere OutputContext.SINGLE\_PRINT; result.successCount deve essere 1 & \ok \\
TU-566 & ExportFacade > printDocument() > chiama setUnavailable per formati non stampabili & state.setUnavailable deve essere stato chiamato 1 volte; gateway.printFile non deve essere stato chiamato & \ok \\
TU-567 & ExportFacade > printDocument() > il messaggio setUnavailable contiene il nome del file & error.message deve contenere 'report.xlsx' & \ok \\
TU-568 & ExportFacade > printDocument() > setUnavailable ha recoverable=false & error.recoverable deve essere false & \ok \\
TU-569 & ExportFacade > printDocument() > chiama setError se getFileDto ritorna null & state.setError deve essere stato chiamato 1 volte & \ok \\
TU-570 & ExportFacade > printDocument() > chiama setError se printFile lancia eccezione & state.setError deve essere stato chiamato 1 volte & \ok \\
TU-571 & ExportFacade > printDocument() > `accetta il formato .\${ext}` & state.setSuccess deve essere stato chiamato; state.setUnavailable non deve essere stato chiamato & \ok \\
TU-572 & ExportFacade > printDocument() > è case-insensitive per l'estensione & state.setSuccess deve essere stato chiamato & \ok \\
TU-573 & ExportFacade > printDocuments() > chiama printFiles per i file stampabili validi & gateway.printFiles deve essere stato chiamato 1 volte; gateway.printFiles deve essere stato chiamato con [1, 2] & \ok \\
TU-574 & ExportFacade > printDocuments() > salta i file con DTO null o formato non stampabile & gateway.printFiles deve essere stato chiamato con [3] & \ok \\
TU-575 & ExportFacade > printDocuments() > chiama setUnavailable se non ci sono formati stampabili tra tutti quelli scelti & state.setUnavailable deve essere stato chiamato 1 volte; gateway.printFiles non deve essere stato chiamato & \ok \\
TU-576 & ExportFacade > printDocuments() > chiama setSuccess al termine & state.setSuccess deve essere stato chiamato 1 volte & \ok \\
TU-577 & ExportFacade > printDocuments() > chiama reset e ritorna se printFiles è canceled & state.reset deve essere stato chiamato 1 volte; state.setSuccess non deve essere stato chiamato & \ok \\
TU-578 & ExportFacade > printDocuments() > chiama setError se printFiles lancia eccezione & state.setError deve essere stato chiamato 1 volte & \ok \\
TU-579 & ExportProgressComponent > dovrebbe mostrare la percentuale corretta nel testo & pctElement.textContent deve contenere '45\%' & \ok \\
TU-580 & ExportProgressComponent > dovrebbe mostrare la label "Generazione PDF…" per il contesto REPORT\_PDF & labelElement.textContent deve essere 'Generazione PDF…' & \ok \\
TU-581 & ExportProgressComponent > dovrebbe mostrare la label "Stampa in corso…" per il contesto MULTI\_PRINT & labelElement.textContent deve essere 'Stampa in corso…' & \ok \\
TU-582 & ExportProgressComponent > dovrebbe aggiornare la larghezza della barra fill & fillElement.style.width deve essere '75\%' & \ok \\
TU-583 & ExportResultComponent > stato iniziale (IDLE) > non renderizza nulla con phase IDLE & fixture.nativeElement.querySelector('.result-box') deve essere null & \ok \\
TU-584 & ExportResultComponent > stato iniziale (IDLE) > il phase di default è IDLE & component.phase() deve essere ExportPhase.IDLE & \ok \\
TU-585 & ExportResultComponent > stato iniziale (IDLE) > result e error sono null di default & component.result() deve essere null; component.error() deve essere null & \ok \\
TU-586 & ExportResultComponent > phase SUCCESS > mostra il box success & fixture.nativeElement.querySelector('.result-success') non deve essere null & \ok \\
TU-587 & ExportResultComponent > phase SUCCESS > non mostra box error o warning & fixture.nativeElement.querySelector('.result-error') deve essere null; fixture.nativeElement.querySelector('.result-warning') deve essere null & \ok \\
TU-588 & ExportResultComponent > phase SUCCESS > ha role="status" e aria-live="polite" & box.getAttribute('role') deve essere 'status'; box.getAttribute('aria-live') deve essere 'polite' & \ok \\
TU-589 & ExportResultComponent > phase SUCCESS > non mostra il contatore elementi falliti se failedCount === 0 & fixture.nativeElement.querySelector('.result-sub') deve essere null & \ok \\
TU-590 & ExportResultComponent > phase SUCCESS > mostra il contatore elementi falliti se failedCount > 0 & sub non deve essere null; sub.textContent deve contenere '3 elementi non salvati' & \ok \\
TU-591 & ExportResultComponent > phase SUCCESS > non mostra nulla se result è null anche con phase SUCCESS & fixture.nativeElement.querySelector('.result-success') deve essere null & \ok \\
TU-592 & ExportResultComponent > successMessage > restituisce stringa vuota se result è null & component.successMessage() deve essere '' & \ok \\
TU-593 & ExportResultComponent > successMessage > `OutputContext.\${ctx}` & component.successMessage() deve essere msg & \ok \\
TU-594 & ExportResultComponent > successMessage > restituisce "Operazione completata" per outputContext sconosciuto & component.successMessage() deve essere 'Operazione completata' & \ok \\
TU-595 & ExportResultComponent > successMessage > il messaggio compare nel template come aria-label & box.getAttribute('aria-label') deve essere 'Documento salvato con successo' & \ok \\
TU-596 & ExportResultComponent > phase ERROR > mostra il box error & fixture.nativeElement.querySelector('.result-error') non deve essere null & \ok \\
TU-597 & ExportResultComponent > phase ERROR > non mostra box success o warning & fixture.nativeElement.querySelector('.result-success') deve essere null; fixture.nativeElement.querySelector('.result-warning') deve essere null & \ok \\
TU-598 & ExportResultComponent > phase ERROR > ha role="alert" e aria-live="assertive" & box.getAttribute('role') deve essere 'alert'; box.getAttribute('aria-live') deve essere 'assertive' & \ok \\
TU-599 & ExportResultComponent > phase ERROR > mostra il messaggio di errore & title.textContent.trim() deve essere 'Qualcosa è andato storto' & \ok \\
TU-600 & ExportResultComponent > phase ERROR > aria-label contiene "Errore:" e il messaggio & box.getAttribute('aria-label') deve essere 'Errore: Qualcosa è andato storto' & \ok \\
TU-601 & ExportResultComponent > phase ERROR > non mostra il bottone Riprova se error.recoverable è false & fixture.nativeElement.querySelector('.retry-btn') deve essere null & \ok \\
TU-602 & ExportResultComponent > phase ERROR > mostra il bottone Riprova se error.recoverable è true & btn non deve essere null; btn.getAttribute('aria-label') deve essere "Riprova l'operazione" & \ok \\
TU-603 & ExportResultComponent > phase ERROR > non mostra nulla se error è null anche con phase ERROR & fixture.nativeElement.querySelector('.result-error') deve essere null & \ok \\
TU-604 & ExportResultComponent > evento retry > emette retry al click del bottone & spy deve essere stato chiamato 1 volte & \ok \\
TU-605 & ExportResultComponent > evento retry > emette retry senza payload (void) & emitted[0] deve essere undefined & \ok \\
TU-606 & ExportResultComponent > phase UNAVAILABLE > mostra il box warning & fixture.nativeElement.querySelector('.result-warning') non deve essere null & \ok \\
TU-607 & ExportResultComponent > phase UNAVAILABLE > non mostra box success o error & fixture.nativeElement.querySelector('.result-success') deve essere null; fixture.nativeElement.querySelector('.result-error') deve essere null & \ok \\
TU-608 & ExportResultComponent > phase UNAVAILABLE > ha role="status" e aria-live="polite" & box.getAttribute('role') deve essere 'status'; box.getAttribute('aria-live') deve essere 'polite' & \ok \\
TU-609 & ExportResultComponent > phase UNAVAILABLE > mostra il messaggio di errore & title.textContent.trim() deve essere 'Funzione non disponibile' & \ok \\
TU-610 & ExportResultComponent > phase UNAVAILABLE > aria-label contiene "Operazione non disponibile:" e il messaggio & box.getAttribute('aria-label') deve essere 'Operazione non disponibile: Funzione non disponibile' & \ok \\
TU-611 & ExportResultComponent > phase UNAVAILABLE > non mostra il bottone Riprova nel warning & fixture.nativeElement.querySelector('.retry-btn') deve essere null & \ok \\
TU-612 & ExportResultComponent > phase UNAVAILABLE > non mostra nulla se error è null anche con phase UNAVAILABLE & fixture.nativeElement.querySelector('.result-warning') deve essere null & \ok \\
TU-613 & ExportResultComponent > accessibilità > lo span result-sub in SUCCESS ha role="alert" e aria-live="assertive" & sub.getAttribute('role') deve essere 'alert'; sub.getAttribute('aria-live') deve essere 'assertive' & \ok \\
TU-614 & ExportResultComponent > accessibilità > aria-atomic è true su tutti i box & box.getAttribute('aria-atomic') deve essere 'true' & \ok \\
TU-615 & ExportPageComponent > ngOnInit > carica i fileIds dal gateway usando il documentId in input & mockGateway.getFilesByDocumentId deve essere stato chiamato con 42 & \ok \\
TU-616 & ExportPageComponent > ngOnInit > non chiama il gateway se documentId non è un numero valido & mockGateway.getFilesByDocumentId non deve essere stato chiamato & \ok \\
TU-617 & ExportPageComponent > ngOnInit > chiama reset della facade al cambio di documentId & mockFacade.reset deve essere stato chiamato & \ok \\
TU-618 & ExportPageComponent > overlay progress > è visibile durante loading + multi-file & fixture.debugElement.query(By.css('app-export-progress')) non deve essere null & \ok \\
TU-619 & ExportPageComponent > overlay progress > non è visibile con un solo file & fixture.debugElement.query(By.css('app-export-progress')) deve essere null & \ok \\
TU-620 & ExportPageComponent > overlay queue > mostra la coda quando ci sono elementi & fixture.debugElement.query(By.css('.queue-container')) non deve essere null & \ok \\
TU-621 & ExportPageComponent > overlay queue > nasconde la coda al click su close & fixture.debugElement.query(By.css('.queue-container')) deve essere null & \ok \\
TU-622 & ExportPageComponent > overlay result > è visibile quando la fase è SUCCESS & fixture.debugElement.query(By.css('app-export-result')) non deve essere null & \ok \\
TU-623 & ExportPageComponent > overlay result > non è visibile in fase IDLE & fixture.debugElement.query(By.css('app-export-result')) deve essere null & \ok \\
TU-624 & ExportPageComponent > overlay result > non è visibile in fase PROCESSING & fixture.debugElement.query(By.css('app-export-result')) deve essere null & \ok \\
TU-625 & ExportPageComponent > overlay result > chiama reset al click su retry & mockFacade.reset deve essere stato chiamato & \ok \\
TU-626 & ExportPageComponent > overlay result > nasconde il result al click su close & fixture.debugElement.query(By.css('app-export-result')) deve essere null & \ok \\
TU-627 & ExportPageComponent > onExport — file singolo > chiama exportFile direttamente senza mostrare il selettore & mockFacade.exportFile deve essere stato chiamato con 7; component.showDownloadSelector() deve essere false & \ok \\
TU-628 & ExportPageComponent > onExport — file multipli > mostra il selettore download con tutti i file preselezionati & component.showDownloadSelector() deve essere true; component.selectedDownloadIds().has(1) deve essere true; component.selectedDownloadIds().has(2) deve essere true & \ok \\
TU-629 & ExportPageComponent > onExport — file multipli > toggleDownloadFile deseleziona e riseleziona un file & component.selectedDownloadIds().has(1) deve essere false; component.selectedDownloadIds().has(1) deve essere true & \ok \\
TU-630 & ExportPageComponent > onExport — file multipli > confirmDownload chiama exportFiles con gli id selezionati & mockFacade.exportFiles deve essere stato chiamato con [1, 2]; component.showDownloadSelector() deve essere false & \ok \\
TU-631 & ExportPageComponent > onExport — file multipli > confirmDownload chiama exportFile (singolo) se è rimasto un solo id selezionato & mockFacade.exportFile deve essere stato chiamato con 1; mockFacade.exportFiles non deve essere stato chiamato & \ok \\
TU-632 & ExportPageComponent > onExport — file multipli > confirmDownload non fa nulla se nessun file è selezionato & mockFacade.exportFile non deve essere stato chiamato; mockFacade.exportFiles non deve essere stato chiamato & \ok \\
TU-633 & ExportPageComponent > onExport — file multipli > cancelDownload chiude il selettore e svuota lo stato & component.showDownloadSelector() deve essere false; component.downloadableFiles() deve avere lunghezza 0; component.selectedDownloadIds().size deve essere 0 & \ok \\
TU-634 & ExportPageComponent > onExport — file multipli > click sull'overlay chiama cancelDownload & component.showDownloadSelector() deve essere false & \ok \\
TU-635 & ExportPageComponent > onPrint — file singolo stampabile > chiama printDocument direttamente senza mostrare il selettore & mockFacade.printDocument deve essere stato chiamato con 3; component.showPrintSelector() deve essere false & \ok \\
TU-636 & ExportPageComponent > onPrint — nessun file stampabile > non apre il selettore e non chiama la facade & mockFacade.printDocument non deve essere stato chiamato; component.showPrintSelector() deve essere false & \ok \\
TU-637 & ExportPageComponent > onPrint — file multipli > mostra il selettore stampa con tutti i file preselezionati & component.showPrintSelector() deve essere true; component.selectedPrintIds().has(10) deve essere true; component.selectedPrintIds().has(11) deve essere true & \ok \\
TU-638 & ExportPageComponent > onPrint — file multipli > togglePrintFile deseleziona e riseleziona un file & component.selectedPrintIds().has(10) deve essere false; component.selectedPrintIds().has(10) deve essere true & \ok \\
TU-639 & ExportPageComponent > onPrint — file multipli > confirmPrint chiama printDocuments con gli id selezionati & mockFacade.printDocuments deve essere stato chiamato con [10, 11]; component.showPrintSelector() deve essere false & \ok \\
TU-640 & ExportPageComponent > onPrint — file multipli > confirmPrint non fa nulla se nessun file è selezionato & mockFacade.printDocuments non deve essere stato chiamato & \ok \\
TU-641 & ExportPageComponent > onPrint — file multipli > cancelPrint chiude il selettore e svuota lo stato & component.showPrintSelector() deve essere false; component.printableFiles() deve avere lunghezza 0; component.selectedPrintIds().size deve essere 0 & \ok \\
TU-642 & ExportPageComponent > computed selectedDownloadCount e selectedPrintCount > selectedDownloadCount riflette la dimensione del Set & component.selectedDownloadCount() deve essere 2; component.selectedDownloadCount() deve essere 1 & \ok \\
TU-643 & ExportPageComponent > computed selectedDownloadCount e selectedPrintCount > selectedPrintCount riflette la dimensione del Set & component.selectedPrintCount() deve essere 2; component.selectedPrintCount() deve essere 1 & \ok \\
TU-644 & NodeFallbackFacade > usa la cache per dip senza chiamare gateway & mockCache.get deve essere stato chiamato con 'node-fallback:DIP:1'; mockGateway.invoke non deve essere stato chiamato; facade.getState()().detail deve essere uguale a cached & \ok \\
TU-645 & NodeFallbackFacade > carica e mappa document class via ipc by-id & mockGateway.invoke deve essere stato chiamato con IpcChannels.BROWSE\_GET\_DOCUMENT\_CLASS\_BY\_ID, 22, null; mockGateway.invoke deve essere stato chiamato con IpcChannels.BROWSE\_GET\_PROCESS\_BY\_DOCUMENT\_CLASS, 22, null; facade.getState()().detail?.title deve essere 'Classe Contratti'; facade.getState()().detail?.relatedSection?.items deve essere uguale a [ { itemType: 'PROCESS', itemId: '31', label: 'Processo Contratti', description: 'UUID: PROC-31 - Stato: UNKNOWN', }, ]; facade .getState()() .detail?.fields.some((field) => field.label.toLowerCase().includes('id interno')) deve essere false; mockCache.set deve essere stato chiamato con 'node-fallback:DOCUMENT\_CLASS:22', any Object, 300000 & \ok \\
TU-646 & NodeFallbackFacade > normalizza il nome file nel fallback FILE & facade.getState()().detail?.title deve essere 'Allegato\_1.pdf' & \ok \\
TU-647 & NodeFallbackFacade > gestisce gli errori tramite error handler & mockErrorHandler.handle deve essere stato chiamato con rawError; facade.getState()().error deve essere uguale a mappedError; mockTelemetry.trackError deve essere stato chiamato con mappedError & \ok \\
TU-648 & NodeFallbackFacade > lancia errore in toNumericId se ID non è numerico & mockGateway.invoke non deve essere stato chiamato; mockErrorHandler.handle deve essere stato chiamato; facade.getState()().error deve essere uguale a mockError & \ok \\
TU-649 & NodeFallbackFacade > carica e mappa il dettaglio DIP via gateway se non in cache & mockGateway.invoke deve essere stato chiamato con IpcChannels.BROWSE\_GET\_DIP\_BY\_ID, 5, null; facade.getState()().detail?.type deve essere 'DIP'; facade.getState()().detail?.fields deve contenere un elemento uguale a object containing { label: 'UUID', value: 'DIP-UUID-123' } & \ok \\
TU-650 & NodeFallbackFacade > solleva errore se IPC restituisce null per un nodo richiesto (es. DIP) & mockErrorHandler.handle deve essere stato chiamato; facade.getState()().error deve essere uguale a mockError & \ok \\
TU-651 & NodeFallbackFacade > applica etichette fallback ai processi correlati se Oggetto manca (resolveProcessLabel) & detail?.title deve essere 'Classe Documentale'; items.length deve essere 3; items[0].label deve essere 'Proc Alpha'; items[1].label deve essere 'Agg Beta'; items[2].label deve essere 'PROC-UUID-X' & \ok \\
TU-652 & NodeFallbackFacade > gestisce processi null/undefined con list vuota in mapRelatedProcesses & detail?.relatedSection?.items deve essere uguale a []; detail?.hint deve essere 'I metadati di processo verranno mostrati qui quando disponibili.' & \ok \\
TU-653 & DocumentToolbarComponent > dovrebbe crearsi correttamente & component deve essere truthy & \ok \\
TU-654 & DocumentToolbarComponent > dovrebbe mostrare il titolo e il formato passati come input & htmlElement.querySelector('.title')?.textContent deve contenere 'Documento di Test'; htmlElement.querySelector('.badge')?.textContent deve contenere 'PDF' & \ok \\
TU-655 & DocumentActionsComponent > should create & component deve essere truthy & \ok \\
TU-656 & DocumentActionsComponent > should display correct context label based on itemType & fixture.nativeElement.textContent deve contenere 'Documento'; fixture.nativeElement.textContent deve contenere 'Processo'; fixture.nativeElement.textContent deve contenere 'Classe Documentale'; fixture.nativeElement.textContent deve contenere 'Fascicolo' & \ok \\
TU-657 & DocumentActionsComponent > should call verifyItem and emit actionCompleted on success & mockIntegrityFacade.verifyItem deve essere stato chiamato con '123', 'DOCUMENT'; component.manualStatus() deve essere 'VALID'; spyEmit deve essere stato chiamato & \ok \\
TU-658 & DocumentActionsComponent > should handle verification failure and set status to UNKNOWN & mockIntegrityFacade.verifyItem deve essere stato chiamato con '456', 'PROCESS'; component.manualStatus() deve essere 'UNKNOWN'; spyEmit deve essere stato chiamato & \ok \\
TU-659 & DocumentActionsComponent > should compute verificationStatus prioritizing manualStatus over initial & component.verificationStatus() deve essere 'INVALID'; component.verificationStatus() deve essere 'VALID' & \ok \\
TU-660 & DocumentActionsComponent > should call exportPdf on export action & mockOutputFacade.exportPdf deve essere stato chiamato con { documentId: '789', tipo: 'PROCESS' } & \ok \\
TU-661 & DocumentActionsComponent > should call print on print action & mockOutputFacade.print deve essere stato chiamato con { documentId: '789', tipo: 'PROCESS' } & \ok \\
TU-662 & DocumentActionsComponent > should reset manualStatus when itemId or itemType changes via effect & component.manualStatus() deve essere 'VALID'; component.manualStatus() deve essere null & \ok \\
TU-663 & ItemDetailPageComponent > dovrebbe inizializzare e caricare i dati per un AGGREGATE & mockAggregateFacade.loadAggregate deve essere stato chiamato con '123'; mockDocumentFacade.loadDocument non deve essere stato chiamato; mockFallbackFacade.loadNode non deve essere stato chiamato; component.pageTitle() deve essere 'Fascicolo affare'; component.isLoading() deve essere false & \ok \\
TU-664 & ItemDetailPageComponent > dovrebbe inizializzare e caricare i dati per un DOCUMENT & mockDocumentFacade.loadDocument deve essere stato chiamato con '456'; mockAggregateFacade.loadAggregate non deve essere stato chiamato; mockProcessFacade.loadProcess non deve essere stato chiamato; mockFallbackFacade.loadNode non deve essere stato chiamato; component.pageTitle() deve essere 'delibera.pdf' & \ok \\
TU-665 & ItemDetailPageComponent > dovrebbe inizializzare e caricare i dati per un PROCESS & mockProcessFacade.loadProcess deve essere stato chiamato con '31'; mockAggregateFacade.loadAggregate non deve essere stato chiamato; mockDocumentFacade.loadDocument non deve essere stato chiamato; mockFallbackFacade.loadNode non deve essere stato chiamato; component.pageTitle() deve essere 'PROC-31'; compiled.querySelector('.process-sidebar') deve essere truthy; compiled.querySelector('.process-main-content') deve essere truthy & \ok \\
TU-666 & ItemDetailPageComponent > dovrebbe caricare fallback detail per tipo DIP & mockFallbackFacade.loadNode deve essere stato chiamato con 'DIP', '1'; component.pageTitle() deve essere 'DIP'; fixture.nativeElement.querySelector('app-node-fallback-panel') deve essere truthy & \ok \\
TU-667 & ItemDetailPageComponent > naviga quando viene selezionato un item correlato dal fallback panel & routerMock.navigate deve essere stato chiamato con ['/detail', 'PROCESS', '31'] & \ok \\
TU-668 & ItemDetailPageComponent > dovrebbe mostrare lo spinner se lo stato è in loading & component.isLoading() deve essere true; compiled.querySelector('.spinner-overlay') deve essere truthy & \ok \\
TU-669 & ItemDetailPageComponent > dovrebbe mostrare un errore e permettere il retryLoad & component.currentError() deve essere uguale a fintoErrore; compiled.querySelector('app-error-dialog') deve essere truthy; mockAggregateFacade.loadAggregate deve essere stato chiamato con 'ERR-1' & \ok \\
TU-670 & ItemDetailPageComponent > dovrebbe restituire una stringa vuota per il pageTitle se non ci sono dettagli & component.pageTitle() deve essere '' & \ok \\
TU-671 & LoadingComponent > should create & component deve essere truthy & \ok \\
TU-672 & LoadingComponent > should show loading overlay by default & overlay deve essere truthy & \ok \\
TU-673 & LoadingService > should be created & service deve essere truthy & \ok \\
TU-674 & navigation-routing > mappa tutti i tipi nodo del tree verso un tipo detail & mapDipNodeTypeToDetailItemType('dip') deve essere 'DIP'; mapDipNodeTypeToDetailItemType('documentClass') deve essere 'DOCUMENT\_CLASS'; mapDipNodeTypeToDetailItemType('process') deve essere 'PROCESS'; mapDipNodeTypeToDetailItemType('document') deve essere 'DOCUMENT'; mapDipNodeTypeToDetailItemType('file') deve essere 'FILE' & \ok \\
TU-675 & navigation-routing > normalizza i tipi risultato ricerca noti & mapSearchResultTypeToDetailItemType('AGGREGAZIONE DOCUMENTALE') deve essere 'AGGREGATE'; mapSearchResultTypeToDetailItemType('documento\_informatico') deve essere 'DOCUMENT'; mapSearchResultTypeToDetailItemType('documento amministrativo informatico') deve essere 'DOCUMENT'; mapSearchResultTypeToDetailItemType('processo') deve essere 'PROCESS'; mapSearchResultTypeToDetailItemType('classe') deve essere 'DOCUMENT\_CLASS'; mapSearchResultTypeToDetailItemType('tipo sconosciuto') deve essere null & \ok \\
TU-676 & navigation-routing > riconosce i tipi rich detail & isRichDetailRouteItemType('AGGREGATE') deve essere true; isRichDetailRouteItemType('DOCUMENT') deve essere true; isRichDetailRouteItemType('PROCESS') deve essere true; isRichDetailRouteItemType('DIP') deve essere false & \ok \\
TU-677 & navigation-routing > costruisce la route detail & buildDetailRoute('DOCUMENT\_CLASS', 77) deve essere uguale a ['/detail', 'DOCUMENT\_CLASS', '77'] & \ok \\
TU-678 & DipFacade > stato iniziale > phase è idle & facade.getState()().phase deve essere 'idle' & \ok \\
TU-679 & DipFacade > stato iniziale > rootNodes è array vuoto & facade.getState()().rootNodes deve avere lunghezza 0 & \ok \\
TU-680 & DipFacade > stato iniziale > nodeCache è mappa vuota & facade.getState()().nodeCache.size deve essere 0 & \ok \\
TU-681 & DipFacade > stato iniziale > loadingNodeIds è set vuoto & facade.getState()().loadingNodeIds.size deve essere 0 & \ok \\
TU-682 & DipFacade > stato iniziale > nodeChildrenErrors è mappa vuota & facade.getState()().nodeChildrenErrors.size deve essere 0 & \ok \\
TU-683 & DipFacade > loadRootNodes — successo > imposta phase=loading durante il caricamento & phaseWhileLoading deve essere 'loading' & \ok \\
TU-684 & DipFacade > loadRootNodes — successo > imposta phase=ready dopo il successo & facade.getState()().phase deve essere 'ready' & \ok \\
TU-685 & DipFacade > loadRootNodes — successo > popola rootNodes con il nodo ritornato dal gateway & facade.getState()().rootNodes deve essere uguale a [rootNode] & \ok \\
TU-686 & DipFacade > loadRootNodes — successo > chiama getRootDip con il dipId corretto & gatewayMock.getRootDip deve essere stato chiamato con 42 & \ok \\
TU-687 & DipFacade > loadRootNodes — successo > azzera rootError al nuovo caricamento & facade.getState()().rootError deve essere definito; facade.getState()().rootError deve essere undefined & \ok \\
TU-688 & DipFacade > loadRootNodes — errore > imposta phase=idle dopo un errore & facade.getState()().phase deve essere 'idle' & \ok \\
TU-689 & DipFacade > loadRootNodes — errore > popola rootError con l'AppError ricevuto & facade.getState()().rootError deve essere uguale a error & \ok \\
TU-690 & DipFacade > loadRootNodes — errore > non modifica rootNodes in caso di errore & facade.getState()().rootNodes deve avere lunghezza 0 & \ok \\
TU-691 & DipFacade > loadChildren — nodo non trovato > non chiama getChildren se il nodo non esiste nel cache né in rootNodes & gatewayMock.getChildren non deve essere stato chiamato & \ok \\
TU-692 & DipFacade > loadChildren — nodo non trovato > non modifica lo stato se il nodo non esiste & facade.getState()().phase deve essere stateBefore.phase & \ok \\
TU-693 & DipFacade > loadChildren — successo > aggiunge il nodeId a loadingNodeIds durante il caricamento & loadingDuringCall deve essere true & \ok \\
TU-694 & DipFacade > loadChildren — successo > rimuove il nodeId da loadingNodeIds dopo il successo & facade.getState()().loadingNodeIds.has(keyFor(parent)) deve essere false & \ok \\
TU-695 & DipFacade > loadChildren — successo > aggiunge i figli alla nodeCache & facade.getState()().nodeCache.get(keyFor(parent)) deve essere uguale a children & \ok \\
TU-696 & DipFacade > loadChildren — successo > rimuove l'errore precedente sullo stesso nodo prima di ricaricare & facade.getState()().nodeChildrenErrors.has(keyFor(parent)) deve essere true; facade.getState()().nodeChildrenErrors.has(keyFor(parent)) deve essere false & \ok \\
TU-697 & DipFacade > loadChildren — successo > trova il nodo nella nodeCache (non solo in rootNodes) & facade.getState()().nodeCache.get(keyFor(child)) deve essere uguale a [grandchild] & \ok \\
TU-698 & DipFacade > loadChildren — errore > rimuove il nodeId da loadingNodeIds dopo un errore & facade.getState()().loadingNodeIds.has(keyFor(parent)) deve essere false & \ok \\
TU-699 & DipFacade > loadChildren — errore > popola nodeChildrenErrors con l'errore del nodo specifico & facade.getState()().nodeChildrenErrors.get(keyFor(parent)) deve essere uguale a error & \ok \\
TU-700 & DipFacade > loadChildren — errore > non modifica la phase globale in caso di errore parziale & facade.getState()().phase deve essere 'ready' & \ok \\
TU-701 & DipFacade > loadChildren — errore > un errore su un nodo non contamina gli altri nodi & facade.getState()().nodeChildrenErrors.has(keyFor(child2)) deve essere true; facade.getState()().nodeChildrenErrors.has(keyFor(root)) deve essere false & \ok \\
TU-702 & IpcGateway > getChildren → dip → chiama getDocumentClasses & mockBridge.invoke deve essere stato chiamato con IpcChannels.BROWSE\_GET\_DOCUMENT\_CLASS\_BY\_DIP\_ID, 1 & \ok \\
TU-703 & IpcGateway > getChildren → file → ritorna array vuoto & result deve essere uguale a [] & \ok \\
TU-704 & IpcGateway > getDocumentClasses usa il nome se presente & result deve essere uguale a [ { id: 10, name: 'Classe Contratti', type: 'documentClass', hasChildren: true, } ] & \ok \\
TU-705 & IpcGateway > getDocumentClasses usa fallback solo tipo se il nome manca & result deve essere uguale a [ { id: 1234567, name: 'Classe Documentale', type: 'documentClass', hasChildren: true, } ] & \ok \\
TU-706 & IpcGateway > getProcesses usa metadata oppure fallback tipo + id & result.map((node) => node.name) deve essere uguale a [ 'Processo Contratti', 'Processo', ] & \ok \\
TU-707 & IpcGateway > getFiles normalizza i nomi con residui di path & result[0].name deve essere 'Allegato\_1.pdf'; result[0].type deve essere 'file'; result[0].hasChildren deve essere false & \ok \\
TU-708 & IpcGateway > usa la cache se presente & result deve essere cached; mockBridge.invoke non deve essere stato chiamato & \ok \\
TU-709 & IpcGateway > salva in cache dopo invoke & mockCache.set deve essere stato chiamato & \ok \\
TU-710 & IpcGateway > gestisce errori tramite IpcErrorHandler & gateway.getChildren(parent) deve lanciare un errore; mockErrorHandler.handle deve essere stato chiamato con error & \ok \\
TU-711 & DipTreeNodeComponent > dovrebbe creare il componente & component deve essere truthy & \ok \\
TU-712 & DipTreeNodeComponent > dovrebbe mostrare il nome del nodo & compiled.textContent deve contenere mockNode.name & \ok \\
TU-713 & DipTreeNodeComponent > dovrebbe emettere toggle quando clicco il bottone & component.toggle.emit deve essere stato chiamato con mockNode; component.nodeSelected.emit deve essere stato chiamato con mockNode & \ok \\
TU-714 & DipTreeNodeComponent > dovrebbe chiamare stopPropagation su toggle & spy deve essere stato chiamato & \ok \\
TU-715 & DipTreeNodeComponent > dovrebbe emettere nodeSelected e toggle quando clicco il nodo espandibile & component.nodeSelected.emit deve essere stato chiamato con mockNode; component.toggle.emit deve essere stato chiamato con mockNode & \ok \\
TU-716 & DipTreeNodeComponent > non dovrebbe emettere toggle al click del nodo se non ha figli & component.nodeSelected.emit deve essere stato chiamato con mockNode; component.toggle.emit non deve essere stato chiamato & \ok \\
TU-717 & DipTreeNodeComponent > non dovrebbe emettere toggle al click del nodo se è in loading & component.nodeSelected.emit deve essere stato chiamato con mockNode; component.toggle.emit non deve essere stato chiamato & \ok \\
TU-718 & DipTreeNodeComponent > dovrebbe mostrare il bottone toggle se ha figli & button deve essere truthy & \ok \\
TU-719 & DipTreeNodeComponent > non dovrebbe mostrare il bottone toggle se non ha figli & button deve essere falsy & \ok \\
TU-720 & DipTreeNodeComponent > dovrebbe mostrare spinner quando isLoading è true & spinner deve essere truthy & \ok \\
TU-721 & DipTreeNodeComponent > non dovrebbe mostrare spinner quando isLoading è false & spinner deve essere falsy & \ok \\
TU-722 & DipTreeNodeComponent > non dovrebbe mostrare inline error se non c’è errore & error deve essere falsy & \ok \\
TU-723 & DipTreeNodeComponent > dovrebbe mostrare dot vuoto quando non è espanso & dot deve essere truthy; dot.classList.contains('node-dot-closed') deve essere true & \ok \\
TU-724 & DipTreeNodeComponent > dovrebbe mostrare dot pieno quando è espanso & dot deve essere truthy; dot.classList.contains('node-dot-open') deve essere true & \ok \\
TU-725 & InlineErrorComponent > dovrebbe creare il componente & component deve essere truthy & \ok \\
TU-726 & InlineErrorComponent > mostra il messaggio di errore & el.textContent deve contenere 'Errore test' & \ok \\
TU-727 & InlineErrorComponent > mostra il pulsante retry se recoverable=true & button deve essere truthy; button.nativeElement.textContent deve contenere 'Riprova' & \ok \\
TU-728 & InlineErrorComponent > non mostra il pulsante retry se recoverable=false & button deve essere null & \ok \\
TU-729 & InlineErrorComponent > emette retry quando si clicca il pulsante & spy deve essere stato chiamato & \ok \\
TU-730 & AppShellComponent > dovrebbe creare il componente & component deve essere truthy & \ok \\
TU-731 & AppShellComponent > dovrebbe mostrare loading quando phase === loading & fixture.nativeElement.textContent deve contenere 'Caricamento' & \ok \\
TU-732 & AppShellComponent > dovrebbe navigare al dettaglio DOCUMENT & routerMock.navigate deve essere stato chiamato con ['/detail', 'DOCUMENT', '1'] & \ok \\
TU-733 & AppShellComponent > dovrebbe navigare al dettaglio PROCESS per nodo process & routerMock.navigate deve essere stato chiamato con ['/detail', 'PROCESS', '42'] & \ok \\
TU-734 & AppShellComponent > dovrebbe navigare al dettaglio DIP per nodo dip & routerMock.navigate deve essere stato chiamato con ['/detail', 'DIP', '7'] & \ok \\
TU-735 & AppShellComponent > dovrebbe navigare al dettaglio DOCUMENT\_CLASS per nodo documentClass & routerMock.navigate deve essere stato chiamato con ['/detail', 'DOCUMENT\_CLASS', '9'] & \ok \\
TU-736 & AppShellComponent > dovrebbe navigare al dettaglio FILE per nodo file & routerMock.navigate deve essere stato chiamato con ['/detail', 'FILE', '11'] & \ok \\
TU-737 & DipExplorerComponent > creazione > dovrebbe creare il componente & component deve essere truthy & \ok \\
TU-738 & DipExplorerComponent > rootNodes > ritorna i rootNodes dallo stato del facade & component.rootNodes deve essere uguale a nodes & \ok \\
TU-739 & DipExplorerComponent > rootNodes > ritorna array vuoto se il facade non ha nodi & component.rootNodes deve avere lunghezza 0 & \ok \\
TU-740 & DipExplorerComponent > rootNodes > si aggiorna quando il facade cambia i rootNodes & component.rootNodes deve avere lunghezza 0; component.rootNodes deve essere uguale a nodes & \ok \\
TU-741 & DipExplorerComponent > state > espone il Signal del facade & component.state() deve corrispondere a { phase: 'ready' } & \ok \\
TU-742 & DipExplorerComponent > state > riflette il phase loading quando il facade è in caricamento & component.state().phase deve essere 'loading' & \ok \\
TU-743 & DipExplorerComponent > state > si aggiorna reattivamente al cambio di phase & component.state().phase deve essere 'idle'; component.state().phase deve essere 'ready' & \ok \\
TU-744 & DipExplorerComponent > node selection navigation > naviga verso il dettaglio fallback per nodo dip & routerMock.navigate deve essere stato chiamato con ['/detail', 'DIP', '3'] & \ok \\
TU-745 & DipExplorerComponent > node selection navigation > naviga verso il dettaglio fallback per nodo documentClass & routerMock.navigate deve essere stato chiamato con ['/detail', 'DOCUMENT\_CLASS', '12'] & \ok \\
TU-746 & DipExplorerComponent > node selection navigation > naviga verso il dettaglio fallback per nodo file & routerMock.navigate deve essere stato chiamato con ['/detail', 'FILE', '44'] & \ok \\
TU-747 & DipTree > creazione > dovrebbe creare il componente & component deve essere truthy & \ok \\
TU-748 & DipTree > flatNodes — sorgente dati > usa rootNodes dal facade se @Input non è fornito & component.flatNodes()[0].node deve essere node & \ok \\
TU-749 & DipTree > flatNodes — sorgente dati > usa @Input rootNodes se fornito, ignorando il facade & ids deve contenere 2; ids non deve contenere 1 & \ok \\
TU-750 & DipTree > flatNodes — sorgente dati > ritorna array vuoto se non ci sono nodi & component.flatNodes() deve avere lunghezza 0 & \ok \\
TU-751 & DipTree > flatNodes — struttura FlatNode > assegna depth=0 ai nodi radice & component.flatNodes()[0].depth deve essere 0 & \ok \\
TU-752 & DipTree > flatNodes — struttura FlatNode > assegna isLoading=true se il nodo è in loadingNodeIds & component.flatNodes()[0].isLoading deve essere true & \ok \\
TU-753 & DipTree > flatNodes — struttura FlatNode > assegna isLoading=false se il nodo non è in loadingNodeIds & component.flatNodes()[0].isLoading deve essere false & \ok \\
TU-754 & DipTree > flatNodes — struttura FlatNode > assegna childrenError se presente nella mappa errori & component.flatNodes()[0].childrenError deve essere error & \ok \\
TU-755 & DipTree > flatNodes — struttura FlatNode > assegna childrenError=undefined se non ci sono errori & component.flatNodes()[0].childrenError deve essere undefined & \ok \\
TU-756 & DipTree > toggleNode — espansione > non include figli se il nodo non è espanso & component.flatNodes() deve avere lunghezza 1 & \ok \\
TU-757 & DipTree > toggleNode — espansione > include i figli con depth=1 dopo toggleNode & flat deve avere lunghezza 2; flat[1].node.id deve essere 2; flat[1].depth deve essere 1 & \ok \\
TU-758 & DipTree > toggleNode — espansione > collassa i figli al secondo toggle sullo stesso nodo & component.flatNodes() deve avere lunghezza 2; component.flatNodes() deve avere lunghezza 1 & \ok \\
TU-759 & DipTree > toggleNode — espansione > segna isExpanded=true dopo il toggle & component.flatNodes()[0].isExpanded deve essere true & \ok \\
TU-760 & DipTree > toggleNode — espansione > segna isExpanded=false dopo il secondo toggle & component.flatNodes()[0].isExpanded deve essere false & \ok \\
TU-761 & DipTree > toggleNode — espansione > gestisce nodi con stesso id ma tipo diverso senza collisioni & component.flatNodes() deve avere lunghezza 2; component.flatNodes()[0].isExpanded deve essere true; component.flatNodes()[1].isExpanded deve essere false; facadeMock.loadChildren deve essere stato chiamato con child & \ok \\
TU-762 & DipTree > toggleNode — loadChildren > chiama facade.loadChildren se i figli non sono in cache & facadeMock.loadChildren deve essere stato chiamato con parent & \ok \\
TU-763 & DipTree > toggleNode — loadChildren > non chiama facade.loadChildren se i figli sono già in cache & facadeMock.loadChildren non deve essere stato chiamato & \ok \\
TU-764 & DipTree > toggleNode — loadChildren > non chiama facade.loadChildren quando si collassa un nodo & facadeMock.loadChildren non deve essere stato chiamato & \ok \\
TU-765 & DipTree > reattività — Signal facade > aggiorna flatNodes quando il facade aggiunge figli alla cache & component.flatNodes() deve avere lunghezza 1; component.flatNodes() deve avere lunghezza 2 & \ok \\
TU-766 & DipTree > reattività — Signal facade > mostra isLoading=true quando facade aggiunge nodo a loadingNodeIds & component.flatNodes()[0].isLoading deve essere true & \ok \\
TU-767 & DipTree > selectNode — @Output > emette nodeSelected quando si chiama selectNode & spy deve essere stato chiamato con node & \ok \\
TU-768 & DipTree > selectNode — @Output > emette il nodo corretto tra più nodi presenti & spy deve essere stato chiamato con nodes[1]; spy non deve essere stato chiamato con nodes[0] & \ok \\
TU-769 & ProcessMetadataComponent > renderizza le sezioni principali con i test id previsti & el.querySelector('[data-testid="process-metadata-card-anagrafica"]') deve essere truthy; el.querySelector('[data-testid="process-metadata-card-overview"]') deve essere truthy; el.querySelector('[data-testid="process-metadata-card-submission"]') deve essere truthy; el.querySelector('[data-testid="process-metadata-card-conservation"]') deve essere truthy; el.querySelector('[data-testid="process-metadata-card-custom"]') deve essere null; el.querySelector('[data-testid="process-metadata-heading-overview"]')?.textContent?.trim() deve essere 'Contesto del Processo'; el.querySelector('[data-testid="process-metadata-row-process-id"]')?.textContent deve contenere '31'; el.querySelector('[data-testid="process-metadata-row-process-status"]')?.textContent deve contenere 'TERMINATO'; el.querySelector('[data-testid="process-metadata-row-document-class-timestamp"]') ?.textContent deve contenere formatReadableDate('2026-04-08T10:30:00Z'); el.querySelector('[data-testid="process-metadata-row-submission-data-inizio"]')?.textContent deve contenere formatReadableDate('2026-04-08'); el.querySelector('[data-testid="process-metadata-row-conservation-data-inizio"]') ?.textContent deve contenere formatReadableDate('2026-04-08'); el.querySelector('[data-testid="process-metadata-row-overview-oggetto"]')?.textContent deve contenere 'Processo Contratti' & \ok \\
TU-770 & ProcessMetadataComponent > mostra i campi opzionali di conservazione quando presenti & el.querySelector('[data-testid="process-metadata-row-conservation-data-fine"]') deve essere truthy; el.querySelector('[data-testid="process-metadata-row-conservation-uuid-attivatore"]') deve essere truthy; el.querySelector('[data-testid="process-metadata-row-conservation-uuid-terminatore"]') deve essere truthy; el.querySelector('[data-testid="process-metadata-row-conservation-canale-attivazione"]') deve essere truthy; el.querySelector('[data-testid="process-metadata-row-conservation-canale-terminazione"]') deve essere truthy; el.querySelector('[data-testid="process-metadata-row-conservation-stato"]') deve essere truthy & \ok \\
TU-771 & ProcessMetadataComponent > mostra i campi opzionali di versamento quando presenti & el.querySelector('[data-testid="process-metadata-row-submission-data-fine"]') deve essere truthy; el.querySelector('[data-testid="process-metadata-row-submission-uuid-attivatore"]') deve essere truthy; el.querySelector('[data-testid="process-metadata-row-submission-uuid-terminatore"]') deve essere truthy; el.querySelector('[data-testid="process-metadata-row-submission-canale-attivazione"]') deve essere truthy; el.querySelector('[data-testid="process-metadata-row-submission-canale-terminazione"]') deve essere truthy; el.querySelector('[data-testid="process-metadata-row-submission-stato"]') deve essere truthy & \ok \\
TU-772 & ProcessMetadataComponent > nasconde il contesto processo quando i campi sono tutti N/A & el.querySelector('[data-testid="process-metadata-card-overview"]') deve essere null; el.querySelector('[data-testid="process-metadata-overview-empty"]') deve essere truthy; el.querySelector('[data-testid="optional-field-absent-message"]')?.textContent deve contenere 'Nessun contesto del processo disponibile nei metadati' & \ok \\
TU-773 & ProcessMetadataComponent > nasconde i campi opzionali di conservazione quando assenti & el.querySelector('[data-testid="process-metadata-row-conservation-data-fine"]') deve essere null; el.querySelector('[data-testid="process-metadata-row-conservation-uuid-attivatore"]') deve essere null; el.querySelector('[data-testid="process-metadata-row-conservation-uuid-terminatore"]') deve essere null; el.querySelector('[data-testid="process-metadata-row-conservation-canale-attivazione"]') deve essere null; el.querySelector('[data-testid="process-metadata-row-conservation-canale-terminazione"]') deve essere null; el.querySelector('[data-testid="process-metadata-row-conservation-stato"]') deve essere null & \ok \\
TU-774 & ProcessMetadataComponent > nasconde i campi opzionali di versamento quando assenti & el.querySelector('[data-testid="process-metadata-row-submission-data-fine"]') deve essere null; el.querySelector('[data-testid="process-metadata-row-submission-uuid-attivatore"]') deve essere null; el.querySelector('[data-testid="process-metadata-row-submission-uuid-terminatore"]') deve essere null; el.querySelector('[data-testid="process-metadata-row-submission-canale-attivazione"]') deve essere null; el.querySelector('[data-testid="process-metadata-row-submission-canale-terminazione"]') deve essere null; el.querySelector('[data-testid="process-metadata-row-submission-stato"]') deve essere null & \ok \\
TU-775 & ProcessMetadataComponent > usa il fallback N/A per campi classe documentale non valorizzati & el.querySelector('[data-testid="process-metadata-row-document-class-name"]')?.textContent deve contenere 'N/A'; el.querySelector('[data-testid="process-metadata-row-document-class-uuid"]')?.textContent deve contenere 'N/A'; el.querySelector('[data-testid="process-metadata-row-document-class-timestamp"]') ?.textContent deve contenere 'N/A' & \ok \\
TU-776 & ProcessMetadataComponent > non visualizza la sezione metadati aggiuntivi anche se presenti & el.querySelector('[data-testid="process-metadata-card-custom"]') deve essere null; el.querySelector('[data-testid="process-metadata-custom-empty"]') deve essere null & \ok \\
TU-777 & process.mapper > mappa correttamente sessioni di versamento e conservazione da metadati annidati & detail.processId deve essere '31'; detail.processUuid deve essere 'PROC-31'; detail.integrityStatus deve essere 'VALID'; detail.metadata.processId deve essere '31'; detail.metadata.processUuid deve essere 'PROC-31'; detail.submission.processo deve essere 'PROC-XML-UUID'; detail.submission.sessione deve essere 'SUB-SESSION-01'; detail.submission.dataInizio deve essere '2025-12-11T16:48:48.087Z'; detail.submission.dataFine deve essere '2025-12-11T16:50:39.671Z'; detail.submission.uuidAttivatore deve essere 'USR-001'; detail.submission.uuidTerminatore deve essere 'USR-001'; detail.submission.canaleAttivazione deve essere 'WebGui'; detail.submission.canaleTerminazione deve essere 'WebGui'; detail.submission.stato deve essere 'TERMINATA'; detail.conservation.processo deve essere 'PROC-XML-UUID'; detail.conservation.sessione deve essere 'PRES-SESSION-01'; detail.conservation.dataInizio deve essere '2025-12-11T16:50:39.671Z'; detail.conservation.dataFine deve essere undefined; detail.conservation.uuidAttivatore deve essere 'USR-001'; detail.conservation.canaleAttivazione deve essere 'Other'; detail.conservation.stato deve essere undefined; detail.metadata.processStatus deve essere 'CHIUSO'; detail.documentClass.id deve essere 22; detail.metadata.documentClassName deve essere 'N/A' & \ok \\
TU-778 & process.mapper > mantiene fallback legacy per metadata piatti di conservazione & detail.submission.processo deve essere 'PRES-31'; detail.submission.sessione deve essere 'SESSION-99'; detail.submission.dataInizio deve essere '2026-04-01'; detail.submission.dataFine deve essere undefined; detail.submission.uuidAttivatore deve essere 'USR-ATT-VERS-001'; detail.submission.canaleAttivazione deve essere 'PortaleWeb'; detail.submission.stato deve essere 'TERMINATA'; detail.conservation.processo deve essere 'PRES-31'; detail.conservation.sessione deve essere 'SESSION-99'; detail.conservation.dataInizio deve essere '2026-04-08'; detail.conservation.dataFine deve essere '2026-04-09'; detail.conservation.uuidAttivatore deve essere 'USR-ATT-CONS-001'; detail.conservation.canaleAttivazione deve essere 'ConservazioneBatch'; detail.conservation.stato deve essere 'TERMINATA'; detail.metadata.processStatus deve essere 'COMPLETATO' & \ok \\
TU-779 & process.mapper > usa fallback robusti con metadata minimali & detail.processId deve essere '7'; detail.processUuid deve essere 'N/A'; detail.integrityStatus deve essere 'UNKNOWN'; detail.overview.oggetto deve essere 'N/D Test'; detail.conservation.processo deve essere 'N/A'; detail.submission.processo deve essere 'N/A'; detail.submission.dataInizio deve essere 'N/A'; detail.documentClass.id deve essere null; detail.documentClass.name deve essere 'N/A'; detail.metadata.documentClassUuid deve essere 'N/A' & \ok \\
TU-780 & ProcessFacade > usa la cache senza chiamare il gateway & mockCache.get deve essere stato chiamato con 'process:31'; mockGateway.invoke non deve essere stato chiamato; facade.getState()().detail deve essere uguale a cachedDetail & \ok \\
TU-781 & ProcessFacade > carica processo, classe e documenti e popola lo stato & mockGateway.invoke deve essere stato chiamato con IpcChannels.BROWSE\_GET\_PROCESS\_BY\_ID, 31, null; mockGateway.invoke deve essere stato chiamato con IpcChannels.BROWSE\_GET\_DOCUMENTS\_BY\_PROCESS, 31, null; mockGateway.invoke deve essere stato chiamato con IpcChannels.BROWSE\_GET\_DOCUMENT\_CLASS\_BY\_ID, 22, null; facade.getState()().detail?.processUuid deve essere 'PROC-31'; facade.getState()().detail?.documentClass.name deve essere 'Classe Contratti'; facade.getState()().detail?.metadata.documentClassName deve essere 'Classe Contratti'; facade.getState()().detail?.indiceDocumenti.length deve essere 1; mockCache.set deve essere stato chiamato con 'process:31', any Object, 300000 & \ok \\
TU-782 & ProcessFacade > espone errore applicativo quando la chiamata IPC fallisce & mockErrorHandler.handle deve essere stato chiamato con rawError; facade.getState()().error deve essere uguale a mappedError; mockTelemetry.trackError deve essere stato chiamato con mappedError & \ok \\
TU-783 & ProcessFacade > isProcess ritorna true quando trova il processo via IPC & result deve essere true; mockGateway.invoke deve essere stato chiamato con IpcChannels.BROWSE\_GET\_PROCESS\_BY\_ID, 31, null & \ok \\
TU-784 & ProcessFacade > isProcess ritorna false con id non numerico & result deve essere false; mockGateway.invoke non deve essere stato chiamato & \ok \\
TU-785 & Broad Integration: Full Search Engine Flow (Servizi Reali) > 1. Ricerca Libera: dalla UI al canale ipc:search:text & mockElectronBridge.invoke deve essere stato chiamato con 'ipc:search:text', object containing { text: 'Ricerca Valida' }, any AbortSignal & \ok \\
TU-786 & Broad Integration: Full Search Engine Flow (Servizi Reali) > 2. Ricerca Avanzata: dal Pannello Filtri al canale ipc:search:advanced & mockElectronBridge.invoke deve essere stato chiamato con 'ipc:search:advanced', object containing { filter: expect.objectContaining({ logicOperator: 'AND', }), }, any AbortSignal & \ok \\
TU-787 & Broad Integration: Full Search Engine Flow (Servizi Reali) > 3. Ricerca Semantica: attivazione toggle e chiamata a ipc:search:semantic & mockElectronBridge.invoke deve essere stato chiamato con 'ipc:search:semantic', object containing { text: 'Concetto chiave', useSemanticSearch: true }, any AbortSignal & \ok \\
TU-788 & Broad Integration: Full Search Engine Flow (Servizi Reali) > 4. Gestione Errori Reale: crash del backend tradotto dalla UI & errorBox deve essere truthy; errorBox.nativeElement.textContent deve contenere 'IPC\_TIMEOUT' & \ok \\
TU-789 & Broad Integration: Full Search Engine Flow (Servizi Reali) > 5. Validazione Reale: la ricerca si blocca se i filtri sono invalidi & mockElectronBridge.invoke non deve essere stato chiamato & \ok \\
TU-790 & SearchIpcGateway > dovrebbe restituire i risultati dalla cache se disponibili (senza chiamare IPC) & mockCache.get deve essere stato chiamato; mockBridge.invoke non deve essere stato chiamato; result deve essere uguale a mockResults & \ok \\
TU-791 & SearchIpcGateway > dovrebbe chiamare IPC e salvare in cache se il dato non è presente & mockBridge.invoke deve essere stato chiamato con 'ipc:search:text', mockQuery, abortController.signal; mockCache.set deve essere stato chiamato; result deve essere uguale a mockResults & \ok \\
TU-792 & SearchIpcGateway > dovrebbe invalidare la cache e lanciare errore tramite l'handler in caso di Abort & mockCache.invalidatePrefix deve essere stato chiamato con 'search:text'; mockErrorHandler.handle non deve essere stato chiamato & \ok \\
TU-793 & SearchIpcGateway > dovrebbe formattare un errore di rete/IPC tramite l'IpcErrorHandlerService & mockErrorHandler.handle deve essere stato chiamato con mockError; (errorResult as any).message deve essere 'IPC timeout' & \ok \\
TU-794 & SearchIpcGateway > searchAdvanced e searchSemantic > dovrebbe usare la cache e chiamare il canale corretto per searchAdvanced & mockBridge.invoke deve essere stato chiamato con 'ipc:search:advanced', mapped, abortController.signal; result deve essere uguale a mockResults; mockCache.get deve essere stato chiamato con `search:advanced:\${JSON.stringify(mapped)}`; cachedResult deve essere uguale a mockResults & \ok \\
TU-795 & SearchIpcGateway > searchAdvanced e searchSemantic > non deve invocare IPC per filtri avanzati non significativi & mockBridge.invoke non deve essere stato chiamato con 'ipc:search:advanced', expect.anything(), expect.anything(); result deve essere uguale a [] & \ok \\
TU-796 & SearchIpcGateway > searchAdvanced e searchSemantic > dovrebbe usare la cache e chiamare il canale corretto per searchSemantic & mockBridge.invoke deve essere stato chiamato con 'ipc:search:semantic', mockQuery, abortController.signal; result deve essere uguale a mockResults; mockCache.get deve essere stato chiamato con `search:semantic:\${JSON.stringify(mockQuery)}`; cachedResult deve essere uguale a mockResults & \ok \\
TU-797 & SearchIpcGateway > searchAdvanced e searchSemantic > dovrebbe caricare le chiavi custom da IPC e cachearle & mockBridge.invoke deve essere stato chiamato con 'search:get-custom-metadata-keys', 1, abortController.signal; mockCache.set deve essere stato chiamato; result deve essere uguale a ['NomeCliente', 'NumeroPratica'] & \ok \\
TU-798 & SearchIpcGateway > searchAdvanced e searchSemantic > dovrebbe restituire le chiavi custom dalla cache senza chiamare IPC & mockCache.get deve essere stato chiamato con 'search:custom-metadata-keys:1'; mockBridge.invoke non deve essere stato chiamato con 'search:get-custom-metadata-keys', expect.anything(), expect.anything(); result deve essere uguale a cachedKeys & \ok \\
TU-799 & SearchIpcGateway > Edge Cases del Segnale di Annullamento (AbortSignal) > dovrebbe fallire immediatamente se il segnale è già stato annullato prima di iniziare & mockBridge.invoke non deve essere stato chiamato; mockCache.invalidatePrefix deve essere stato chiamato con 'search:text'; errorResult deve essere undefined & \ok \\
TU-800 & SearchIpcGateway > Edge Cases del Segnale di Annullamento (AbortSignal) > NON deve emettere risultati se la promise IPC si risolve in ritardo (dopo un abort) & nextCalled deve essere false & \ok \\
TU-801 & SearchIpcGateway > Edge Cases del Segnale di Annullamento (AbortSignal) > NON deve propagare errori se la promise IPC fallisce in ritardo (dopo un abort) & errorCount deve essere 0 & \ok \\
TU-802 & search-request.mapper custom metadata > treats custom metadata array as meaningful advanced filters & hasMeaningfulAdvancedFilters(filters) deve essere true & \ok \\
TU-803 & search-request.mapper custom metadata > maps a custom metadata array into LIKE conditions & dto non deve essere null; dto?.filter.items deve contenere un elemento uguale a object containing { logicOperator: 'AND', items: expect.arrayContaining([ expect.objectContaining({ path: 'DocumentoInformatico.CustomMetadata.NomeCliente', operator: 'LIKE', value: '\%Alfa Solutions S.p.A.\%', }), ]), } & \ok \\
TU-804 & search-request.mapper custom metadata > maps a plain key to the selected document root custom metadata path & dto non deve essere null; dto?.filter.items deve contenere un elemento uguale a object containing { logicOperator: 'AND', items: expect.arrayContaining([ expect.objectContaining({ path: 'DocumentoAmministrativoInformatico.CustomMetadata.NomeCliente', operator: 'LIKE', value: '\%Alfa Solutions S.p.A.\%', }), ]), } & \ok \\
TU-805 & search-request.mapper custom metadata > maps legacy full path DocumentoInformatico.CustomMetadata.Key preserving rooted path & dto non deve essere null; dto?.filter.items deve contenere un elemento uguale a object containing { logicOperator: 'AND', items: expect.arrayContaining([ expect.objectContaining({ path: 'DocumentoInformatico.CustomMetadata.NomeCliente', operator: 'LIKE', value: '\%Alfa Solutions S.p.A.\%', }), ]), } & \ok \\
TU-806 & search-request.mapper custom metadata > normalizes legacy aggregate root to AggregazioneDocumentaliInformatiche.CustomMetadata.Key & dto non deve essere null; dto?.filter.items deve contenere un elemento uguale a object containing { logicOperator: 'AND', items: expect.arrayContaining([ expect.objectContaining({ path: 'AggregazioneDocumentaliInformatiche.CustomMetadata.NomeCliente', operator: 'LIKE', value: '\%Alfa Solutions S.p.A.\%', }), ]), } & \ok \\
TU-807 & search-request.mapper custom metadata > maps multiple custom metadata entries with AND semantics & conditions deve contenere un elemento uguale a object containing { path: 'DocumentoInformatico.CustomMetadata.NomeCliente', operator: 'LIKE', value: '\%Alfa\%', }; conditions deve contenere un elemento uguale a object containing { path: 'DocumentoInformatico.CustomMetadata.Pagato', operator: 'LIKE', value: '\%true\%', } & \ok \\
TU-808 & search-request.mapper custom metadata > maps plain key to all supported roots when document type is not selected & dto non deve essere null; dto?.filter.items deve contenere un elemento uguale a object containing { logicOperator: 'AND', items: expect.arrayContaining([ expect.objectContaining({ logicOperator: 'OR', items: expect.arrayContaining([ expect.objectContaining({ path: 'DocumentoInformatico.CustomMetadata.NomeCliente', operator: 'LIKE', value: '\%Alfa\%', }), expect.objectContaining({ path: 'DocumentoAmministrativoInformatico.CustomMetadata.NomeCliente', operator: 'LIKE', value: '\%Alfa\%', }), expect.objectContaining({ path: 'AggregazioneDocumentaliInformatiche.CustomMetadata.NomeCliente', operator: 'LIKE', value: '\%Alfa\%', }), ]), }), ]), } & \ok \\
TU-809 & search-request.mapper > returns null for empty advanced filters & dto deve essere null; hasMeaningfulAdvancedFilters({ common: {}, diDai: {}, aggregate: {}, customMeta: null, subject: [], } as any) deve essere false & \ok \\
TU-810 & search-request.mapper > maps text filters with LIKE contains & dto non deve essere null; dto?.filter.items deve contenere un elemento uguale a object containing { logicOperator: 'AND', items: expect.arrayContaining([ expect.objectContaining({ path: 'DocumentoInformatico.Note', operator: 'LIKE', value: '\%protocollo\%', }), ]), } & \ok \\
TU-811 & search-request.mapper > normalizes custom metadata path to DocumentoRoot.CustomMetadata.Key and applies LIKE & dto?.filter.items deve contenere un elemento uguale a object containing { logicOperator: 'AND', items: expect.arrayContaining([ expect.objectContaining({ path: 'DocumentoInformatico.CustomMetadata.NomeCliente', operator: 'LIKE', value: '\%Acme\%', }), ]), } & \ok \\
TU-812 & search-request.mapper > maps aggregate filters under AggregazioneDocumentale root & dto?.filter.items deve contenere un elemento uguale a object containing { logicOperator: 'AND', items: expect.arrayContaining([ expect.objectContaining({ path: 'AggregazioneDocumentale.TipoAgg.TipoAggregazione', operator: 'EQ', value: 'Fascicolo', }), expect.objectContaining({ path: 'AggregazioneDocumentale.TipoAgg.IdAggregazione', operator: 'EQ', value: 'A-123', }), ]), } & \ok \\
TU-813 & search-request.mapper > maps subject criteria including role/type/details with ELEM\_MATCH & subjectGroup deve essere definito; subjectGroup.items[0] deve essere uguale a object containing { path: 'DocumentoInformatico.Soggetti.Ruolo', operator: 'ELEM\_MATCH', } & \ok \\
TU-814 & search-request.mapper > does not emit container-level conditions for subject type wrappers & nestedItems deve contenere un elemento uguale a object containing { path: 'Altro.TipoRuolo', operator: 'EQ', value: 'Altro', }; nestedItems non deve contenere un elemento uguale a object containing { path: 'Altro.PF', }; nestedItems deve contenere un elemento uguale a object containing { path: 'Altro.PF.Cognome', operator: 'LIKE', value: '\%Rossi\%', }; nestedItems deve contenere un elemento uguale a object containing { path: 'Altro.PF.Nome', operator: 'LIKE', value: '\%Paolo\%', } & \ok \\
TU-815 & search-request.mapper > maps role and type wrappers aligned with Soggetti XML structure & nestedItems deve contenere un elemento uguale a object containing { path: 'ResponsabileGestioneDocumentale.TipoRuolo', operator: 'EQ', value: 'Responsabile della Gestione Documentale', }; nestedItems deve contenere un elemento uguale a object containing { path: 'ResponsabileGestioneDocumentale.PF.Cognome', operator: 'LIKE', value: '\%Colombo\%', }; nestedItems deve contenere un elemento uguale a object containing { path: 'ResponsabileGestioneDocumentale.PF.Nome', operator: 'LIKE', value: '\%Marco\%', } & \ok \\
TU-816 & search-request.mapper > maps subject ELEM\_MATCH under DAI root when tipoDocumento is DAI & subjectGroup deve essere definito; subjectGroup.items deve avere lunghezza 1; subjectGroup.items[0] deve essere uguale a object containing { path: 'DocumentoAmministrativoInformatico.Soggetti.Ruolo', operator: 'ELEM\_MATCH', } & \ok \\
TU-817 & search-request.mapper > maps subject ELEM\_MATCH across all roots when tipoDocumento is not specified & subjectGroup deve essere definito; paths deve essere uguale a [ 'AggregazioneDocumentale.Soggetti.Ruolo', 'DocumentoAmministrativoInformatico.Soggetti.Ruolo', 'DocumentoInformatico.Soggetti.Ruolo', ] & \ok \\
TU-818 & search-request.mapper > maps PG Codice Fiscale / Partita IVA to CodiceFiscale\_PartitaIva & nestedItems deve contenere un elemento uguale a object containing { path: 'Destinatario.PG.CodiceFiscale\_PartitaIva', operator: 'EQ', value: '31140103768', } & \ok \\
TU-819 & search-request.mapper > maps predefined diDai values to canonical XSD literals & items deve contenere un elemento uguale a object containing { path: 'DocumentoInformatico.DatiDiRegistrazione.TipoRegistro.Repertorio\_Registro.TipoRegistro', operator: 'EQ', value: String.raw`Repertorio-Registro`, }; items deve contenere un elemento uguale a object containing { path: 'DocumentoInformatico.ModalitaDiFormazione', operator: 'EQ', value: 'creazione tramite utilizzo di strumenti software che assicurino la produzione di documenti nei formati previsti in allegato 2', } & \ok \\
TU-820 & search-request.mapper > maps predefined aggregate values to canonical XSD literals & items deve contenere un elemento uguale a object containing { path: 'AggregazioneDocumentale.TipoAgg.TipoAggregazione', operator: 'EQ', value: 'Serie Di Fascicoli', } & \ok \\
TU-821 & search-request.mapper > maps Riservatezza display labels to boolean values & items deve contenere un elemento uguale a object containing { path: 'DocumentoInformatico.Riservato', operator: 'EQ', value: true, } & \ok \\
TU-822 & SearchFacade > EC1: Errore IPC > search deve gestire gli errori, tracciarli e resettare isSearching & state.error deve essere uguale a appError; state.loading deve essere false; state.isSearching deve essere false; mockErrorHandler.handle deve essere stato chiamato con rawError; mockTelemetry.trackError deve essere stato chiamato con appError & \ok \\
TU-823 & SearchFacade > EC2: Concorrenza > search() deve ignorare le chiamate se isSearching è già true & mockSearchChannel.search non deve essere stato chiamato & \ok \\
TU-824 & SearchFacade > EC3: Validazione Filtri Avanzati > searchAdvanced() NON deve chiamare l'IPC se la pre-validazione fallisce & state.validationErrors.get('dataDa') deve essere uguale a validationError; mockSearchChannel.searchAdvanced non deve essere stato chiamato & \ok \\
TU-825 & SearchFacade > EC4: Annullamento > cancelSearch() deve resettare i flag e triggerare AbortController & state.loading deve essere false; state.isSearching deve essere false; abortSpy deve essere stato chiamato & \ok \\
TU-826 & SearchFacade > EC5: Ricerca Semantica > searchSemantic() non deve eseguire se l'indicizzazione non è READY & mockSearchChannel.searchSemantic non deve essere stato chiamato & \ok \\
TU-827 & SearchFacade > EC5: Ricerca Semantica > searchSemantic() deve eseguire se l'indicizzazione è READY & mockSearchChannel.searchSemantic deve essere stato chiamato; facade.getState()().results deve essere uguale a mockResults & \ok \\
TU-828 & SearchFacade > Successo della ricerca > search() deve popolare results e annunciare il numero di risultati & state.results deve essere uguale a mockResults; state.loading deve essere false; state.isSearching deve essere false; mockTelemetry.trackEvent deve essere stato chiamato con any String; mockLiveAnnouncer.announce deve essere stato chiamato con 'Trovati 2 risultati', 'polite' & \ok \\
TU-829 & SearchFacade > Successo della ricerca > searchAdvanced() deve popolare results dopo validazione positiva & state.results deve essere uguale a mockResults; state.validationErrors.size deve essere 0; mockSearchChannel.searchAdvanced deve essere stato chiamato con dummyFilters, any AbortSignal & \ok \\
TU-830 & SearchFacade > Successo della ricerca > searchAdvanced() non deve chiamare IPC con filtri vuoti e deve svuotare i risultati & mockSearchChannel.searchAdvanced non deve essere stato chiamato; facade.getState()().results deve essere uguale a []; facade.getState()().loading deve essere false; facade.getState()().isSearching deve essere false & \ok \\
TU-831 & SearchFacade > AbortError handling > search() deve ignorare AbortError e non tracciare il telemetry & mockErrorHandler.handle non deve essere stato chiamato; mockTelemetry.trackError non deve essere stato chiamato; facade.getState()().loading deve essere false & \ok \\
TU-832 & SearchFacade > setQuery e setFilters > setQuery() deve aggiornare lo stato della query & facade.getState()().query deve essere uguale a newQuery & \ok \\
TU-833 & SearchFacade > setQuery e setFilters > setFilters() deve aggiornare lo stato dei filtri & facade.getState()().filters deve essere uguale a newFilters & \ok \\
TU-834 & SearchFacade > retry() > retry() deve triggerare una nuova ricerca & mockSearchChannel.search deve essere stato chiamato; facade.getState()().results deve essere uguale a mockResults & \ok \\
TU-835 & SearchFacade > Debounce behavior > search() deve debounce le chiamate IPC di 300ms & mockSearchChannel.search non deve essere stato chiamato; mockSearchChannel.search non deve essere stato chiamato; mockSearchChannel.search deve essere stato chiamato 1 volte & \ok \\
TU-836 & SearchFacade > Validazione errors cleanup > searchAdvanced() deve pulire gli errori di validazione precedenti su validazione positiva & facade.getState()().validationErrors.size deve essere 0 & \ok \\
TU-837 & SearchFacade > Concurrent search prevention > executeFullTextSearch() non deve eseguire se isSearching è già true & mockSearchChannel.search deve essere stato chiamato 1 volte; facade.getState()().isSearching deve essere true; mockSearchChannel.search deve essere stato chiamato 1 volte & \ok \\
TU-838 & SearchFacade > Guardie di concorrenza (Mutex) e Coverage > dovrebbe bloccare executeAdvancedSearch se un'altra ricerca è in corso & facade.getState()().isSearching deve essere true; mockSearchChannel.searchAdvanced deve essere stato chiamato 1 volte & \ok \\
TU-839 & SearchFacade > Guardie di concorrenza (Mutex) e Coverage > dovrebbe bloccare executeSemanticSearch se un'altra ricerca è in corso & facade.getState()().isSearching deve essere true; mockSearchChannel.searchSemantic deve essere stato chiamato 1 volte & \ok \\
TU-840 & SearchFacade > Guardie di concorrenza (Mutex) e Coverage > dovrebbe bloccare executeFullTextSearch se il mutex si attiva durante il debounce & mockSearchChannel.search non deve essere stato chiamato & \ok \\
TU-841 & SearchFacade > Copertura Rami (Branch Coverage) > cancelSearch() non deve fallire o lanciare eccezioni se abortController è null & () => facade.cancelSearch() non deve lanciare un errore; facade.getState()().isSearching deve essere false & \ok \\
TU-842 & SearchFacade > Copertura Rami (Branch Coverage) > prepareForSearch() deve abortire un vecchio controller se ancora presente & abortSpy deve essere stato chiamato & \ok \\
TU-843 & SearchFacade > Copertura Rami (Branch Coverage) > executeAdvancedSearch() deve gestire correttamente gli errori nel blocco catch & mockErrorHandler.handle deve essere stato chiamato con rawError; facade.getState()().error deve essere uguale a { code: 'ERR\_ADV', message: 'Errore' } & \ok \\
TU-844 & SearchFacade > Copertura Rami (Branch Coverage) > executeSemanticSearch() deve gestire correttamente gli errori nel blocco catch & mockErrorHandler.handle deve essere stato chiamato con rawError; facade.getState()().error deve essere uguale a { code: 'ERR\_SEM', message: 'Errore' } & \ok \\
TU-845 & SearchFacade > Copertura Rami (Branch Coverage) > handleError() deve processare correttamente un errore nullo o senza proprietà name & mockErrorHandler.handle deve essere stato chiamato con null & \ok \\
TU-846 & SearchResultFactoryService > dovrebbe restituire ProcessResultCardComponent per il tipo PROCESSO & componentClass deve essere ProcessResultCardComponent & \ok \\
TU-847 & SearchResultFactoryService > dovrebbe restituire ClassResultCardComponent per il tipo CLASSE & componentClass deve essere ClassResultCardComponent & \ok \\
TU-848 & SearchResultFactoryService > dovrebbe restituire DocumentResultCardComponent per ElementType.DOCUMENTO\_INFORMATICO & componentClass deve essere DocumentResultCardComponent & \ok \\
TU-849 & SearchResultFactoryService > dovrebbe restituire DocumentResultCardComponent per ElementType.DOCUMENTO\_AMMINISTRATIVO\_INFORMATICO & componentClass deve essere DocumentResultCardComponent & \ok \\
TU-850 & SearchResultFactoryService > dovrebbe restituire AggregateResultCardComponent per ElementType.AGGREGAZIONE\_DOCUMENTALE & componentClass deve essere AggregateResultCardComponent & \ok \\
TU-851 & SearchResultFactoryService > dovrebbe restituire ProcessResultCardComponent per ElementType.PROCESS & componentClass deve essere ProcessResultCardComponent & \ok \\
TU-852 & SearchResultFactoryService > CORNER CASE: dovrebbe usare DocumentResultCardComponent come fallback per tipi sconosciuti & componentClass deve essere DocumentResultCardComponent & \ok \\
TU-853 & SearchResultFactoryService > CORNER CASE: dovrebbe usare DocumentResultCardComponent se il tipo è vuoto & componentClass deve essere DocumentResultCardComponent & \ok \\
TU-854 & SemanticIndexFacade > dovrebbe inizializzarsi con lo stato IDLE & currentState.status deve essere IndexingStatus.IDLE; currentState.progressPercentage deve essere 0; mockIpcGateway.getIndexingStatus deve essere stato chiamato 1 volte & \ok \\
TU-855 & SemanticIndexFacade > dovrebbe aggiornare lo stato quando il gateway emette un nuovo valore & currentState.status deve essere IndexingStatus.INDEXING; currentState.progressPercentage deve essere 45 & \ok \\
TU-856 & SemanticIndexFacade > dovrebbe passare allo stato TIMEOUT se non riceve aggiornamenti per 120 secondi & facade.getStatus()().status deve essere IndexingStatus.TIMEOUT & \ok \\
TU-857 & SemanticIndexFacade > dovrebbe resettare il timer se riceve un aggiornamento prima della scadenza & facade.getStatus()().status deve essere IndexingStatus.INDEXING; facade.getStatus()().status deve essere IndexingStatus.TIMEOUT & \ok \\
TU-858 & SemanticIndexFacade > dovrebbe fermare il timer se riceve uno stato terminale (READY) & facade.getStatus()().status deve essere IndexingStatus.READY & \ok \\
TU-859 & SemanticIndexFacade > dovrebbe fermare il timer e impostare lo stato su ERROR se l'observable va in errore & facade.getStatus()().status deve essere IndexingStatus.ERROR; facade.getStatus()().status deve essere IndexingStatus.ERROR & \ok \\
TU-860 & AggContextStrategy > should return available roles mapped from the roleMap keys & roles.length deve essere maggiore di 0; titolareRole deve essere definito; titolareRole?.label deve essere SubjectRoleType.AMMINISTRAZIONE\_TITOLARE & \ok \\
TU-861 & AggContextStrategy > should return the correct allowed types for a valid role & types deve essere uguale a [SubjectType.PG, SubjectType.PAI, SubjectType.PAE] & \ok \\
TU-862 & AggContextStrategy > should return an empty array for an undefined or unknown role & types deve essere uguale a [] & \ok \\
TU-863 & AllContextStrategy > should return a specific list of available roles with custom labels & roles.length deve essere 16; autoreRole?.label deve essere 'Autore'; regRole?.label deve essere 'Sogg. Registrazione (DI)' & \ok \\
TU-864 & AllContextStrategy > should return the correct allowed types for a valid role & types deve essere uguale a [SubjectType.PAI, SubjectType.PAE] & \ok \\
TU-865 & AllContextStrategy > should return an empty array for an undefined or unknown role & types deve essere uguale a [] & \ok \\
TU-866 & DaiContextStrategy > should return available roles mapped from the roleMap keys & roles.length deve essere maggiore di 0; regRole deve essere definito; regRole?.label deve essere SubjectRoleType.AMMINISTRAZIONE\_REGISTRAZIONE & \ok \\
TU-867 & DaiContextStrategy > should return the correct allowed types for a valid role & types deve essere uguale a [SubjectType.PAI] & \ok \\
TU-868 & DaiContextStrategy > should return an empty array for an undefined or unknown role & types deve essere uguale a [] & \ok \\
TU-869 & DiContextStrategy > should return available roles mapped from the roleMap keys & roles.length deve essere maggiore di 0; autoreRole deve essere definito; autoreRole?.label deve essere SubjectRoleType.AUTORE & \ok \\
TU-870 & DiContextStrategy > should return the correct allowed types for a valid role & types deve essere uguale a [ SubjectType.PF, SubjectType.PG, SubjectType.PAI, SubjectType.PAE ] & \ok \\
TU-871 & DiContextStrategy > should return an empty array for an undefined or unknown role & types deve essere uguale a [] & \ok \\
TU-872 & AggregateFiltersComponent > dovrebbe istanziarsi con la struttura gerarchica corretta e array vuoto & component deve essere truthy; component.form.contains('procedimento') deve essere true; component.form.contains('assegnazione') deve essere true; component.fasiFormArray.length deve essere 0 & \ok \\
TU-873 & AggregateFiltersComponent > dovrebbe emettere filtersChanged quando il form cambia & emitSpy deve essere stato chiamato; emitted?.idAggregazione deve essere 'AGG-001' & \ok \\
TU-874 & AggregateFiltersComponent > Gestione Fasi (FormArray) > dovrebbe aggiungere e rimuovere fasi correttamente & component.fasiFormArray.length deve essere 1; component.fasiFormArray.at(0).get('tipoFase') deve essere definito; component.fasiFormArray.length deve essere 0 & \ok \\
TU-875 & AggregateFiltersComponent > Gestione Fasi (FormArray) > dovrebbe innescare add/remove tramite i bottoni HTML per la copertura & addBtn deve essere truthy; component.fasiFormArray.length deve essere 1; removeBtn deve essere truthy; component.fasiFormArray.length deve essere 0 & \ok \\
TU-876 & AggregateFiltersComponent > ngOnChanges (Integrazione Dati) > dovrebbe aggiornare i campi annidati e ricostruire le fasi & component.form.get('tipoAggregazione')?.value deve essere AggregationType.FASCICOLO; component.form.get('procedimento.materia')?.value deve essere 'Ambiente'; component.fasiFormArray.length deve essere 1 & \ok \\
TU-877 & AggregateFiltersComponent > ngOnChanges (Integrazione Dati) > dovrebbe ignorare cambiamenti se filters non è presente o nullo & patchSpy non deve essere stato chiamato & \ok \\
TU-878 & AggregateFiltersComponent > ngOnDestroy() deve pulire le sottoscrizioni & nextSpy deve essere stato chiamato; completeSpy deve essere stato chiamato & \ok \\
TU-879 & AsyncStateWrapperComponent > dovrebbe proiettare il contenuto quando non è in caricamento e non ci sono errori & projected deve essere truthy; projected.nativeElement.textContent deve contenere 'Contenuto Reale'; fixture.debugElement.query(By.css('.async-loading-state')) deve essere null; fixture.debugElement.query(By.css('.async-error-state')) deve essere null & \ok \\
TU-880 & AsyncStateWrapperComponent > dovrebbe mostrare il loader e nascondere il contenuto quando isLoading è true & loader deve essere truthy; loader.nativeElement.textContent deve contenere 'Caricamento in corso'; fixture.debugElement.query(By.css('.projected-content')) deve essere null & \ok \\
TU-881 & AsyncStateWrapperComponent > dovrebbe mostrare l'errore, nascondere il contenuto e innescare il retry al click & errorState deve essere truthy; errorState.nativeElement.textContent deve contenere 'Il server non risponde'; hostComponent.onRetry deve essere stato chiamato & \ok \\
TU-882 & AsyncStateWrapperComponent > dovrebbe usare il messaggio di errore di fallback se error.message è vuoto & errorState.nativeElement.textContent deve contenere 'Si è verificato un errore imprevisto.' & \ok \\
TU-883 & CommonFiltersComponent > dovrebbe istanziarsi correttamente con la struttura base annidata & component deve essere truthy; component.form.contains('chiaveDescrittiva') deve essere true; component.form.contains('classificazione') deve essere true; component.form.contains('conservazione') deve essere true; component.form.contains('note') deve essere true; component.form.contains('tipoDocumento') deve essere true & \ok \\
TU-884 & CommonFiltersComponent > dovrebbe emettere filtersChanged quando un valore del form cambia dal DOM & emitSpy deve essere stato chiamato; emittedValue?.note deve essere 'Test di nota' & \ok \\
TU-885 & CommonFiltersComponent > ngOnChanges (Patching dei dati esterni) > dovrebbe aggiornare i campi del form quando riceve nuovi filtri & component.form.get('tipoDocumento')?.value deve essere DocumentType.DOCUMENTO\_INFORMATICO; component.form.get('note')?.value deve essere 'Da verificare'; component.form.get('classificazione.codice')?.value deve essere 'TIT.1' & \ok \\
TU-886 & CommonFiltersComponent > ngOnChanges (Patching dei dati esterni) > non dovrebbe fallire se ngOnChanges non riceve il parametro filters & () => component.ngOnChanges({}) non deve lanciare un errore & \ok \\
TU-887 & CommonFiltersComponent > getError (Gestione Errori di Validazione) > dovrebbe restituire l'errore corretto navigando il path annidato e stamparlo nell'HTML & error deve essere truthy; error?.message deve essere 'Oggetto obbligatorio'; errorDiv deve essere truthy; errorDiv.nativeElement.textContent deve contenere 'Oggetto obbligatorio' & \ok \\
TU-888 & CommonFiltersComponent > getError (Gestione Errori di Validazione) > dovrebbe restituire undefined se l'errore non esiste o validationResult è null & component.getError('note') deve essere undefined; component.getError('classificazione.codice') deve essere undefined & \ok \\
TU-889 & CustomMetaFiltersComponent logic > does not rebuild controls when incoming filters are equivalent & component.entries.at(0) deve essere original & \ok \\
TU-890 & CustomMetaFiltersComponent logic > rebuilds controls when incoming filters change & component.entries.at(0) non deve essere original; component.entries.at(0).value deve essere uguale a { field: 'MetaKey2', value: 'MetaValue2' } & \ok \\
TU-891 & CustomMetaFiltersComponent logic > keeps local blank draft rows when parent echoes only meaningful entries & component.entries.length deve essere 2; component.entries.length deve essere 2; component.entries.at(1) deve essere draftControl & \ok \\
TU-892 & CustomMetaFiltersComponent logic > emits filtered values while typing and null when entries are blank & emitSpy deve essere stato chiamato con [{ field: 'Cliente', value: '' }]; emitSpy deve essere stato chiamato con null & \ok \\
TU-893 & CustomMetaFiltersComponent > dovrebbe istanziare un form con entries vuoto & component.entries.length deve essere 0 & \ok \\
TU-894 & CustomMetaFiltersComponent > ngOnChanges() dovrebbe sincronizzare le entries con i nuovi filtri senza emettere & component.entries.length deve essere 1; component.entries.at(0).value deve essere uguale a { field: 'Chiave1', value: 'Valore1' }; emitSpy non deve essere stato chiamato & \ok \\
TU-895 & CustomMetaFiltersComponent > ngOnChanges() non dovrebbe ricreare i controlli se il valore in ingresso è equivalente & component.entries.at(0) deve essere firstControl & \ok \\
TU-896 & CustomMetaFiltersComponent > ngOnChanges() dovrebbe resettare le entries se i filtri sono null & component.entries.length deve essere 1; component.entries.length deve essere 0 & \ok \\
TU-897 & CustomMetaFiltersComponent > dovrebbe emettere l'array filtrato se almeno un campo è valorizzato & emitSpy deve essere stato chiamato con [{ field: 'Chiave', value: '' }] & \ok \\
TU-898 & CustomMetaFiltersComponent > dovrebbe emettere null se l'utente svuota i campi delle entries & emitSpy deve essere stato chiamato con null & \ok \\
TU-899 & CustomMetaFiltersComponent > dovrebbe restituire gli errori customMeta con getError e hasAnyError & component.getError(0, 'field')?.message deve essere 'Chiave non valida'; component.hasAnyError() deve essere true & \ok \\
TU-900 & CustomMetaFiltersComponent > hasAnyError dovrebbe restituire false senza validationResult & component.hasAnyError() deve essere false & \ok \\
TU-901 & CustomMetaFiltersComponent > getSelectableKeys dovrebbe usare availableKeys ordinate e includere il valore corrente se assente & component.getSelectableKeys(0) deve essere uguale a ['Zeta', 'A', 'B'] & \ok \\
TU-902 & CustomMetaFiltersComponent > getSelectableKeyLabel dovrebbe formattare la chiave come nel dettaglio metadati & component.getSelectableKeyLabel('Process.PreservationSession.DocumentsStats.DocumentsFilesCount') deve essere 'Documents Files Count' & \ok \\
TU-903 & DiDaiFiltersComponent > dovrebbe istanziarsi correttamente con la struttura base & component deve essere truthy; component.form.contains('registrazione') deve essere true; component.form.contains('identificativoFormato') deve essere true; component.form.contains('verifica') deve essere true; component.tracciatureFormArray.length deve essere 0 & \ok \\
TU-904 & DiDaiFiltersComponent > dovrebbe emettere filtersChanged quando un valore del form cambia & emitSpy deve essere stato chiamato; emittedValue?.tipologia deve essere 'Test Tipo' & \ok \\
TU-905 & DiDaiFiltersComponent > Gestione FormArray Tracciature > dovrebbe aggiungere una nuova tracciatura vuota & component.tracciatureFormArray.length deve essere 0; component.tracciatureFormArray.length deve essere 1; component.tracciatureFormArray.at(0).value deve essere uguale a { tipoModifica: null, dataModifica: null, oraModifica: null, idVersionePrec: null, } & \ok \\
TU-906 & DiDaiFiltersComponent > Gestione FormArray Tracciature > dovrebbe rimuovere una tracciatura esistente & component.tracciatureFormArray.length deve essere 2; component.tracciatureFormArray.length deve essere 1 & \ok \\
TU-907 & DiDaiFiltersComponent > ngOnChanges (Patching dei dati esterni) > dovrebbe aggiornare i campi statici e ricostruire l'array tracciature & component.form.get('tipologia')?.value deve essere 'Bando'; component.form.get('registrazione.tipologiaRegistro')?.value deve essere RegisterType.PROTOCOLLO; component.tracciatureFormArray.length deve essere 2; component.tracciatureFormArray.at(0).get('tipoModifica')?.value deve essere ModificationType.INTEGRAZIONE; component.tracciatureFormArray.at(1).get('tipoModifica')?.value deve essere ModificationType.ANNULLAMENTO & \ok \\
TU-908 & DiDaiFiltersComponent > ngOnChanges (Patching dei dati esterni) > non dovrebbe fallire se ngOnChanges non riceve il parametro filters & () => component.ngOnChanges({}) non deve lanciare un errore & \ok \\
TU-909 & DiDaiFiltersComponent > ngOnChanges (Patching dei dati esterni) > dovrebbe svuotare l'array se i nuovi filtri non hanno tracciatureModifiche & component.tracciatureFormArray.length deve essere 1; component.tracciatureFormArray.length deve essere 0 & \ok \\
TU-910 & DiDaiFiltersComponent > getError (Gestione Errori di Validazione) > dovrebbe restituire l'errore corretto navigando il path & error deve essere truthy; error?.message deve essere 'Numero invalido' & \ok \\
TU-911 & DiDaiFiltersComponent > getError (Gestione Errori di Validazione) > dovrebbe restituire undefined se l'errore non esiste o validationResult è null & component.getError('campoInesistente') deve essere undefined; component.getError('registrazione.tipologiaFlusso') deve essere undefined & \ok \\
TU-912 & DiDaiFiltersComponent > getError (Gestione Errori di Validazione) > dovrebbe restituire errore corretto per tracciature indicizzate & error?.code deve essere 'ERR\_CONF\_DT\_002' & \ok \\
TU-913 & DiDaiFiltersComponent > dovrebbe innescare i metodi di tracciatura cliccando fisicamente i bottoni nel template HTML & addBtn deve essere truthy; component.tracciatureFormArray.length deve essere 1; removeBtn deve essere truthy; component.tracciatureFormArray.length deve essere 0 & \ok \\
TU-914 & FieldErrorComponent > should create & component deve essere truthy & \ok \\
TU-915 & FilterValueInputComponent > dovrebbe emettere valueChanged quando l'input testo cambia & spy deve essere stato chiamato con 'ciao' & \ok \\
TU-916 & FilterValueInputComponent > dovrebbe gestire il tipo DATE e stampare l'errore se presente & input deve essere truthy; errorCmp.nativeElement.textContent deve contenere 'Data obbligatoria' & \ok \\
TU-917 & FilterValueInputComponent > dovrebbe renderizzare una select se il tipo è SELECT & select deve essere truthy; options.length deve essere 2 & \ok \\
TU-918 & FilterValueInputComponent > dovrebbe emettere il valore quando l'input DATE cambia (copertura evento ngModelChange) & spy deve essere stato chiamato & \ok \\
TU-919 & FilterValueInputComponent > dovrebbe gestire il tipo NUMBER ed emettere il valore al cambiamento (copertura righe 25-30) & input deve essere truthy; spy deve essere stato chiamato & \ok \\
TU-920 & FilterValueInputComponent > dovrebbe gestire il tipo ENUM ed emettere il valore al cambiamento (copertura righe 32-35) & select deve essere truthy; spy deve essere stato chiamato & \ok \\
TU-921 & InlineErrorComponent > dovrebbe istanziarsi correttamente e mostrare il messaggio di default & component deve essere truthy; component.message deve essere 'Si è verificato un errore.'; wrapperEl deve essere truthy; iconEl.nativeElement.className deve contenere 'bi-exclamation-triangle-fill'; msgEl.nativeElement.textContent deve contenere 'Si è verificato un errore.' & \ok \\
TU-922 & InlineErrorComponent > dovrebbe aggiornare il DOM quando riceve un nuovo messaggio in input & msgEl.nativeElement.textContent deve contenere 'Credenziali non valide.' & \ok \\
TU-923 & AggregateResultCardComponent > dovrebbe renderizzare il nome e l'UUID del fascicolo & nameEl.textContent deve contenere 'Fascicolo Personale Mario Rossi'; idEl.textContent deve contenere 'ID Fascicolo: FAS-456'; badgeEl.textContent deve contenere 'AGGREGAZIONE DOCUMENTALE' & \ok \\
TU-924 & AggregateResultCardComponent > dovrebbe mostrare lo score semantico formattato se attivo & scoreBadge.textContent deve contenere '0.72' & \ok \\
TU-925 & AggregateResultCardComponent > A11Y: dovrebbe invocare onSelectAction al click & mockAction deve essere stato chiamato con mockAggregate & \ok \\
TU-926 & ClassResultCardComponent > dovrebbe renderizzare il nome e l'UUID della classe & nameEl.textContent deve contenere 'Risorse Umane'; idEl.textContent deve contenere 'ID Classe: CLS-789'; badgeEl.textContent deve contenere 'CLASSE' & \ok \\
TU-927 & ClassResultCardComponent > A11Y: dovrebbe invocare onSelectAction premendo Invio (keydown.enter) & mockAction deve essere stato chiamato con mockClass & \ok \\
TU-928 & DocumentResultCardComponent > dovrebbe renderizzare il nome e l'UUID del documento & nameEl.textContent deve contenere 'Relazione Tecnica'; idEl.textContent deve contenere 'DOC-999' & \ok \\
TU-929 & DocumentResultCardComponent > dovrebbe mostrare lo score solo se isSemanticSearch è true & scoreBadge deve essere truthy & \ok \\
TU-930 & DocumentResultCardComponent > non dovrebbe mostrare lo score se isSemanticSearch è false & scoreBadge deve essere falsy & \ok \\
TU-931 & DocumentResultCardComponent > A11Y: dovrebbe invocare onSelectAction al click & mockAction deve essere stato chiamato con mockDocument & \ok \\
TU-932 & DocumentResultCardComponent > A11Y: dovrebbe invocare onSelectAction premendo Invio (keydown.enter) & mockAction deve essere stato chiamato con mockDocument & \ok \\
TU-933 & ProcessResultCardComponent > dovrebbe renderizzare il nome e l'UUID del processo & nameEl.textContent deve contenere 'Approvazione Bilancio'; idEl.textContent deve contenere 'ID Processo: PRC-123'; badgeEl.textContent deve contenere 'PROCESS' & \ok \\
TU-934 & ProcessResultCardComponent > A11Y: dovrebbe invocare onSelectAction al click & mockAction deve essere stato chiamato con mockProcess & \ok \\
TU-935 & ProcessResultCardComponent > A11Y: dovrebbe invocare onSelectAction premendo Invio (keydown.enter) & mockAction deve essere stato chiamato con mockProcess & \ok \\
TU-936 & SearchBarRulesManager > evaluate() > should force semantic state and set activeType to FREE when isSemantic is true & result deve essere uguale a { isSemanticForced: true, activeType: SearchQueryType.FREE, } & \ok \\
TU-937 & SearchBarRulesManager > evaluate() > should NOT force semantic state and should retain the requestedType when isSemantic is false & result deve essere uguale a { isSemanticForced: false, activeType: requestedType, } & \ok \\
TU-938 & SearchBarComponent > dovrebbe istanziare il componente con i valori di default & component deve essere truthy; component.form.value.text deve essere ''; component.form.value.type deve essere MockQueryType.FREE; component.form.value.useSemanticSearch deve essere false & \ok \\
TU-939 & SearchBarComponent > ngOnChanges() dovrebbe aggiornare il form se riceve una nuova query in input senza emettere & component.form.getRawValue().text deve essere 'Fascicolo 123'; component.form.getRawValue().type deve essere MockQueryType.PROCESS\_ID; component.form.getRawValue().useSemanticSearch deve essere true; emitSpy non deve essere stato chiamato & \ok \\
TU-940 & SearchBarComponent > ngOnChanges() dovrebbe disabilitare il form se isSearching diventa true & component.form.disabled deve essere true & \ok \\
TU-941 & SearchBarComponent > ngOnChanges() dovrebbe riabilitare il form se isSearching diventa false & component.form.enabled deve essere true & \ok \\
TU-942 & SearchBarComponent > dovrebbe mostrare il loader nel template quando isSearching è true & spinner deve essere truthy; spinner.nativeElement.textContent deve contenere 'Ricerca in corso' & \ok \\
TU-943 & SearchBarComponent > ngOnDestroy() dovrebbe ripulire le sottoscrizioni chiudendo il subject & nextSpy deve essere stato chiamato; completeSpy deve essere stato chiamato & \ok \\
TU-944 & SearchResultsComponent > dovrebbe istanziarsi correttamente con array vuoto di default & component deve essere truthy; component.results.length deve essere 0 & \ok \\
TU-945 & SearchResultsComponent > dovrebbe mostrare il componente EmptyState se non ci sono risultati & emptyState deve essere truthy; emptyState.componentInstance.message deve essere 'Nessun dato'; emptyState.nativeElement.textContent deve contenere 'Nessun dato' & \ok \\
TU-946 & SearchResultsComponent > dovrebbe NON mostrare EmptyState quando ci sono risultati & emptyState deve essere null & \ok \\
TU-947 & SearchResultsComponent > dovrebbe delegare al factory il componente corretto per ogni tipo di risultato & mockFactory.getComponentForType deve essere stato chiamato con ElementType.DOCUMENTO\_INFORMATICO; mockFactory.getComponentForType deve essere stato chiamato con ElementType.DOCUMENTO\_AMMINISTRATIVO\_INFORMATICO & \ok \\
TU-948 & SearchResultsComponent > dovrebbe renderizzare un container dinamico per ogni risultato & stubs.length deve essere 2 & \ok \\
TU-949 & SearchResultsComponent > getInputs dovrebbe passare result, isSemanticSearch e onSelectAction & inputs['result'] deve essere result; inputs['isSemanticSearch'] deve essere true; typeof inputs['onSelectAction'] deve essere 'function' & \ok \\
TU-950 & SearchResultsComponent > dovrebbe emettere resultSelected quando onSelectAction viene invocata dal componente figlio & emitSpy deve essere stato chiamato 1 volte; emitSpy deve essere stato chiamato con result & \ok \\
TU-951 & SearchResultsComponent > dovrebbe propagare isSemanticSearch=false di default agli input del componente figlio & inputs['isSemanticSearch'] deve essere false & \ok \\
TU-952 & SemanticIndexStatusComponent > dovrebbe renderizzare lo stato sconosciuto se state è null & component.statusClass deve essere 'status-unknown'; component.statusLabel deve essere 'Sconosciuto'; textEl.nativeElement.textContent deve contenere 'Sconosciuto' & \ok \\
TU-953 & SemanticIndexStatusComponent > dovrebbe renderizzare lo stato READY & component.statusClass deve essere 'status-ready'; component.statusLabel deve essere 'Pronto' & \ok \\
TU-954 & SemanticIndexStatusComponent > dovrebbe renderizzare lo stato ERROR & component.statusClass deve essere 'status-error'; component.statusLabel deve essere 'Errore Indice' & \ok \\
TU-955 & SemanticIndexStatusComponent > dovrebbe renderizzare lo stato INDEXING senza progress se non fornito & component.statusClass deve essere 'status-indexing'; component.statusLabel deve essere 'In Aggiornamento...'; progressEl deve essere null & \ok \\
TU-956 & SemanticIndexStatusComponent > dovrebbe renderizzare lo stato INDEXING mostrando la percentuale se progress è presente & progressEl deve essere truthy; progressEl.nativeElement.textContent deve contenere '45\%' & \ok \\
TU-957 & SemanticIndexStatusComponent > dovrebbe gestire uno status non mappato restituendo il valore originale come label & component.statusClass deve essere 'status-unknown'; component.statusLabel deve essere 'CUSTOM\_STATE' & \ok \\
TU-958 & SubjectDetailFormComponent > dovrebbe nascere vuoto e mostrare il messaggio di fallback & component deve essere truthy; component.fields.length deve essere 0; renderedText deve contenere 'Seleziona una tipologia' & \ok \\
TU-959 & SubjectDetailFormComponent > Ricostruzione Dinamica (Strategy Pattern) > dovrebbe ricostruire il form e stampare gli input corretti per Persona Fisica (PF) & component.fields.length deve essere 3; component.form.contains('cognomePF') deve essere true; component.form.contains('nomePF') deve essere true; inputs.length deve essere 3 & \ok \\
TU-960 & SubjectDetailFormComponent > Ricostruzione Dinamica (Strategy Pattern) > dovrebbe applicare il validatore required solo ai campi obbligatori & cognomeControl?.hasValidator deve essere truthy; cognomeControl?.invalid deve essere true & \ok \\
TU-961 & SubjectDetailFormComponent > Ricostruzione Dinamica (Strategy Pattern) > dovrebbe mostrare un unico campo PG Codice Fiscale / Partita IVA & component.form.contains('codiceFiscalePartitaIvaPG') deve essere true; component.form.contains('codiceFiscalePG') deve essere false; component.form.contains('partitaIvaPG') deve essere false; component.fields.some((f) => f.key === 'codiceFiscalePartitaIvaPG') deve essere true; component.fields.some((f) => f.label === 'Codice Fiscale / Partita IVA') deve essere true & \ok \\
TU-962 & SubjectDetailFormComponent > Patching dei Dati > dovrebbe valorizzare il form quando arrivano details esterni & component.form.value.cognomePF deve essere 'Rossi'; component.form.value.nomePF deve essere 'Mario' & \ok \\
TU-963 & SubjectDetailFormComponent > Patching dei Dati > non dovrebbe fallire se arrivano details prima che il form sia costruito & () => { component.ngOnChanges({ details: new SimpleChange(null, { cognomePF: 'Rossi' }, true), }); } non deve lanciare un errore & \ok \\
TU-964 & SubjectDetailFormComponent > Emissione Eventi > dovrebbe emettere detailsChanged quando un utente compila il form & emitSpy deve essere stato chiamato con object containing { cognomePF: 'Bianchi' } & \ok \\
TU-965 & SubjectDetailFormComponent > ngOnDestroy() deve pulire le sottoscrizioni per evitare memory leak & nextSpy deve essere stato chiamato; completeSpy deve essere stato chiamato & \ok \\
TU-966 & SubjectFiltersComponent (Wizard) > dovrebbe partire allo step 0 con stato pulito & component.currentStep() deve essere 0; component.selectedRole() deve essere null & \ok \\
TU-967 & SubjectFiltersComponent (Wizard) > dovrebbe chiamare setRole cliccando i bottoni dello step 1 & expectedRole deve essere definito; setRoleSpy deve essere stato chiamato con expectedRole & \ok \\
TU-968 & SubjectFiltersComponent (Wizard) > dovrebbe chiamare setType cliccando i bottoni dello step 2 & setTypeSpy deve essere stato chiamato con allowedTypes[idx] & \ok \\
TU-969 & SubjectFiltersComponent (Wizard) > ngOnChanges() dovrebbe chiamare resetWizard() quando resetCounter viene incrementato & resetSpy deve essere stato chiamato & \ok \\
TU-970 & SubjectFiltersComponent (Wizard) > ngOnChanges() NON dovrebbe chiamare resetWizard() al primo caricamento & resetSpy non deve essere stato chiamato & \ok \\
TU-971 & SubjectFiltersComponent (Wizard) > setRole() dovrebbe salvare il ruolo e passare allo step 2 & component.selectedRole() deve essere SubjectRoleType.AUTORE; component.currentStep() deve essere 2 & \ok \\
TU-972 & SubjectFiltersComponent (Wizard) > setType() dovrebbe salvare il tipo, resettare i dettagli e passare allo step 3 & component.selectedType() deve essere SubjectType.PAI; component.currentDetails() deve essere null; component.currentStep() deve essere 3 & \ok \\
TU-973 & SubjectFiltersComponent (Wizard) > prevStep() e nextStep() dovrebbero scorrere gli step & component.currentStep() deve essere 1; component.currentStep() deve essere 3; component.currentStep() deve essere 3 & \ok \\
TU-974 & SubjectFiltersComponent (Wizard) > resetWizard() dovrebbe ripristinare tutto allo step 0 & component.currentStep() deve essere 0; component.selectedRole() deve essere null; component.selectedType() deve essere null & \ok \\
TU-975 & SubjectFiltersComponent (Wizard) > ngOnChanges() dovrebbe popolare subjectsList se riceve un array subject & component.subjectsList() deve essere uguale a mockSubject & \ok \\
TU-976 & SubjectFiltersComponent (Wizard) > onDetailsChanged() dovrebbe solo salvare i dati temporaneamente nel signal & component.currentDetails() deve essere uguale a mockDetails & \ok \\
TU-977 & SubjectFiltersComponent (Wizard) > addSubjectToList() dovrebbe emettere e aggiornare la lista se i dati sono completi & emitSpy deve essere stato chiamato; component.subjectsList().length deve essere 1; component.subjectsList()[0].role deve essere SubjectRoleType.AUTORE; component.subjectsList()[0].details deve essere uguale a mockDetails; component.currentStep() deve essere 0 & \ok \\
TU-978 & SubjectFiltersComponent (Wizard) > removeSubject() dovrebbe rimuovere l'elemento e lanciare l'evento & component.subjectsList().length deve essere 1; component.subjectsList()[0].role deve essere 'DESTINATARIO'; emitSpy deve essere stato chiamato con component.subjectsList() & \ok \\
TU-979 & SubjectFiltersComponent (Wizard) > dovrebbe renderizzare lo Step 1 e coprire click e annulla & buttons.length deve essere component.availableRoles().length; actionBtns.length deve essere 1; component.currentStep() deve essere 0 & \ok \\
TU-980 & SubjectFiltersComponent (Wizard) > dovrebbe renderizzare lo Step 2 e coprire click, indietro e annulla & buttons.length deve essere component.allowedTypes().length; actionBtns.length deve essere 2; component.currentStep() deve essere 1; component.currentStep() deve essere 0 & \ok \\
TU-981 & SubjectFiltersComponent (Wizard) > dovrebbe renderizzare lo Step 3 e coprire azioni e binding del figlio & detailForm deve essere truthy; actionBtns.length deve essere 3 & \ok \\
TU-982 & AdvancedFilterPanelComponent > Interazione con i componenti figli (Eventi DOM HTML) > dovrebbe intercettare onEntriesChanged da app-custom-meta-filters (HTML righe 39-40) & emitSpy deve essere stato chiamato con object containing { customMeta: mockEntries } & \ok \\
TU-983 & AdvancedFilterPanelComponent > Interazione con i componenti figli (Eventi DOM HTML) > dovrebbe intercettare onSubjectChanged da app-subject-filters & emitSpy deve essere stato chiamato con object containing { subject: mockSubject } & \ok \\
TU-984 & AdvancedFilterPanelComponent > Submit e Disabilitazione Bottone > onSubmit() NON dovrebbe emettere se la validazione esterna ha trovato errori & emitSubmitSpy non deve essere stato chiamato & \ok \\
TU-985 & AdvancedFilterPanelComponent > Submit e Disabilitazione Bottone > dovrebbe disabilitare il bottone Applica se currentValidationResult è invalido (HTML riga 104) & submitBtn.disabled deve essere true & \ok \\
TU-986 & AdvancedFilterPanelComponent > Submit e Disabilitazione Bottone > dovrebbe emettere filtersSubmit se non ci sono errori e si preme Applica (TS riga 100) & emitSubmitSpy deve essere stato chiamato & \ok \\
TU-987 & AdvancedFilterPanelComponent > Logica Reattiva (validateAndEmit) > dovrebbe chiamare validateAndEmit al variare del form, aggiornando validationResult e filtersChanged (TS 108-120) & mockValidatorFn deve essere stato chiamato; emitValidationSpy deve essere stato chiamato; emitFiltersSpy deve essere stato chiamato con object containing { common: expect.objectContaining({ testo: 'prova reattività' }), } & \ok \\
TU-988 & AdvancedFilterPanelComponent > Altre Funzioni Base > dovrebbe resettare il form e il validationResult con onReset & component.currentValidationResult deve essere null; emitResetSpy deve essere stato chiamato & \ok \\
TU-989 & AdvancedFilterPanelComponent > Getter/Setter filters e Logica di Reset > dovrebbe eseguire patchValue (ramo else) se i filtri NON sono vuoti & patchSpy deve essere stato chiamato con newFilters, { emitEvent: false } & \ok \\
TU-990 & AdvancedFilterPanelComponent > Getter/Setter filters e Logica di Reset > dovrebbe eseguire reset (ramo if) e incrementare subjectResetCounter se i filtri SONO vuoti & resetSpy deve essere stato chiamato con {}, { emitEvent: false }; component.subjectResetCounter deve essere initialCounter + 1 & \ok \\
TU-991 & AdvancedFilterPanelComponent > Getter/Setter filters e Logica di Reset > dovrebbe gestire in modo sicuro i valori undefined (copertura dei rami || {}) & resetSpy deve essere stato chiamato con {}, { emitEvent: false } & \ok \\
TU-992 & AdvancedFilterPanelComponent > Getter/Setter filters e Logica di Reset > NON dovrebbe fare nulla se panelForm non è ancora inizializzato & () => { component.filters = { common: { t: 1 } } as any; } non deve lanciare un errore & \ok \\
TU-993 & FilterRulesManager > calculateUIState() > should not block any filters when form is empty or null & state.disableAggregate deve essere false; state.disableDiDai deve essere false; state.motivoBloccoAggregate deve essere ''; state.motivoBloccoDiDai deve essere '' & \ok \\
TU-994 & FilterRulesManager > calculateUIState() > should block Aggregate filters when tipoDocumento is DOCUMENTO\_INFORMATICO & state.disableAggregate deve essere true; state.motivoBloccoAggregate deve essere 'Filtri bloccati: hai selezionato il Tipo Documento.'; state.disableDiDai deve essere false & \ok \\
TU-995 & FilterRulesManager > calculateUIState() > should block Aggregate filters when tipoDocumento is DOCUMENTO\_AMMINISTRATIVO\_INFORMATICO & state.disableAggregate deve essere true; state.motivoBloccoAggregate deve essere 'Filtri bloccati: hai selezionato il Tipo Documento.' & \ok \\
TU-996 & FilterRulesManager > calculateUIState() > should block Aggregate filters when diDai fields are compiled & state.disableAggregate deve essere true; state.motivoBloccoAggregate deve essere 'Filtri bloccati: stai compilando i parametri DiDai.' & \ok \\
TU-997 & FilterRulesManager > calculateUIState() > should block DiDai filters when tipoDocumento is AGGREGAZIONE\_DOCUMENTALE & state.disableDiDai deve essere true; state.motivoBloccoDiDai deve essere 'Filtri bloccati: hai selezionato il Tipo Fascicolo.'; state.disableAggregate deve essere false & \ok \\
TU-998 & FilterRulesManager > calculateUIState() > should block DiDai filters when aggregate fields are compiled & state.disableDiDai deve essere true; state.motivoBloccoDiDai deve essere 'Filtri bloccati: stai compilando i parametri Aggregazione.' & \ok \\
TU-999 & FilterRulesManager > hasValoriCompilati() > should return false for null, undefined, or primitive types (not an object) & FilterRulesManager.hasValoriCompilati(null) deve essere false; FilterRulesManager.hasValoriCompilati(undefined) deve essere false; FilterRulesManager.hasValoriCompilati('just a string') deve essere false; FilterRulesManager.hasValoriCompilati(123) deve essere false; FilterRulesManager.hasValoriCompilati(true) deve essere false & \ok \\
TU-1000 & FilterRulesManager > hasValoriCompilati() > should return true if the object itself is a Date & FilterRulesManager.hasValoriCompilati(new Date()) deve essere true & \ok \\
TU-1001 & FilterRulesManager > hasValoriCompilati() > should return false for objects containing only null, undefined, or empty strings & FilterRulesManager.hasValoriCompilati({ a: null, b: undefined, c: '', d: ' ' }) deve essere false & \ok \\
TU-1002 & FilterRulesManager > hasValoriCompilati() > should return false for objects containing false boolean values & FilterRulesManager.hasValoriCompilati({ a: false }) deve essere false & \ok \\
TU-1003 & FilterRulesManager > hasValoriCompilati() > should return false for objects containing empty arrays & FilterRulesManager.hasValoriCompilati({ items: [] }) deve essere false & \ok \\
TU-1004 & FilterRulesManager > hasValoriCompilati() > should return false recursively for nested objects with only empty values & FilterRulesManager.hasValoriCompilati({ wrapper: { nested: { prop: '' } } }) deve essere false & \ok \\
TU-1005 & FilterRulesManager > hasValoriCompilati() > should return true for objects containing a nested Date & FilterRulesManager.hasValoriCompilati({ dateProp: new Date() }) deve essere true & \ok \\
TU-1006 & FilterRulesManager > hasValoriCompilati() > should return true for nested objects with at least one valid value & FilterRulesManager.hasValoriCompilati({ wrapper: { nested: { prop: 'valid' } } }) deve essere true & \ok \\
TU-1007 & FilterRulesManager > hasValoriCompilati() > should return true for objects containing numbers, true booleans, or populated arrays & FilterRulesManager.hasValoriCompilati({ numberProp: 0 }) deve essere true; FilterRulesManager.hasValoriCompilati({ booleanProp: true }) deve essere true; FilterRulesManager.hasValoriCompilati({ arrayProp: [1, 2, 3] }) deve essere true & \ok \\
TU-1008 & SearchPageComponent > dovrebbe crearsi correttamente e leggere lo stato iniziale & component deve essere truthy; mockFacade.getState deve essere stato chiamato; mockFacade.loadCustomMetadataKeys deve essere stato chiamato con & \ok \\
TU-1009 & SearchPageComponent > Interazione con la Barra di Ricerca (Eventi DOM) > dovrebbe aggiornare la query intercettando (queryChanged) dalla search-bar & mockFacade.setQuery deve essere stato chiamato & \ok \\
TU-1010 & SearchPageComponent > Interazione con la Barra di Ricerca (Eventi DOM) > dovrebbe lanciare la ricerca intercettando (searchRequested) dalla search-bar & mockFacade.search deve essere stato chiamato & \ok \\
TU-1011 & SearchPageComponent > Orchestrazione della Ricerca > dovrebbe chiamare searchSemantic se useSemanticSearch è true & mockFacade.searchSemantic deve essere stato chiamato; mockFacade.search non deve essere stato chiamato & \ok \\
TU-1012 & SearchPageComponent > Orchestrazione della Ricerca > dovrebbe chiamare search se useSemanticSearch è false & mockFacade.search deve essere stato chiamato; mockFacade.searchSemantic non deve essere stato chiamato & \ok \\
TU-1013 & SearchPageComponent > Orchestrazione della Ricerca > dovrebbe mostrare externalValidation se la facade espone errori di validazione nello stato & component.externalValidation?.isValid deve essere false; component.externalValidation?.errors.get('aggregate.dataApertura')?.[0] deve essere uguale a validationError; mockFacade.search deve essere stato chiamato & \ok \\
TU-1014 & SearchPageComponent > Orchestrazione della Ricerca > validateFilters() dovrebbe delegare alla facade & mockFacade.validateFilters deve essere stato chiamato con filters & \ok \\
TU-1015 & SearchPageComponent > Interazione con i Filtri (Eventi DOM) > dovrebbe passare i nuovi filtri alla facade intercettando (filtersChanged) dal figlio & mockFacade.setFilters deve essere stato chiamato con mockFilters & \ok \\
TU-1016 & SearchPageComponent > Interazione con i Filtri (Eventi DOM) > dovrebbe resettare i filtri intercettando (filtersReset) dal figlio & mockFacade.setFilters deve essere stato chiamato con { common: {}, diDai: {}, aggregate: {}, customMeta: null, subject: [], } & \ok \\
TU-1017 & SearchPageComponent > Interazione con i Filtri (Eventi DOM) > dovrebbe intercettare (filtersSubmit) dal figlio ed eseguire la ricerca avanzata & mockFacade.searchAdvanced deve essere stato chiamato con mockFilters; mockFacade.search non deve essere stato chiamato & \ok \\
TU-1018 & SearchPageComponent > Interazione con i Filtri (Eventi DOM) > dovrebbe aggiornare externalValidation intercettando (validationResult) dal figlio & component.externalValidation deve essere uguale a result & \ok \\
TU-1019 & SearchPageComponent > Rendering Dinamico (Control Flow e Risultati) > dovrebbe mostrare il loader se loading è true & loader deve essere truthy; loader.nativeElement.textContent deve contenere 'Ricerca in corso' & \ok \\
TU-1020 & SearchPageComponent > Rendering Dinamico (Control Flow e Risultati) > dovrebbe mostrare il box di errore se presente un errore nello stato & errorBox deve essere truthy; errorBox.nativeElement.textContent deve contenere 'Errore di Rete'; retryBtn deve essere truthy & \ok \\
TU-1021 & SearchPageComponent > Rendering Dinamico (Control Flow e Risultati) > dovrebbe mostrare il messaggio "Nessun risultato" se array vuoto e query testuale presente & emptyState deve essere truthy; emptyState.nativeElement.textContent deve contenere 'Nessun risultato trovato' & \ok \\
TU-1022 & SearchPageComponent > Rendering Dinamico (Control Flow e Risultati) > dovrebbe mostrare i risultati se ci sono elementi trovati & resultsComponent deve essere truthy; pageTitle.nativeElement.textContent deve contenere 'Trovati 1 risultati' & \ok \\
TU-1023 & getIntegrityClass > should return "integrity--valid" when status is VALID & getIntegrityClass(IntegrityStatusEnum.VALID) deve essere 'integrity--valid' & \ok \\
TU-1024 & getIntegrityClass > should return "integrity--invalid" when status is INVALID & getIntegrityClass(IntegrityStatusEnum.INVALID) deve essere 'integrity--invalid' & \ok \\
TU-1025 & getIntegrityClass > should return "integrity--unknown" when status is UNKNOWN & getIntegrityClass(IntegrityStatusEnum.UNKNOWN) deve essere 'integrity--unknown' & \ok \\
TU-1026 & getIntegrityClass > should return "integrity--unknown" for any unmapped or default status & getIntegrityClass('NON\_EXISTENT\_STATUS' as IntegrityStatusEnum) deve essere 'integrity--unknown' & \ok \\
TU-1027 & FilterValidatorService > dovrebbe inizializzarsi automaticamente con le 5 strategie di ricerca & strategies.length deve essere 5 & \ok \\
TU-1028 & FilterValidatorService > dovrebbe aggregare correttamente gli errori di una nuova strategia aggiunta dinamicamente & result.isValid deve essere false; result.errors.has('testField') deve essere true & \ok \\
TU-1029 & FilterValidatorService > dovrebbe restituire isValid: true se tutte le strategie passano con successo & result.isValid deve essere true; result.errors.size deve essere 0 & \ok \\
TU-1030 & AggregationContradictionStrategy > Casi di Successo (Nessuna Contraddizione) > dovrebbe passare se non ci sono filtri aggregate & result.size deve essere 0 & \ok \\
TU-1031 & AggregationContradictionStrategy > Casi di Successo (Nessuna Contraddizione) > dovrebbe passare se l'utente non specifica il tipo ma compila i campi (ricerca broad) & result.size deve essere 0 & \ok \\
TU-1032 & AggregationContradictionStrategy > Casi di Successo (Nessuna Contraddizione) > dovrebbe passare con una combinazione valida completa & result.size deve essere 0 & \ok \\
TU-1033 & AggregationContradictionStrategy > Casi di Fallimento (Contraddizioni Logiche) > dovrebbe fallire se Tipo Aggregazione NON è Fascicolo ma viene specificata una Tipologia Fascicolo & result.size deve essere 1; errors deve essere definito; errors?.[0].code deve essere 'AGG\_CNTR\_001' & \ok \\
TU-1034 & AggregationContradictionStrategy > Casi di Fallimento (Contraddizioni Logiche) > dovrebbe fallire se Tipo Fascicolo NON è Procedimento ma vengono compilati i dati del Procedimento & result.size deve essere 1; errors deve essere definito; errors?.[0].code deve essere 'AGG\_CNTR\_002' & \ok \\
TU-1035 & AggregationContradictionStrategy > Casi di Fallimento (Contraddizioni Logiche) > dovrebbe fallire se Tipo Aggregazione NON è Fascicolo ma viene compilata l'Assegnazione & result.size deve essere 1; errors deve essere definito; errors?.[0].code deve essere 'AGG\_CNTR\_003' & \ok \\
TU-1036 & AggregationContradictionStrategy > Casi di Fallimento (Contraddizioni Logiche) > dovrebbe accumulare errori multipli se ci sono più contraddizioni contemporaneamente & result.size deve essere 2; result.has('aggregate.tipoFascicolo') deve essere true; result.has('aggregate.assegnazione') deve essere true & \ok \\
TU-1037 & DateTimeCoherenceValidationStrategy > dovrebbe segnalare errore se oraRegistrazione è valorizzata ma dataRegistrazione è assente & result.has('diDai.registrazione.dataRegistrazione') deve essere true; result.get('diDai.registrazione.dataRegistrazione')?.[0].code deve essere 'ERR\_CONF\_DT\_001' & \ok \\
TU-1038 & DateTimeCoherenceValidationStrategy > dovrebbe segnalare errore se oraModifica è valorizzata ma dataModifica è assente & result.has('diDai.tracciatureModifiche.0.dataModifica') deve essere true; result.get('diDai.tracciatureModifiche.0.dataModifica')?.[0].code deve essere 'ERR\_CONF\_DT\_002' & \ok \\
TU-1039 & DateTimeCoherenceValidationStrategy > NON dovrebbe segnalare errori quando data e ora sono coerenti & result.size deve essere 0 & \ok \\
TU-1040 & FormationModeContradictionStrategy > Assenza di campi scatenanti (Uscite anticipate) > NON dovrebbe segnalare errore se manca completamente il blocco diDai & result.size deve essere 0 & \ok \\
TU-1041 & FormationModeContradictionStrategy > Assenza di campi scatenanti (Uscite anticipate) > NON dovrebbe segnalare errore se c'è la modalità ma conformitaCopie è null/undefined & result.size deve essere 0 & \ok \\
TU-1042 & FormationModeContradictionStrategy > Assenza di campi scatenanti (Uscite anticipate) > NON dovrebbe segnalare errore se c'è conformitaCopie ma manca la modalità & result.size deve essere 0 & \ok \\
TU-1043 & FormationModeContradictionStrategy > Contraddizioni Logiche > dovrebbe segnalare errore se cerco modalità EX\_NOVO e imposto conformitaCopie a true & result.has('diDai.verifica.conformitaCopie') deve essere true; result.get('diDai.verifica.conformitaCopie')?.[0].code deve essere 'ERR\_CONF\_FORM\_001' & \ok \\
TU-1044 & FormationModeContradictionStrategy > Contraddizioni Logiche > dovrebbe segnalare errore se cerco modalità MEMORIZZAZIONE e imposto conformitaCopie a false & result.has('diDai.verifica.conformitaCopie') deve essere true & \ok \\
TU-1045 & FormationModeContradictionStrategy > Ricerche Valide > NON dovrebbe segnalare errore se cerco modalità ACQUISIZIONE e imposto conformitaCopie a true & result.size deve essere 0 & \ok \\
TU-1046 & FormationModeContradictionStrategy > Ricerche Valide > NON dovrebbe segnalare errore se cerco modalità legacy "b" e imposto conformitaCopie a true & result.size deve essere 0 & \ok \\
TU-1047 & FormationModeContradictionStrategy > Ricerche Valide > NON dovrebbe segnalare errore se cerco modalità "B" (maiuscolo) e imposto conformitaCopie & result.size deve essere 0 & \ok \\
TU-1048 & RegistryContradictionValidationStrategy > Assenza del blocco principale > NON dovrebbe eseguire validazioni se manca il blocco diDai o registrazione & result.size deve essere 0 & \ok \\
TU-1049 & RegistryContradictionValidationStrategy > Contraddizioni Logiche > dovrebbe segnalare errore se cerco Tipo "Nessuno" ma inserisco un codiceRegistro & result.has('diDai.registrazione.codiceRegistro') deve essere true; result.get('diDai.registrazione.codiceRegistro')?.[0].code deve essere 'ERR\_CONF\_REG\_001' & \ok \\
TU-1050 & RegistryContradictionValidationStrategy > Contraddizioni Logiche > dovrebbe segnalare errore se cerco Tipo "Nessuno" ma inserisco un numeroRegistrazione (copre ramo hasValue su numero) & result.has('diDai.registrazione.numeroRegistrazione') deve essere true; result.get('diDai.registrazione.numeroRegistrazione')?.[0].code deve essere 'ERR\_CONF\_REG\_002' & \ok \\
TU-1051 & RegistryContradictionValidationStrategy > Contraddizioni Logiche > dovrebbe segnalare entrambi gli errori se tutti e due i campi sono compilati assieme a Tipo "Nessuno" & result.size deve essere 2; result.has('diDai.registrazione.codiceRegistro') deve essere true; result.has('diDai.registrazione.numeroRegistrazione') deve essere true & \ok \\
TU-1052 & RegistryContradictionValidationStrategy > Ricerche Valide (Filtri Parziali) > NON dovrebbe segnalare errore se cerco Tipo "Nessuno" e lascio i campi vuoti (ricerca lecita) & result.size deve essere 0 & \ok \\
TU-1053 & RegistryContradictionValidationStrategy > Ricerche Valide (Filtri Parziali) > NON dovrebbe segnalare errore se cerco un Tipo di protocollo valido (es. ORDINARIO) e inserisco i campi & result.size deve essere 0 & \ok \\
TU-1054 & SearchRangeValidationStrategy > Assenza di campi scatenanti (Uscite anticipate) > NON dovrebbe eseguire validazioni se manca completamente il blocco common & result.size deve essere 0 & \ok \\
TU-1055 & SearchRangeValidationStrategy > Assenza di campi scatenanti (Uscite anticipate) > NON dovrebbe segnalare errore se è presente solo la dataDa o solo la dataA & result.size deve essere 0 & \ok \\
TU-1056 & SearchRangeValidationStrategy > Assenza di campi scatenanti (Uscite anticipate) > NON dovrebbe segnalare errore se è presente solo il minimo o solo il massimo numerico & result.size deve essere 0 & \ok \\
TU-1057 & SearchRangeValidationStrategy > Contraddizioni di Range Temporali > dovrebbe segnalare errore se dataDa è successiva a dataA & result.has('common.dataDa') deve essere true; result.get('common.dataDa')?.[0].code deve essere 'ERR\_RANGE\_001' & \ok \\
TU-1058 & SearchRangeValidationStrategy > Contraddizioni di Range Temporali > NON dovrebbe segnalare errore se dataDa è precedente o uguale a dataA & strategy.validate(filters1).size deve essere 0; strategy.validate(filters2).size deve essere 0 & \ok \\
TU-1059 & SearchRangeValidationStrategy > Contraddizioni di Range Temporali > dovrebbe segnalare errore se dataApertura è successiva a dataChiusura & result.has('aggregate.dataApertura') deve essere true; result.get('aggregate.dataApertura')?.[0].code deve essere 'ERR\_RANGE\_001' & \ok \\
TU-1060 & SearchRangeValidationStrategy > Contraddizioni di Range Temporali > NON dovrebbe segnalare errore se dataApertura è precedente o uguale a dataChiusura & strategy.validate(filters1).size deve essere 0; strategy.validate(filters2).size deve essere 0 & \ok \\
TU-1061 & SearchRangeValidationStrategy > Contraddizioni di Range Temporali > dovrebbe segnalare errore se dataInizioAssegn è successiva a dataFineAssegn & result.has('aggregate.assegnazione.dataInizioAssegn') deve essere true; result.get('aggregate.assegnazione.dataInizioAssegn')?.[0].code deve essere 'ERR\_RANGE\_001' & \ok \\
TU-1062 & SearchRangeValidationStrategy > Contraddizioni di Range Temporali > NON dovrebbe segnalare errore se dataInizioAssegn è precedente o uguale a dataFineAssegn & strategy.validate(filters1).size deve essere 0; strategy.validate(filters2).size deve essere 0 & \ok \\
TU-1063 & SearchRangeValidationStrategy > Contraddizioni di Range Temporali > dovrebbe segnalare errore se dataInizioFase è successiva a dataFineFase & result.has('aggregate.procedimento.fasi.0.dataInizioFase') deve essere true; result.get('aggregate.procedimento.fasi.0.dataInizioFase')?.[0].code deve essere 'ERR\_RANGE\_001' & \ok \\
TU-1064 & SearchRangeValidationStrategy > Contraddizioni di Range Temporali > NON dovrebbe segnalare errore se dataInizioFase è precedente o uguale a dataFineFase & strategy.validate(filters1).size deve essere 0; strategy.validate(filters2).size deve essere 0 & \ok \\
TU-1065 & SearchRangeValidationStrategy > Contraddizioni di Range Numerici > dovrebbe segnalare errore se il minimo è maggiore del massimo & result.has('common.numeroAllegatiMin') deve essere true; result.get('common.numeroAllegatiMin')?.[0].code deve essere 'ERR\_RANGE\_002' & \ok \\
TU-1066 & SearchRangeValidationStrategy > Contraddizioni di Range Numerici > NON dovrebbe segnalare errore se il minimo è minore o uguale al massimo & strategy.validate(filters1).size deve essere 0; strategy.validate(filters2).size deve essere 0 & \ok \\
TU-1067 & IntegrityIpcGateway > dovrebbe istanziare il servizio & gateway deve essere truthy & \ok \\
TU-1068 & IntegrityIpcGateway > dovrebbe lanciare un errore se la Electron API non è disponibile & gateway.checkDipIntegrity(1) deve lanciare `Electron API not available for channel IpcChannels.CHECK\_DIP\_INTEGRITY\_STATUS` & \ok \\
TU-1069 & IntegrityIpcGateway > dovrebbe usare (window as any).api come fallback se electronAPI non cè & mockInvoke deve essere stato chiamato con IpcChannels.CHECK\_DIP\_INTEGRITY\_STATUS, 1; result deve essere IntegrityStatusEnum.VALID & \ok \\
TU-1070 & IntegrityIpcGateway > Command IPC Calls > checkDipIntegrity & mockInvoke deve essere stato chiamato con IpcChannels.CHECK\_DIP\_INTEGRITY\_STATUS, 10; res deve essere IntegrityStatusEnum.VALID & \ok \\
TU-1071 & IntegrityIpcGateway > Command IPC Calls > checkDocumentClassIntegrity & mockInvoke deve essere stato chiamato con IpcChannels.CHECK\_DOCUMENT\_CLASS\_INTEGRITY\_STATUS, 11; res deve essere IntegrityStatusEnum.INVALID & \ok \\
TU-1072 & IntegrityIpcGateway > Command IPC Calls > checkProcessIntegrity & mockInvoke deve essere stato chiamato con IpcChannels.CHECK\_PROCESS\_INTEGRITY\_STATUS, 12; res deve essere IntegrityStatusEnum.UNKNOWN & \ok \\
TU-1073 & IntegrityIpcGateway > Command IPC Calls > checkDocumentIntegrity & mockInvoke deve essere stato chiamato con IpcChannels.CHECK\_DOCUMENT\_INTEGRITY\_STATUS, 13; res deve essere IntegrityStatusEnum.VALID & \ok \\
TU-1074 & IntegrityIpcGateway > Command IPC Calls > checkFileIntegrity & mockInvoke deve essere stato chiamato con IpcChannels.CHECK\_FILE\_INTEGRITY\_STATUS, 14; res deve essere IntegrityStatusEnum.INVALID & \ok \\
TU-1075 & IntegrityIpcGateway > Query IPC Calls > getClassesByDipId & mockInvoke deve essere stato chiamato con IpcChannels.BROWSE\_GET\_DOCUMENT\_CLASS\_BY\_DIP\_ID, 20; res deve essere uguale a mockResult & \ok \\
TU-1076 & IntegrityIpcGateway > Query IPC Calls > getProcessesByClassId & mockInvoke deve essere stato chiamato con IpcChannels.BROWSE\_GET\_PROCESS\_BY\_DOCUMENT\_CLASS, 21; res deve essere uguale a mockResult & \ok \\
TU-1077 & IntegrityIpcGateway > Query IPC Calls > getDocumentsByProcessId & mockInvoke deve essere stato chiamato con IpcChannels.BROWSE\_GET\_DOCUMENTS\_BY\_PROCESS, 22; res deve essere uguale a mockResult & \ok \\
TU-1078 & IntegrityIpcGateway > Query IPC Calls > getFilesByDocumentId & mockInvoke deve essere stato chiamato con IpcChannels.BROWSE\_GET\_FILE\_BY\_DOCUMENT, 23; res deve essere uguale a mockResult & \ok \\
TU-1079 & IntegrityFacade > verifyDip invalida tutti i prefissi cache del dettaglio & mockCache.invalidatePrefix deve essere stato chiamato con 'aggregate:'; mockCache.invalidatePrefix deve essere stato chiamato con 'process:'; mockCache.invalidatePrefix deve essere stato chiamato con 'document:'; mockCache.invalidatePrefix deve essere stato chiamato con 'node-fallback:'; facade.currentDipStatus()() deve essere IntegrityStatusEnum.VALID & \ok \\
TU-1080 & IntegrityFacade > verifyItem DOCUMENT invalida la chiave documento specifica & result deve essere IntegrityStatusEnum.VALID; mockGateway.checkDocumentIntegrity deve essere stato chiamato con 12; mockCache.invalidate deve essere stato chiamato con 'document:12' & \ok \\
TU-1081 & IntegrityFacade > verifyItem PROCESS invalida process/aggregate e prefisso document & result deve essere IntegrityStatusEnum.INVALID; mockGateway.checkProcessIntegrity deve essere stato chiamato con 31; mockCache.invalidate deve essere stato chiamato con 'process:31'; mockCache.invalidate deve essere stato chiamato con 'aggregate:31'; mockCache.invalidatePrefix deve essere stato chiamato con 'document:' & \ok \\
TU-1082 & IntegrityFacade > verifyItem DOCUMENT\_CLASS invalida cache fallback e prefissi correlati & result deve essere IntegrityStatusEnum.UNKNOWN; mockGateway.checkDocumentClassIntegrity deve essere stato chiamato con 22; mockCache.invalidate deve essere stato chiamato con 'node-fallback:DOCUMENT\_CLASS:22'; mockCache.invalidatePrefix deve essere stato chiamato con 'process:'; mockCache.invalidatePrefix deve essere stato chiamato con 'aggregate:'; mockCache.invalidatePrefix deve essere stato chiamato con 'document:' & \ok \\
TU-1083 & IntegrityFacade > verifyItem AGGREGATE invalida aggregate/process e prefisso document & result deve essere IntegrityStatusEnum.VALID; mockGateway.checkProcessIntegrity deve essere stato chiamato con 44; mockCache.invalidate deve essere stato chiamato con 'aggregate:44'; mockCache.invalidate deve essere stato chiamato con 'process:44'; mockCache.invalidatePrefix deve essere stato chiamato con 'document:' & \ok \\
TU-1084 & IntegrityFacade > verifyItem in errore ritorna UNKNOWN e aggiorna errore tramite handler & result deve essere 'UNKNOWN'; mockErrorHandler.handle deve essere stato chiamato 1 volte; handledError.source deve essere 'IntegrityFacade.verifyItem'; handledError.context deve essere uguale a { itemId: '55', itemType: 'DOCUMENT', action: 'CHECK\_ITEM\_INTEGRITY', }; facade.error()() deve essere 'errore test' & \ok \\
TU-1085 & IntegrityFacade > clearResults resetta stato, classi ed errore quando non sta verificando & facade.currentDipStatus()() deve essere IntegrityStatusEnum.INVALID; facade.dipClasses()() deve essere uguale a [{ id: 1, name: 'Classe A' }]; facade.currentDipStatus()() deve essere null; facade.dipClasses()() deve essere uguale a []; facade.error()() deve essere null & \ok \\
TU-1086 & IntegrityFacade > clearResults non modifica i dati se verifica in corso & facade.currentDipStatus()() deve essere IntegrityStatusEnum.VALID; facade.dipClasses()() deve essere uguale a [{ id: 2, name: 'Classe B' }] & \ok \\
TU-1087 & IntegritySummaryCardsComponent > should create & component deve essere truthy & \ok \\
TU-1088 & IntegritySummaryCardsComponent > should render the stats correctly & validCard?.textContent deve contenere '5'; validCard?.textContent deve contenere 'Elementi Integri'; invalidCard?.textContent deve contenere '3'; invalidCard?.textContent deve contenere 'Elementi Corrotti'; unknownCard?.textContent deve contenere '2'; unknownCard?.textContent deve contenere 'In Attesa di Verifica' & \ok \\
TU-1089 & IntegrityValidPanelComponent > should create & component deve essere truthy & \ok \\
TU-1090 & IntegrityValidPanelComponent > should not render anything if nodes is empty & compiled.querySelector('.valid-container') deve essere null & \ok \\
TU-1091 & IntegrityValidPanelComponent > should render valid nodes when provided & compiled.querySelector('.valid-container') deve essere truthy; items.length deve essere 1; items[0].textContent deve contenere 'Classe'; items[0].textContent deve contenere 'Node Valid 1'; items[0].textContent deve contenere 'Valido' & \ok \\
TU-1092 & IntegrityValidPanelComponent > should render invalid label when node status is INVALID & compiled.querySelector('.status-pill')?.textContent deve contenere 'Invalido'; compiled.querySelector('.status-pill--invalid') deve essere truthy; compiled.querySelector('.error-group') deve essere truthy; compiled.textContent deve contenere 'Processo A'; compiled.textContent deve contenere 'Doc 1' & \ok \\
TU-1093 & IntegrityValidPanelComponent > should not render error details for valid nodes & compiled.querySelector('.error-group') deve essere null & \ok \\
TU-1094 & IntegrityValidPanelComponent > should format type correctly & component.formatType('CLASS') deve essere 'Classe'; component.formatType('PROCESS') deve essere 'Processo'; component.formatType('DOCUMENT') deve essere 'Documento'; component.formatType('SOME\_UNKNOWN\_TYPE') deve essere 'Documento' & \ok \\
TU-1095 & IntegrityValidPanelComponent > should map status labels and css classes correctly & component.getStatusLabel(IntegrityStatusEnum.VALID) deve essere 'Valido'; component.getStatusLabel(IntegrityStatusEnum.INVALID) deve essere 'Invalido'; component.getStatusLabel(IntegrityStatusEnum.UNKNOWN) deve essere 'Non verificato'; component.getStatusClass(IntegrityStatusEnum.VALID) deve essere 'status-pill status-pill--valid'; component.getStatusClass(IntegrityStatusEnum.INVALID) deve essere 'status-pill status-pill--invalid'; component.getStatusClass(IntegrityStatusEnum.UNKNOWN) deve essere 'status-pill status-pill--unknown' & \ok \\
TU-1096 & IntegrityValidPanelComponent > should correctly sort classNodes: INVALID first, then alphabetically by name & computedNodes.length deve essere 3; computedNodes[0].name deve essere 'Delta'; computedNodes[1].name deve essere 'Alpha'; computedNodes[2].name deve essere 'Zebra' & \ok \\
TU-1097 & IntegrityValidPanelComponent > should parse various contextPath formats correctly in groupedErrors & errors.get('Classe non disponibile')?.processes[0].processName deve essere 'Processo non disponibile'; errors.get('C1')?.processes[0].processName deve essere 'P1'; errors.get('C1')?.totalDocuments deve essere 1; errors.get('C2')?.processes[0].processName deve essere 'P2'; errors.get('C3')?.processes[0].processName deve essere 'Processo non disponibile' & \ok \\
TU-1098 & IntegrityDashboardComponent > should create and load overview on init & component deve essere truthy; mockFacade.loadOverview deve essere stato chiamato con component.currentDipId & \ok \\
TU-1099 & IntegrityDashboardComponent > should display progress section when verifying & compiled.querySelector('.progress-section') deve essere truthy; compiled.querySelector('app-integrity-summary-cards') deve essere null; compiled.querySelector('app-integrity-valid-panel') deve essere null & \ok \\
TU-1100 & IntegrityDashboardComponent > should display summary and invalid-class panel when not verifying & compiled.querySelector('.progress-section') deve essere null; compiled.querySelector('app-integrity-summary-cards') deve essere truthy; compiled.querySelector('app-integrity-valid-panel') deve essere truthy & \ok \\
TU-1101 & IntegrityDashboardComponent > should start verification and reload overview when startVerification is called & mockFacade.verifyDip deve essere stato chiamato con component.currentDipId; mockFacade.loadOverview deve essere stato chiamato 2 volte & \ok \\
TU-1102 & ErrorDialogComponent > UI and Severity Renderings > should render ERROR severity correctly & component.isFatal() deve essere false; component.getIcon() deve essere 'bi-x-octagon-fill'; component.getTitle() deve essere 'Si è verificato un errore'; headerTitle.textContent.trim() deve essere 'Si è verificato un errore'; dialog.classList.contains('fatal-dialog') deve essere false & \ok \\
TU-1103 & ErrorDialogComponent > UI and Severity Renderings > should render FATAL severity and apply the fatal-dialog CSS class & component.isFatal() deve essere true; component.getIcon() deve essere 'bi-bug-fill'; component.getTitle() deve essere 'Errore Critico del Sistema'; dialog.classList.contains('fatal-dialog') deve essere true & \ok \\
TU-1104 & ErrorDialogComponent > UI and Severity Renderings > should render WARNING severity correctly & component.getIcon() deve essere 'bi-exclamation-triangle-fill'; component.getTitle() deve essere 'Attenzione' & \ok \\
TU-1105 & ErrorDialogComponent > UI and Severity Renderings > should render default (INFO/fallback) severity correctly & component.getIcon() deve essere 'bi-info-circle-fill'; component.getTitle() deve essere 'Avviso' & \ok \\
TU-1106 & ErrorDialogComponent > Optional Elements (@if conditions) > should not render code, source, details, or retry button when omitted & monoElements.length deve essere 0; techDetails.length deve essere 0; retryButtons.length deve essere 0 & \ok \\
TU-1107 & ErrorDialogComponent > Optional Elements (@if conditions) > should render code, source, and details when provided & errorMeta.textContent deve contenere 'SEARCH\_ENGINE\_ERROR'; errorMeta.textContent deve contenere 'Backend API'; techDetails.textContent deve contenere 'Stack trace details...' & \ok \\
TU-1108 & ErrorDialogComponent > Event Outputs > should emit onClose when close button is clicked & closeSpy deve essere stato chiamato 1 volte & \ok \\
TU-1109 & ErrorDialogComponent > Event Outputs > should emit onRetry when retry button is clicked (if recoverable) & retrySpy deve essere stato chiamato 1 volte & \ok \\
TU-1110 & ElectronLoggingGateway > dovrebbe chiamare il bridge IPC corretto e completarsi con successo & mockBridge.invoke deve essere stato chiamato con 'ipc:log:write', mockLogEntry & \ok \\
TU-1111 & ElectronLoggingGateway > dovrebbe propagare l'errore se la scrittura IPC fallisce & errorResult deve essere mockError; mockBridge.invoke deve essere stato chiamato 1 volte & \ok \\
TU-1112 & IpcErrorHandlerService > createError() > dovrebbe creare un AppError con la corretta mappatura (es. DOC\_BLOB\_LOAD\_ERROR) & error.code deve essere ErrorCode.DOC\_BLOB\_LOAD\_ERROR; error.message deve essere 'File non trovato'; error.source deve essere 'DocFacade'; error.category deve essere ErrorCategory.IO; error.severity deve essere ErrorSeverity.ERROR; error.recoverable deve essere true & \ok \\
TU-1113 & IpcErrorHandlerService > createError() > dovrebbe creare un AppError fatale e non recuperabile per IPC\_ERROR & error.severity deve essere ErrorSeverity.FATAL; error.recoverable deve essere false & \ok \\
TU-1114 & IpcErrorHandlerService > handle() / toErrorCode() > dovrebbe estrarre correttamente il codice se fornito nell'oggetto raw & result.code deve essere ErrorCode.DOC\_FORMAT\_UNSUPPORTED; result.message deve essere 'Formato .xyz non supportato'; result.source deve essere 'ElectronRenderer'; result.category deve essere ErrorCategory.IO & \ok \\
TU-1115 & IpcErrorHandlerService > handle() / toErrorCode() > dovrebbe fare fallback su SEARCH\_ENGINE\_ERROR se l'oggetto non ha un codice riconosciuto & result.code deve essere ErrorCode.SEARCH\_ENGINE\_ERROR; result.message deve essere 'Errore misterioso nel motore di ricerca'; result.detail deve essere rawError.stack & \ok \\
TU-1116 & IpcErrorHandlerService > handle() / toErrorCode() > dovrebbe dedurre IPC\_ERROR se il messaggio contiene la stringa "ipc" & result.code deve essere ErrorCode.IPC\_ERROR; result.category deve essere ErrorCategory.IPC & \ok \\
TU-1117 & IpcErrorHandlerService > handle() / toErrorCode() > dovrebbe gestire stringhe primitive in modo sicuro & result.code deve essere ErrorCode.SEARCH\_ENGINE\_ERROR; result.message deve essere 'Qualcosa è andato storto'; result.source deve essere 'IPC Gateway' & \ok \\
TU-1118 & IpcErrorHandlerService > Casi limite (Edge Cases) e Coverage > dovrebbe usare la configurazione di default in createError per un codice non mappato (righe 67-77) & error.category deve essere ErrorCategory.IPC; error.severity deve essere ErrorSeverity.ERROR; error.recoverable deve essere false & \ok \\
TU-1119 & IpcErrorHandlerService > Casi limite (Edge Cases) e Coverage > dovrebbe fare fallback se l'oggetto raw ha una proprietà code non presente nell'enum (riga 98) & result.code deve essere ErrorCode.SEARCH\_ENGINE\_ERROR & \ok \\
TU-1120 & IpcCacheService > dovrebbe salvare e recuperare un valore se non è scaduto & result deve essere uguale a { id: 1, name: 'Documento' } & \ok \\
TU-1121 & IpcCacheService > dovrebbe restituire null e invalidare la cache se il TTL è scaduto & result deve essere null & \ok \\
TU-1122 & IpcCacheService > dovrebbe restituire null per una chiave mai inserita & result deve essere null & \ok \\
TU-1123 & IpcCacheService > dovrebbe invalidare esplicitamente una chiave specifica & cacheService.get('manual-key') deve essere null & \ok \\
TU-1124 & IpcCacheService > dovrebbe invalidare solo le chiavi che corrispondono a un prefisso (invalidatePrefix) & cacheService.get('search/123') deve essere null; cacheService.get('search/456') deve essere null; cacheService.get('doc/789') deve essere 'dati documento' & \ok \\
TU-1125 & TelemetryService > trackEvent() dovrebbe inviare un log di livello INFO senza payload & mockLoggingChannel.writeLog deve essere stato chiamato con { level: LogLevel.INFO, event: TelemetryEvent.SEARCH\_EXECUTED, timestamp: MOCK\_TIME, payload: null, } & \ok \\
TU-1126 & TelemetryService > trackTiming() dovrebbe inviare un log di livello INFO con la durata nel payload & mockLoggingChannel.writeLog deve essere stato chiamato con { level: LogLevel.INFO, event: TelemetryMetric.SEARCH\_LATENCY\_MS, timestamp: MOCK\_TIME, payload: { durationMs: 350 }, } & \ok \\
TU-1127 & TelemetryService > trackError() dovrebbe inviare un log di livello ERROR con l'oggetto AppError & mockLoggingChannel.writeLog deve essere stato chiamato con { level: LogLevel.ERROR, event: 'APP\_ERROR', timestamp: MOCK\_TIME, payload: mockError, } & \ok \\
TU-1128 & TelemetryService > dovrebbe gestire in modo sicuro gli errori del canale di logging (fire-and-forget) & () => service.trackEvent(TelemetryEvent.SEARCH\_EXECUTED) non deve lanciare un errore; console.error deve essere stato chiamato con 'Fallimento scrittura telemetria via IPC:', logError & \ok \\
TU-1129 & ElectronLoggingGateway > dovrebbe chiamare il bridge IPC corretto e completarsi con successo & mockBridge.invoke deve essere stato chiamato con 'ipc:log:write', mockLogEntry & \ok \\
TU-1130 & ElectronLoggingGateway > dovrebbe propagare l'errore se la scrittura IPC fallisce & errorResult deve essere mockError; mockBridge.invoke deve essere stato chiamato 1 volte & \ok \\
TU-1131 & IpcErrorHandlerService > createError() > dovrebbe creare un AppError con la corretta mappatura (es. DOC\_BLOB\_LOAD\_ERROR) & error.code deve essere ErrorCode.DOC\_BLOB\_LOAD\_ERROR; error.message deve essere 'File non trovato'; error.source deve essere 'DocFacade'; error.category deve essere ErrorCategory.IO; error.severity deve essere ErrorSeverity.ERROR; error.recoverable deve essere true & \ok \\
TU-1132 & IpcErrorHandlerService > createError() > dovrebbe creare un AppError fatale e non recuperabile per IPC\_ERROR & error.severity deve essere ErrorSeverity.FATAL; error.recoverable deve essere false & \ok \\
TU-1133 & IpcErrorHandlerService > handle() / toErrorCode() > dovrebbe estrarre correttamente il codice se fornito nell'oggetto raw & result.code deve essere ErrorCode.DOC\_FORMAT\_UNSUPPORTED; result.message deve essere 'Formato .xyz non supportato'; result.source deve essere 'ElectronRenderer'; result.category deve essere ErrorCategory.IO & \ok \\
TU-1134 & IpcErrorHandlerService > handle() / toErrorCode() > dovrebbe fare fallback su SEARCH\_ENGINE\_ERROR se l'oggetto non ha un codice riconosciuto & result.code deve essere ErrorCode.SEARCH\_ENGINE\_ERROR; result.message deve essere 'Errore misterioso nel motore di ricerca'; result.detail deve essere rawError.stack & \ok \\
TU-1135 & IpcErrorHandlerService > handle() / toErrorCode() > dovrebbe dedurre IPC\_ERROR se il messaggio contiene la stringa "ipc" & result.code deve essere ErrorCode.IPC\_ERROR; result.category deve essere ErrorCategory.IPC & \ok \\
TU-1136 & IpcErrorHandlerService > handle() / toErrorCode() > dovrebbe gestire stringhe primitive in modo sicuro & result.code deve essere ErrorCode.SEARCH\_ENGINE\_ERROR; result.message deve essere 'Qualcosa è andato storto'; result.source deve essere 'IPC Gateway' & \ok \\
TU-1137 & IpcErrorHandlerService > Casi limite (Edge Cases) e Coverage > dovrebbe usare la configurazione di default in createError per un codice non mappato (righe 67-77) & error.category deve essere ErrorCategory.IPC; error.severity deve essere ErrorSeverity.ERROR; error.recoverable deve essere false & \ok \\
TU-1138 & IpcErrorHandlerService > Casi limite (Edge Cases) e Coverage > dovrebbe fare fallback se l'oggetto raw ha una proprietà code non presente nell'enum (riga 98) & result.code deve essere ErrorCode.SEARCH\_ENGINE\_ERROR & \ok \\
TU-1139 & IpcCacheService > dovrebbe salvare e recuperare un valore se non è scaduto & result deve essere uguale a { id: 1, name: 'Documento' } & \ok \\
TU-1140 & IpcCacheService > dovrebbe restituire null e invalidare la cache se il TTL è scaduto & result deve essere null & \ok \\
TU-1141 & IpcCacheService > dovrebbe restituire null per una chiave mai inserita & result deve essere null & \ok \\
TU-1142 & IpcCacheService > dovrebbe invalidare esplicitamente una chiave specifica & cacheService.get('manual-key') deve essere null & \ok \\
TU-1143 & IpcCacheService > dovrebbe invalidare solo le chiavi che corrispondono a un prefisso (invalidatePrefix) & cacheService.get('search/123') deve essere null; cacheService.get('search/456') deve essere null; cacheService.get('doc/789') deve essere 'dati documento' & \ok \\
TU-1144 & LiveAnnouncerService > dovrebbe chiamare il CDK LiveAnnouncer con il livello polite omettendo il parametro & mockCdkAnnouncer.announce deve essere stato chiamato con 'Ricerca completata', 'polite' & \ok \\
TU-1145 & LiveAnnouncerService > dovrebbe usare il livello polite come fallback se viene passato esplicitamente undefined (copertura riga 6) & mockCdkAnnouncer.announce deve essere stato chiamato con 'Ricerca completata', 'polite' & \ok \\
TU-1146 & LiveAnnouncerService > dovrebbe rispettare il livello assertive se esplicitamente richiesto & mockCdkAnnouncer.announce deve essere stato chiamato con 'Errore di sistema critico', 'assertive' & \ok \\
TU-1147 & TelemetryService > trackEvent() dovrebbe inviare un log di livello INFO senza payload & mockLoggingChannel.writeLog deve essere stato chiamato con { level: LogLevel.INFO, event: TelemetryEvent.SEARCH\_EXECUTED, timestamp: MOCK\_TIME, payload: null, } & \ok \\
TU-1148 & TelemetryService > trackTiming() dovrebbe inviare un log di livello INFO con la durata nel payload & mockLoggingChannel.writeLog deve essere stato chiamato con { level: LogLevel.INFO, event: TelemetryMetric.SEARCH\_LATENCY\_MS, timestamp: MOCK\_TIME, payload: { durationMs: 350 }, } & \ok \\
TU-1149 & TelemetryService > trackError() dovrebbe inviare un log di livello ERROR con l'oggetto AppError & mockLoggingChannel.writeLog deve essere stato chiamato con { level: LogLevel.ERROR, event: 'APP\_ERROR', timestamp: MOCK\_TIME, payload: mockError, } & \ok \\
TU-1150 & TelemetryService > dovrebbe gestire in modo sicuro gli errori del canale di logging (fire-and-forget) & () => service.trackEvent(TelemetryEvent.SEARCH\_EXECUTED) non deve lanciare un errore; console.error deve essere stato chiamato con 'Fallimento scrittura telemetria via IPC:', logError & \ok \\
TU-1151 & EmptyStateComponent > should create & component deve essere truthy & \ok \\
TU-1152 & simplifyCustomMetadataLabel > mantiene il segmento finale significativo & simplifyCustomMetadataLabel('Process.PreservationSession.DocumentsStats.DocumentsFilesCount') deve essere 'Documents Files Count' & \ok \\
TU-1153 & simplifyCustomMetadataLabel > mantiene due segmenti se il finale e generico & simplifyCustomMetadataLabel('Process.End.Date') deve essere 'End Date' & \ok \\
TU-1154 & simplifyCustomMetadataLabel > rimuove prefissi rumorosi e gestisce valori semplici & simplifyCustomMetadataLabel('Root.CustomMetadata.Process.Status') deve essere 'Status'; simplifyCustomMetadataLabel('DocumentsFilesCount') deve essere 'Documents Files Count' & \ok \\
TU-1155 & toSpacedMetadataLabel > aggiunge spazi a camelCase, PascalCase e underscore & toSpacedMetadataLabel('SoggettoAutoreDellaModifica') deve essere 'Soggetto Autore Della Modifica'; toSpacedMetadataLabel('CodiceFiscale\_PartitaIva') deve essere 'Codice Fiscale Partita Iva' & \ok \\
TU-1156 & normalizeDisplayFileName > estrae il nome file da un path con UUID & normalizeDisplayFileName('/869e1069-e50d-48e6-8191-1c677f3053a2/Allegato\_1.pdf') deve essere 'Allegato\_1.pdf' & \ok \\
TU-1157 & normalizeDisplayFileName > rimuove il prefisso ./ & normalizeDisplayFileName('./Allegato\_1.pdf') deve essere 'Allegato\_1.pdf' & \ok \\
TU-1158 & normalizeDisplayFileName > gestisce path windows & normalizeDisplayFileName('folder\\sub\\FilePrincipale.pdf') deve essere 'FilePrincipale.pdf' & \ok \\
TU-1159 & normalizeDisplayFileName > lascia invariato un nome gia semplice & normalizeDisplayFileName('documento.pdf') deve essere 'documento.pdf' & \ok \\
TU-1160 & normalizeDisplayFileName > ritorna stringa vuota se il valore e vuoto & normalizeDisplayFileName(' ') deve essere '' & \ok \\
TU-1161 & document-index.util > resolveDocumentLabel usa NomeDelDocumento quando presente & label deve essere 'Contratto.pdf' & \ok \\
TU-1162 & document-index.util > resolveDocumentLabel usa fallback su uuid quando metadata e vuoto & label deve essere 'DOC-2' & \ok \\
TU-1163 & document-index.util > mapDocumentsToIndexEntries supporta metadata wrapper & entries deve essere uguale a [ { tipoDocumento: 'Documento', identificativo: 'Documento Wrapper', routeId: '10', }, ] & \ok \\
TU-1164 & document-index.util > mapDocumentsToIndexEntries ritorna array vuoto con input nullo & mapDocumentsToIndexEntries(null) deve essere uguale a [] & \ok \\
TU-1165 & normalizeMetadataNodes > ritorna il contenuto metadata se e gia un array valido & result deve avere lunghezza 2; result[0].name deve essere 'Oggetto' & \ok \\
TU-1166 & normalizeMetadataNodes > estrae i nodi da wrapper con proprietà value & result deve essere uguale a [{ name: 'NomeDelDocumento', value: 'Doc.pdf' }] & \ok \\
TU-1167 & normalizeMetadataNodes > ritorna array vuoto con input non valido & normalizeMetadataNodes(null) deve essere uguale a []; normalizeMetadataNodes(undefined) deve essere uguale a []; normalizeMetadataNodes({ name: 'x', value: 'y' }) deve essere uguale a [] & \ok \\

\hline
\rowcolor{white}
\caption{Report test di unità}
\label{tab:test-unita}
\end{longtable}

\end{document}